%% Generated by Sphinx.
\def\sphinxdocclass{report}
\documentclass[letterpaper,10pt,french]{sphinxmanual}
\ifdefined\pdfpxdimen
   \let\sphinxpxdimen\pdfpxdimen\else\newdimen\sphinxpxdimen
\fi \sphinxpxdimen=.75bp\relax
\ifdefined\pdfimageresolution
    \pdfimageresolution= \numexpr \dimexpr1in\relax/\sphinxpxdimen\relax
\fi
%% let collapsible pdf bookmarks panel have high depth per default
\PassOptionsToPackage{bookmarksdepth=5}{hyperref}

\PassOptionsToPackage{booktabs}{sphinx}
\PassOptionsToPackage{colorrows}{sphinx}

\PassOptionsToPackage{warn}{textcomp}
\usepackage[utf8]{inputenc}
\ifdefined\DeclareUnicodeCharacter
% support both utf8 and utf8x syntaxes
  \ifdefined\DeclareUnicodeCharacterAsOptional
    \def\sphinxDUC#1{\DeclareUnicodeCharacter{"#1}}
  \else
    \let\sphinxDUC\DeclareUnicodeCharacter
  \fi
  \sphinxDUC{00A0}{\nobreakspace}
  \sphinxDUC{2500}{\sphinxunichar{2500}}
  \sphinxDUC{2502}{\sphinxunichar{2502}}
  \sphinxDUC{2514}{\sphinxunichar{2514}}
  \sphinxDUC{251C}{\sphinxunichar{251C}}
  \sphinxDUC{2572}{\textbackslash}
\fi
\usepackage{cmap}
\usepackage[T1]{fontenc}
\usepackage{amsmath,amssymb,amstext}
\usepackage{babel}



\usepackage{tgtermes}
\usepackage{tgheros}
\renewcommand{\ttdefault}{txtt}



\usepackage[Sonny]{fncychap}
\ChNameVar{\Large\normalfont\sffamily}
\ChTitleVar{\Large\normalfont\sffamily}
\usepackage{sphinx}

\fvset{fontsize=auto}
\usepackage{geometry}


% Include hyperref last.
\usepackage{hyperref}
% Fix anchor placement for figures with captions.
\usepackage{hypcap}% it must be loaded after hyperref.
% Set up styles of URL: it should be placed after hyperref.
\urlstyle{same}


\usepackage{sphinxmessages}
\setcounter{tocdepth}{1}



\title{MLang}
\date{janv. 09, 2026}
\release{}
\author{Direction Générale des FInances Publiques}
\newcommand{\sphinxlogo}{\vbox{}}
\renewcommand{\releasename}{}
\makeindex
\begin{document}

\ifdefined\shorthandoff
  \ifnum\catcode`\=\string=\active\shorthandoff{=}\fi
  \ifnum\catcode`\"=\active\shorthandoff{"}\fi
\fi

\pagestyle{empty}
\sphinxmaketitle
\pagestyle{plain}
\sphinxtableofcontents
\pagestyle{normal}
\phantomsection\label{\detokenize{index::doc}}





\chapter{Le langage M}
\label{\detokenize{index:le-langage-m}}
\sphinxstepscope


\section{Introduction au langage M}
\label{\detokenize{intro:introduction-au-langage-m}}\label{\detokenize{intro:intro-m}}\label{\detokenize{intro::doc}}

\subsection{Les éléments de base}
\label{\detokenize{intro:les-elements-de-base}}
\sphinxAtStartPar
Un programme M se caracterise par la définition successive :
\begin{itemize}
\item {} 
\sphinxAtStartPar
de types de variables avec leurs attributs ;

\item {} 
\sphinxAtStartPar
de domaines et d’événements ;

\item {} 
\sphinxAtStartPar
d’espaces de variables ;

\item {} 
\sphinxAtStartPar
d’applications ;

\item {} 
\sphinxAtStartPar
de variables et de constantes ;

\item {} 
\sphinxAtStartPar
des fonctions ;

\item {} 
\sphinxAtStartPar
de règles de calcul associées ou non à une application.

\end{itemize}

\sphinxAtStartPar
Exécuter d’un tel programme pour une application donnée correspond à
exécuter l’ensemble des règles de calcul associée la dite application.

\sphinxAtStartPar
L’ordre d’exécution des règles de calcul dépend des dépendances entre les
variables. Si une variable est affectée dans une règle, alors cette règle
sera exécutée avant toutes celles utilisant la valeur de la dite variable.

\sphinxAtStartPar
Si une variable est lue et écrite cycliquement, l’ordre d’exécution des règles
n’est pas garanti.


\subsection{Définitions par défaut}
\label{\detokenize{intro:definitions-par-defaut}}
\sphinxAtStartPar
Les types de variables, leurs attributs, les domaines et les événements
sont historiquement des mots clés propres au langage M.

\sphinxAtStartPar
Il est aujourd’hui possible de les définir à la main.


\subsection{Aperçu de la syntaxe}
\label{\detokenize{intro:apercu-de-la-syntaxe}}
\sphinxAtStartPar
Voici un exemple simple de programme M minimal.

\begin{sphinxVerbatim}[commandchars=\\\{\}]
\PYG{c+c1}{\PYGZsh{} Fichier: test.m}

\PYG{c+c1}{\PYGZsh{} On définit un espace de variable global dans lequel nos variables}
\PYG{c+c1}{\PYGZsh{} seront stockées par défaut.}
\PYG{n}{espace\PYGZus{}variables} \PYG{n}{GLOBAL} \PYG{p}{:} \PYG{n}{par\PYGZus{}defaut}\PYG{p}{;}

\PYG{c+c1}{\PYGZsh{} On définit deux domaines obligatoires:}
\PYG{c+c1}{\PYGZsh{} * un domaine pour les règles;}
\PYG{c+c1}{\PYGZsh{} * un domaine pour les vérificateurs.}
\PYG{n}{domaine} \PYG{n}{regle} \PYG{n}{mon\PYGZus{}domaine\PYGZus{}de\PYGZus{}regle}\PYG{p}{:} \PYG{n}{par\PYGZus{}defaut}\PYG{p}{:} \PYG{n}{calculable}\PYG{p}{;}
\PYG{n}{domaine} \PYG{n}{verif} \PYG{n}{mon\PYGZus{}domaine\PYGZus{}de\PYGZus{}verifications}\PYG{p}{:} \PYG{n}{par\PYGZus{}defaut}\PYG{p}{;}

\PYG{c+c1}{\PYGZsh{} On définit une ou plusiseurs application pour nos règles.}
\PYG{n}{application} \PYG{n}{mon\PYGZus{}application}\PYG{p}{;}

\PYG{c+c1}{\PYGZsh{} On définit une cible, un ensemble d\PYGZsq{}instructions à exécuter.}
\PYG{n}{cible} \PYG{n}{hello\PYGZus{}world}\PYG{p}{:}
\PYG{n}{application} \PYG{p}{:} \PYG{n}{mon\PYGZus{}application}\PYG{p}{;}
\PYG{n}{afficher} \PYG{l+s+s2}{\PYGZdq{}}\PYG{l+s+s2}{Bonjour, monde!}\PYG{l+s+se}{\PYGZbs{}n}\PYG{l+s+s2}{\PYGZdq{}}\PYG{p}{;}
\end{sphinxVerbatim}

\sphinxAtStartPar
Cet exemple jouet n’utilise ni ne définit aucune variable, les domaines
et les espaces de variables ne sont pas utilisés ici mais ils seront nécessaires
plus tard.

\sphinxAtStartPar
Pour essayer notre exemple, nous allons créer un fichier “test.irj” avec le
contenu suivant :

\begin{sphinxVerbatim}[commandchars=\\\{\}]
\PYG{c+c1}{\PYGZsh{}NOM}
\PYG{n}{MON}\PYG{o}{\PYGZhy{}}\PYG{n}{TEST}
\PYG{c+c1}{\PYGZsh{}ENTREES\PYGZhy{}PRIMITIF}
\PYG{c+c1}{\PYGZsh{}CONTROLES\PYGZhy{}PRIMITIF}
\PYG{c+c1}{\PYGZsh{}RESULTATS\PYGZhy{}PRIMITIF}
\PYG{c+c1}{\PYGZsh{}ENTREES\PYGZhy{}CORRECTIF}
\PYG{c+c1}{\PYGZsh{}CONTROLES\PYGZhy{}CORRECTIF}
\PYG{c+c1}{\PYGZsh{}RESULTATS\PYGZhy{}CORRECTIF}
\PYG{c+c1}{\PYGZsh{}\PYGZsh{}}
\end{sphinxVerbatim}

\sphinxAtStartPar
Le détail de la syntaxe irj peut être sur la page dédiée : {\hyperref[\detokenize{syntax_irj:syntax-irj}]{\sphinxcrossref{\DUrole{std}{\DUrole{std-ref}{Les fichiers IRJ}}}}}

\sphinxAtStartPar
Enfin, on peut lancer mlang.

\begin{sphinxVerbatim}[commandchars=\\\{\}]
 \PYGZdl{} mlang \PYGZhy{}\PYGZhy{}without\PYGZus{}dfgip\PYGZus{}m test.m \PYGZhy{}A mon\PYGZus{}application \PYGZhy{}\PYGZhy{}mpp\PYGZus{}function hello\PYGZus{}world \PYGZhy{}\PYGZhy{}run\PYGZus{}test test.irj
 Parsing: completed!
 Bonjour, monde!
 [RESULT] test.irj
 [RESULT] No failure!
 [RESULT] Test passed!
\end{sphinxVerbatim}


\subsection{Définition de variables}
\label{\detokenize{intro:definition-de-variables}}
\sphinxAtStartPar
Les variables sont divisées en deux catégories.
\begin{itemize}
\item {} 
\sphinxAtStartPar
Les variables \sphinxcode{\sphinxupquote{saisie}}s sont les variables d’entrées du
programme M.

\item {} 
\sphinxAtStartPar
Les variables \sphinxcode{\sphinxupquote{calculee}}s sont les variables sur lesquelles seront
calculées les données.

\end{itemize}

\sphinxAtStartPar
Chacune de ces catégories principales doivent être dotées d’attributs,
un ensemble d’entiers constants pour une variable donnée.
Ainsi, on peut rajouter à notre programme test les lignes suivantes :

\begin{sphinxVerbatim}[commandchars=\\\{\}]
\PYG{n}{variable} \PYG{n}{saisie} \PYG{p}{:} \PYG{n}{attribut} \PYG{n}{mon\PYGZus{}attribut}\PYG{p}{;}
\PYG{n}{variable} \PYG{n}{calculee} \PYG{p}{:} \PYG{n}{attribut} \PYG{n}{mon\PYGZus{}attribut}\PYG{p}{;}
\PYG{n}{X} \PYG{p}{:} \PYG{n}{saisie} \PYG{n}{mon\PYGZus{}attribut} \PYG{o}{=} \PYG{l+m+mi}{0} \PYG{n}{alias} \PYG{n}{ALIAS\PYGZus{}DE\PYGZus{}X} \PYG{p}{:} \PYG{l+s+s2}{\PYGZdq{}}\PYG{l+s+s2}{Cette variable s}\PYG{l+s+s2}{\PYGZsq{}}\PYG{l+s+s2}{appelle X}\PYG{l+s+s2}{\PYGZdq{}}\PYG{p}{;}
\PYG{n}{Y} \PYG{p}{:} \PYG{n}{calculee} \PYG{n}{mon\PYGZus{}attribut} \PYG{o}{=} \PYG{l+m+mi}{1} \PYG{p}{:} \PYG{l+s+s2}{\PYGZdq{}}\PYG{l+s+s2}{Cette variable s}\PYG{l+s+s2}{\PYGZsq{}}\PYG{l+s+s2}{appelle Y}\PYG{l+s+s2}{\PYGZdq{}}\PYG{p}{;}
\end{sphinxVerbatim}

\sphinxAtStartPar
Notez que la variable \sphinxcode{\sphinxupquote{X}} a un alias \sphinxcode{\sphinxupquote{ALIAS\_DE\_X}}.
Toutes les variables saisies doivent être parées d’un alias qui peut être
utilisé de la même façon que son nom original.
Cet alias existe initialement pour faire le lien entre la variable utilisée
dans le M (nom original) et le code de la variable dans le formulaire de
déclaration de l’impot sur le revenu tel qu’on le retrouve aujourd’hui sur
le site de déclaration (1AP, 8TV, …).

\sphinxAtStartPar
Ajoutons une nouvelle cible à notre calcul :

\begin{sphinxVerbatim}[commandchars=\\\{\}]
\PYG{n}{cible} \PYG{n}{hello\PYGZus{}world2}\PYG{p}{:}
\PYG{n}{application} \PYG{p}{:} \PYG{n}{mon\PYGZus{}application}\PYG{p}{;}
\PYG{n}{afficher} \PYG{l+s+s2}{\PYGZdq{}}\PYG{l+s+s2}{Bonjour, monde, X = }\PYG{l+s+s2}{\PYGZdq{}}\PYG{p}{;}
\PYG{n}{afficher} \PYG{p}{(}\PYG{n}{X}\PYG{p}{)}\PYG{p}{;}
\PYG{n}{afficher} \PYG{l+s+s2}{\PYGZdq{}}\PYG{l+s+s2}{ !}\PYG{l+s+se}{\PYGZbs{}n}\PYG{l+s+s2}{\PYGZdq{}}\PYG{p}{;}
\PYG{n}{Y} \PYG{o}{=} \PYG{n}{X} \PYG{o}{+} \PYG{l+m+mi}{1}\PYG{p}{;}
\PYG{n}{afficher} \PYG{l+s+s2}{\PYGZdq{}}\PYG{l+s+s2}{Y = }\PYG{l+s+s2}{\PYGZdq{}}\PYG{p}{;}
\PYG{n}{afficher} \PYG{p}{(}\PYG{n}{Y}\PYG{p}{)}\PYG{p}{;}
\PYG{n}{afficher} \PYG{l+s+s2}{\PYGZdq{}}\PYG{l+s+s2}{ !}\PYG{l+s+se}{\PYGZbs{}n}\PYG{l+s+s2}{\PYGZdq{}}\PYG{p}{;}
\end{sphinxVerbatim}

\sphinxAtStartPar
Et lançons le calcul :

\begin{sphinxVerbatim}[commandchars=\\\{\}]
 \PYGZdl{} mlang \PYGZhy{}\PYGZhy{}without\PYGZus{}dfgip\PYGZus{}m test.m \PYGZhy{}A mon\PYGZus{}application \PYGZhy{}\PYGZhy{}mpp\PYGZus{}function hello\PYGZus{}world2 \PYGZhy{}\PYGZhy{}run\PYGZus{}test test.irj
Parsing: completed!
Bonjour, monde, X = indefini !
Y = 1 !
[RESULT] test.irj
[RESULT] No failure!
[RESULT] Test passed!
\end{sphinxVerbatim}

\sphinxAtStartPar
Pour comprendre la valeur finale de Y, référez\sphinxhyphen{}vous à la
section {\hyperref[\detokenize{syntax:valeurs}]{\sphinxcrossref{\DUrole{std}{\DUrole{std-ref}{Les valeurs}}}}}.


\subsection{Règles de calcul}
\label{\detokenize{intro:regles-de-calcul}}
\sphinxAtStartPar
Les règles en M sont des unités de calcul de variables associées à une ou
plusieurs application et optionnellement à un domaine de règles.
Elles sont composées d’une successions d’affectations.
Voici la définition de deux règles simples que nous rajoutons à notre fichier test :

\begin{sphinxVerbatim}[commandchars=\\\{\}]
\PYG{n}{Z} \PYG{p}{:} \PYG{n}{calculee} \PYG{n}{mon\PYGZus{}attribut} \PYG{o}{=} \PYG{l+m+mi}{1} \PYG{p}{:} \PYG{l+s+s2}{\PYGZdq{}}\PYG{l+s+s2}{Cette variable s}\PYG{l+s+s2}{\PYGZsq{}}\PYG{l+s+s2}{appelle Z}\PYG{l+s+s2}{\PYGZdq{}}\PYG{p}{;}

\PYG{n}{regle} \PYG{n}{mon\PYGZus{}domaine\PYGZus{}de\PYGZus{}regles} \PYG{l+m+mi}{1}\PYG{p}{:}
\PYG{n}{application} \PYG{p}{:} \PYG{n}{mon\PYGZus{}application}\PYG{p}{;}
\PYG{n}{Z} \PYG{o}{=} \PYG{n}{Y} \PYG{o}{+} \PYG{l+m+mi}{1}\PYG{p}{;}

\PYG{n}{regle} \PYG{l+m+mi}{2}\PYG{p}{:}
\PYG{n}{application} \PYG{p}{:} \PYG{n}{mon\PYGZus{}application}\PYG{p}{;}
\PYG{n}{Y} \PYG{o}{=} \PYG{n}{X} \PYG{o}{+} \PYG{l+m+mi}{1}\PYG{p}{;}
\end{sphinxVerbatim}

\sphinxAtStartPar
La première règle est associée au domaine \sphinxcode{\sphinxupquote{mon\_domaine\_de\_regles}} tandis que la
seconde n’étant pas spécifié, sera associée au domaine de règle par défaut.
Dans notre exemple, il s’agissait également de \sphinxcode{\sphinxupquote{mon\_domaine\_de\_regles}}.

\sphinxAtStartPar
Le calcul d’un domaine correspond au calcul de l’ensemble de ses règles.
L’ordre d’application des règles dépend de l’ordre d’affectation des variables.
Dans notre cas, \sphinxcode{\sphinxupquote{X}} est une entrée dont \sphinxcode{\sphinxupquote{Y}} dépend (règle 2), et \sphinxcode{\sphinxupquote{Z}} dépend de
\sphinxcode{\sphinxupquote{Y}}.

\sphinxAtStartPar
Par conséquent, la règle 2 sera appliquée avant la règle 1.
On peut ainsi rajouter une cible qui calcule le domaine de règles :

\begin{sphinxVerbatim}[commandchars=\\\{\}]
\PYG{n}{cible} \PYG{n}{calc\PYGZus{}test}\PYG{p}{:}
\PYG{n}{application} \PYG{p}{:} \PYG{n}{mon\PYGZus{}application}\PYG{p}{;}
\PYG{n}{calculer} \PYG{n}{domaine} \PYG{n}{mon\PYGZus{}domaine\PYGZus{}de\PYGZus{}regles}\PYG{p}{;}
\PYG{n}{afficher} \PYG{l+s+s2}{\PYGZdq{}}\PYG{l+s+s2}{X = }\PYG{l+s+s2}{\PYGZdq{}}\PYG{p}{;}
\PYG{n}{afficher} \PYG{p}{(}\PYG{n}{X}\PYG{p}{)}\PYG{p}{;}
\PYG{n}{afficher} \PYG{l+s+s2}{\PYGZdq{}}\PYG{l+s+se}{\PYGZbs{}n}\PYG{l+s+s2}{Y = }\PYG{l+s+s2}{\PYGZdq{}}\PYG{p}{;}
\PYG{n}{afficher} \PYG{p}{(}\PYG{n}{Y}\PYG{p}{)}\PYG{p}{;}
\PYG{n}{afficher} \PYG{l+s+s2}{\PYGZdq{}}\PYG{l+s+se}{\PYGZbs{}n}\PYG{l+s+s2}{Z = }\PYG{l+s+s2}{\PYGZdq{}}\PYG{p}{;}
\PYG{n}{afficher} \PYG{p}{(}\PYG{n}{Z}\PYG{p}{)}\PYG{p}{;}
\PYG{n}{afficher} \PYG{l+s+s2}{\PYGZdq{}}\PYG{l+s+se}{\PYGZbs{}n}\PYG{l+s+s2}{\PYGZdq{}}\PYG{p}{;}
\end{sphinxVerbatim}

\sphinxAtStartPar
Et lancer le calcul :

\begin{sphinxVerbatim}[commandchars=\\\{\}]
 \PYGZdl{} mlang \PYGZhy{}\PYGZhy{}without\PYGZus{}dfgip\PYGZus{}m test.m \PYGZhy{}A mon\PYGZus{}application \PYGZhy{}\PYGZhy{}mpp\PYGZus{}function calc\PYGZus{}test \PYGZhy{}\PYGZhy{}run\PYGZus{}test test.irj
 Parsing: completed!
 X = indefini
 Y = 1
 Z = 2
 [RESULT] test.irj
 [RESULT] No failure!
 [RESULT] Test passed!
\end{sphinxVerbatim}

\sphinxstepscope


\section{La syntaxe du M}
\label{\detokenize{syntax:la-syntaxe-du-m}}\label{\detokenize{syntax:syntax}}\label{\detokenize{syntax::doc}}

\subsection{Programme M et applications}
\label{\detokenize{syntax:programme-m-et-applications}}
\sphinxAtStartPar
Un programme M est formé d’une suite de caractères codés sur 8 bits.
Les 128 premiers codes de caractères correspondent aux codes ASCII.
Un programme M est constitué d’une suite d’éléments parmi lesquels on compte :
\begin{itemize}
\item {} 
\sphinxAtStartPar
des déclarations d’applications ;

\item {} 
\sphinxAtStartPar
des définitions de constantes ;

\item {} 
\sphinxAtStartPar
des déclarations d’enchaîneurs ;

\item {} 
\sphinxAtStartPar
des définitions de catégories de variables ;

\item {} 
\sphinxAtStartPar
des déclarations de variables ;

\item {} 
\sphinxAtStartPar
des déclarations d’erreurs ;

\item {} 
\sphinxAtStartPar
des déclarations de fonctions externes ;

\item {} 
\sphinxAtStartPar
des définitions de domaines de règles ;

\item {} 
\sphinxAtStartPar
des définitions de domaines de vérifications ;

\item {} 
\sphinxAtStartPar
des déclarations de sorties ;

\item {} 
\sphinxAtStartPar
des règles ;

\item {} 
\sphinxAtStartPar
des vérifications ;

\item {} 
\sphinxAtStartPar
des fonctions ;

\item {} 
\sphinxAtStartPar
des cibles.

\end{itemize}

\sphinxAtStartPar
Un programme M peut définir plusieurs applications.
Pour être traité (compilation, interprétation, etc.), il nécessite la donnée
d’une liste d’applications.
Cette liste est typiquement fournie comme argument du compilateur, de
l’interpréteur ou de tout autre outil de traitement.
Cette liste sera appelée la liste des applications sélectionnées.


\subsection{Définitions liminaires}
\label{\detokenize{syntax:definitions-liminaires}}
\sphinxAtStartPar
Les formes de Backus\sphinxhyphen{}Naur étendues sont utilisées pour préciser la morphologie
du langage.
Un lexème est une suite de caractères et chaque forme est susceptible d’en
reconnaître un ensemble.

\sphinxAtStartPar
Ces formes sont étendues avec les notations suivantes :
\begin{itemize}
\item {} 
\sphinxAtStartPar
\sphinxcode{\sphinxupquote{\textless{}non\sphinxhyphen{}terminal:identifiant\textgreater{}}} représente le non\sphinxhyphen{}terminal, pour lequel la
chaîne de caractères qu’il reconnaît est représentée par identifiant ;

\item {} 
\sphinxAtStartPar
\sphinxcode{\sphinxupquote{forme 1 \textasciicircum{} … \textasciicircum{} forme n}} qui représente une suite de lexèmes correspondants
aux formes 1 à n, dans n’importe quel ordre;

\item {} 
\sphinxAtStartPar
et : \sphinxcode{\sphinxupquote{(forme 0) séparateur …}} qui représente une liste non\sphinxhyphen{}vide de lexèmes
reeconnus par la forme 0, tous séparés par des lexèmes correspondants au
séparateur.

\end{itemize}

\sphinxAtStartPar
Les lexèmes suivants sont les mots réservés :
\sphinxstylestrong{BOOLEEN},
\sphinxstylestrong{DATE\_AAAA},
\sphinxstylestrong{DATE\_JJMMAAAA},
\sphinxstylestrong{DATE\_MM},
\sphinxstylestrong{ENTIER},
\sphinxstylestrong{REEL},
\sphinxstylestrong{afficher},
\sphinxstylestrong{afficher\_erreur},
\sphinxstylestrong{aiguillage},
\sphinxstylestrong{ajouter},
\sphinxstylestrong{alias},
\sphinxstylestrong{alors},
\sphinxstylestrong{anomalie},
\sphinxstylestrong{application},
\sphinxstylestrong{apres},
\sphinxstylestrong{argument},
\sphinxstylestrong{arranger\_evenements},
\sphinxstylestrong{attribut},
\sphinxstylestrong{autorise},
\sphinxstylestrong{avec},
\sphinxstylestrong{base},
\sphinxstylestrong{calculable},
\sphinxstylestrong{calculee},
\sphinxstylestrong{calculer},
\sphinxstylestrong{cas},
\sphinxstylestrong{categorie},
\sphinxstylestrong{champ\_evenement},
\sphinxstylestrong{cible},
\sphinxstylestrong{const},
\sphinxstylestrong{dans},
\sphinxstylestrong{dans\_domaine},
\sphinxstylestrong{discordance},
\sphinxstylestrong{domaine},
\sphinxstylestrong{enchaineur},
\sphinxstylestrong{entre},
\sphinxstylestrong{erreur},
\sphinxstylestrong{espace},
\sphinxstylestrong{espace\_variables},
\sphinxstylestrong{et},
\sphinxstylestrong{evenement},
\sphinxstylestrong{evenements},
\sphinxstylestrong{exporte\_erreurs},
\sphinxstylestrong{faire},
\sphinxstylestrong{filtrer},
\sphinxstylestrong{finalise\_erreurs},
\sphinxstylestrong{finquand},
\sphinxstylestrong{finsi},
\sphinxstylestrong{fonction},
\sphinxstylestrong{increment},
\sphinxstylestrong{indefini},
\sphinxstylestrong{indenter},
\sphinxstylestrong{informative},
\sphinxstylestrong{iterer},
\sphinxstylestrong{leve\_erreur},
\sphinxstylestrong{meme\_variable},
\sphinxstylestrong{nb\_anomalies},
\sphinxstylestrong{nb\_bloquantes},
\sphinxstylestrong{nb\_categorie},
\sphinxstylestrong{nb\_discordances},
\sphinxstylestrong{nb\_informatives},
\sphinxstylestrong{neant},
\sphinxstylestrong{nettoie\_erreurs},
\sphinxstylestrong{nettoie\_erreurs\_finalisees},
\sphinxstylestrong{nom},
\sphinxstylestrong{non},
\sphinxstylestrong{numero\_compl},
\sphinxstylestrong{numero\_verif},
\sphinxstylestrong{ou},
\sphinxstylestrong{par\_defaut},
\sphinxstylestrong{pour},
\sphinxstylestrong{puis\_quand},
\sphinxstylestrong{quand},
\sphinxstylestrong{reference},
\sphinxstylestrong{regle},
\sphinxstylestrong{restaurer},
\sphinxstylestrong{restituee},
\sphinxstylestrong{resultat},
\sphinxstylestrong{saisie},
\sphinxstylestrong{si},
\sphinxstylestrong{sinon},
\sphinxstylestrong{sinon\_si},
\sphinxstylestrong{sortie},
\sphinxstylestrong{specialise},
\sphinxstylestrong{stop},
\sphinxstylestrong{tableau},
\sphinxstylestrong{taille},
\sphinxstylestrong{trier},
\sphinxstylestrong{type},
\sphinxstylestrong{un},
\sphinxstylestrong{valeur}
\sphinxstylestrong{variable},
\sphinxstylestrong{verif},
\sphinxstylestrong{verifiable},
et \sphinxstylestrong{verifier}.

\sphinxAtStartPar
Les lexèmes numériques sont définis comme suit :
\begin{itemize}
\item {} 
\sphinxAtStartPar
\sphinxcode{\sphinxupquote{\textless{}naturel\textgreater{} ::= {[}0\sphinxhyphen{}9{]} {[}0\sphinxhyphen{}9 \_{]}*}}

\item {} 
\sphinxAtStartPar
\sphinxcode{\sphinxupquote{\textless{}réel\textgreater{} ::= \textless{}naturel\textgreater{} (. \textless{}naturel\textgreater{})?}}

\item {} 
\sphinxAtStartPar
\sphinxcode{\sphinxupquote{\textless{}symbole\textgreater{}}} est la forme \sphinxcode{\sphinxupquote{{[}a\sphinxhyphen{}z A\sphinxhyphen{}Z 0\sphinxhyphen{}9 \_{]}+}} dont sont exclus \sphinxcode{\sphinxupquote{\textless{}naturel\textgreater{}}}s et
les mots réservés ;

\item {} 
\sphinxAtStartPar
\sphinxcode{\sphinxupquote{\textless{}variable\textgreater{}}} est un \sphinxcode{\sphinxupquote{\textless{}symbole\textgreater{}}} que l’on distingue pour représenter les
mots pouvant prendre le nom d’une constante.

\end{itemize}

\sphinxAtStartPar
On définit également les atomes numériques, prenant une valeur soit par un
nombre littéral, soit par un symbole représentant une constante ou une
variable :
\begin{itemize}
\item {} 
\sphinxAtStartPar
\sphinxcode{\sphinxupquote{\textless{}atome naturel\textgreater{} ::= \textless{}naturel\textgreater{} | \textless{}variable\textgreater{}}}

\item {} 
\sphinxAtStartPar
\sphinxcode{\sphinxupquote{\textless{}atome reel\textgreater{} ::= \textless{}réel\textgreater{} | \textless{}variable\textgreater{}}}

\end{itemize}

\sphinxAtStartPar
Les chaînes de caractères littérales ont la forme suivante :
\begin{itemize}
\item {} 
\sphinxAtStartPar
\sphinxcode{\sphinxupquote{\textless{}chaîne\textgreater{} :: = " {[}\textasciicircum{} "{]}* "}}

\end{itemize}


\subsubsection{Déclaration une d’application}
\label{\detokenize{syntax:declaration-une-d-application}}
\sphinxAtStartPar
Les \sphinxstyleemphasis{déclarations d’applications} ont la forme suivante :

\begin{sphinxVerbatim}[commandchars=\\\{\}]
\PYG{n}{application} \PYG{o}{\PYGZlt{}}\PYG{n}{symbole}\PYG{o}{\PYGZgt{}}\PYG{p}{;}
\end{sphinxVerbatim}

\sphinxAtStartPar
où \sphinxcode{\sphinxupquote{\textless{}symbole\textgreater{}}} est le nom de l’application
déclarée.


\subsubsection{Définition d’une constante}
\label{\detokenize{syntax:definition-d-une-constante}}
\sphinxAtStartPar
Les \sphinxstyleemphasis{définitions de constantes} ont la forme suivante :

\begin{sphinxVerbatim}[commandchars=\\\{\}]
\PYG{o}{\PYGZlt{}}\PYG{n}{symbole}\PYG{o}{\PYGZgt{}} \PYG{p}{:} \PYG{n}{const} \PYG{o}{=} \PYG{o}{\PYGZlt{}}\PYG{n}{atome} \PYG{n}{réel}\PYG{o}{\PYGZgt{}} \PYG{p}{;}
\end{sphinxVerbatim}

\sphinxAtStartPar
Le symbole est le nom de la constante.


\subsubsection{Déclaration d’un enchaîneur}
\label{\detokenize{syntax:declaration-d-un-enchaineur}}
\sphinxAtStartPar
Les \sphinxstyleemphasis{déclarations d’enchaîneurs} ont la forme suivante :

\begin{sphinxVerbatim}[commandchars=\\\{\}]
\PYG{n}{enchaineur} \PYG{o}{\PYGZlt{}}\PYG{n}{symbole}\PYG{o}{\PYGZgt{}} \PYG{p}{:} \PYG{o}{\PYGZlt{}}\PYG{n}{symboles}\PYG{o}{\PYGZgt{}} \PYG{p}{;}
\end{sphinxVerbatim}

\sphinxAtStartPar
avec :
\sphinxcode{\sphinxupquote{\textless{}symboles\textgreater{} ::= \textless{}symbole\textgreater{} , …}} représente l’ensemble des applications incluant cet enchaîneur.
Le \sphinxcode{\sphinxupquote{\textless{}symbole\textgreater{}}} de l’enchaineur est le nom de l’application déclarée.


\subsubsection{Définition d’une catégorie de variables}
\label{\detokenize{syntax:definition-d-une-categorie-de-variables}}
\sphinxAtStartPar
Les \sphinxstyleemphasis{définitions de catégories de variables} ont la forme suivante :

\begin{sphinxVerbatim}[commandchars=\\\{\}]
variable (saisie | calculee) \PYGZlt{}nom\PYGZgt{} (: attribut \PYGZlt{}symboles\PYGZgt{})? ;
\end{sphinxVerbatim}

\sphinxAtStartPar
avec :
\begin{itemize}
\item {} 
\sphinxAtStartPar
\sphinxcode{\sphinxupquote{\textless{}nom\textgreater{} ::= \textless{}symbole\textgreater{}+}} est le nom de la catégorie de variables;

\item {} 
\sphinxAtStartPar
\sphinxcode{\sphinxupquote{\textless{}symboles\textgreater{} ::= \textless{}symbole\textgreater{} , …}} représente l’ensemble de ses attributs.

\end{itemize}


\subsubsection{Déclaration d’une variable}
\label{\detokenize{syntax:declaration-d-une-variable}}
\sphinxAtStartPar
La forme des *types de variables *est comme suit :

\begin{sphinxVerbatim}[commandchars=\\\{\}]
\PYG{o}{\PYGZlt{}}\PYG{n+nb}{type} \PYG{n}{variable}\PYG{o}{\PYGZgt{}} \PYG{p}{:}\PYG{o}{:=} \PYG{n}{BOOLEEN} \PYG{o}{|} \PYG{n}{DATE\PYGZus{}AAAA} \PYG{o}{|} \PYG{n}{DATE\PYGZus{}JJMMAAAA} \PYG{o}{|} \PYG{n}{DATE\PYGZus{}MM} \PYG{o}{|} \PYG{n}{ENTIER} \PYG{o}{|} \PYG{n}{REEL}
\end{sphinxVerbatim}


\paragraph{Variable saisie}
\label{\detokenize{syntax:variable-saisie}}
\sphinxAtStartPar
Les \sphinxstyleemphasis{déclarations de variables saisies} ont la forme suivante :

\begin{sphinxVerbatim}[commandchars=\\\{\}]
\PYG{o}{\PYGZlt{}}\PYG{n}{symbole}\PYG{o}{\PYGZgt{}} \PYG{p}{:} \PYG{n}{saisie} \PYG{o}{\PYGZlt{}}\PYG{n}{catégorie}\PYG{o}{\PYGZgt{}} \PYG{o}{\PYGZlt{}}\PYG{n}{attributs}\PYG{o}{\PYGZgt{}} \PYG{n}{alias} \PYG{o}{\PYGZlt{}}\PYG{n}{symbole}\PYG{o}{\PYGZgt{}} \PYG{p}{:} \PYG{o}{\PYGZlt{}}\PYG{n}{chaîne}\PYG{o}{\PYGZgt{}} \PYG{n+nb}{type} \PYG{o}{\PYGZlt{}}\PYG{n+nb}{type} \PYG{n}{variable}\PYG{o}{\PYGZgt{}} \PYG{p}{;}
\end{sphinxVerbatim}

\sphinxAtStartPar
avec :
\begin{itemize}
\item {} 
\sphinxAtStartPar
\sphinxcode{\sphinxupquote{\textless{}catégorie\textgreater{} ::= \textless{}symbole\textgreater{}+}}

\item {} 
\sphinxAtStartPar
\sphinxcode{\sphinxupquote{\textless{}attributs\textgreater{} ::= \textless{}attribut\textgreater{} restituee? \textless{}attribut\textgreater{}}}

\item {} 
\sphinxAtStartPar
\sphinxcode{\sphinxupquote{\textless{}attribut\textgreater{} ::= \textless{}symbole\textgreater{} = \textless{}atome naturel\textgreater{}}}

\end{itemize}

\sphinxAtStartPar
Le \sphinxcode{\sphinxupquote{\textless{}symbole\textgreater{}}} est le nom de la variable.


\paragraph{Variable calculée}
\label{\detokenize{syntax:variable-calculee}}
\sphinxAtStartPar
Les \sphinxstyleemphasis{déclarations de variables calculées} ont la forme suivante :

\begin{sphinxVerbatim}[commandchars=\\\{\}]
\PYGZlt{}symbole\PYGZgt{}
: (table [ \PYGZlt{}atome naturel\PYGZgt{} ])? calculee (base? \PYGZca{} restituee?)
: \PYGZlt{}chaîne\PYGZgt{} type \PYGZlt{}type variable\PYGZgt{} ;
\end{sphinxVerbatim}

\sphinxAtStartPar
Le  est le nom de la variable. Si c’est un tableau, l”
est sa taille.


\subsubsection{Déclaration d’une erreur}
\label{\detokenize{syntax:declaration-d-une-erreur}}
\sphinxAtStartPar
Les \sphinxstyleemphasis{déclarations d’erreurs} ont la forme suivante :

\begin{sphinxVerbatim}[commandchars=\\\{\}]
\PYG{o}{\PYGZlt{}}\PYG{n}{symbole}\PYG{o}{\PYGZgt{}} \PYG{p}{:} \PYG{o}{\PYGZlt{}}\PYG{n+nb}{type} \PYG{n}{erreur}\PYG{o}{\PYGZgt{}} \PYG{p}{:} \PYG{o}{\PYGZlt{}}\PYG{n}{chaîne}\PYG{o}{\PYGZgt{}} \PYG{p}{:} \PYG{o}{\PYGZlt{}}\PYG{n}{chaîne}\PYG{o}{\PYGZgt{}} \PYG{p}{:} \PYG{o}{\PYGZlt{}}\PYG{n}{chaîne}\PYG{o}{\PYGZgt{}} \PYG{p}{:} \PYG{o}{\PYGZlt{}}\PYG{n}{chaîne}\PYG{o}{\PYGZgt{}} \PYG{p}{:} \PYG{o}{\PYGZlt{}}\PYG{n}{chaîne}\PYG{o}{\PYGZgt{}} \PYG{p}{;}
\end{sphinxVerbatim}

\sphinxAtStartPar
avec :
\begin{itemize}
\item {} 
\sphinxAtStartPar
\sphinxcode{\sphinxupquote{\textless{}type erreur\textgreater{} ::= anomalie | discordance | informative}}

\end{itemize}

\sphinxAtStartPar
Le symbole est le nom de l’erreur.


\subsubsection{Déclaration d’une fonction externe}
\label{\detokenize{syntax:declaration-d-une-fonction-externe}}
\sphinxAtStartPar
Les \sphinxstyleemphasis{déclarations de fonctions} externes ont la forme suivante :

\begin{sphinxVerbatim}[commandchars=\\\{\}]
\PYG{o}{\PYGZlt{}}\PYG{n}{symbole}\PYG{o}{\PYGZgt{}} \PYG{p}{:} \PYG{n}{fonction} \PYG{o}{\PYGZlt{}}\PYG{n}{atome} \PYG{n}{naturel}\PYG{o}{\PYGZgt{}} \PYG{p}{;}
\end{sphinxVerbatim}

\sphinxAtStartPar
Le \sphinxcode{\sphinxupquote{\textless{}symbole\textgreater{}}} est le nom de la fonction, et l” son arité.


\subsubsection{Définition d’un domaine de règles}
\label{\detokenize{syntax:definition-d-un-domaine-de-regles}}
\sphinxAtStartPar
Les \sphinxstyleemphasis{défintions de domaines de règles} ont la forme suivante :

\begin{sphinxVerbatim}[commandchars=\\\{\}]
\PYG{n}{domaine} \PYG{n}{regle} \PYG{o}{\PYGZlt{}}\PYG{n}{nom}\PYG{o}{\PYGZgt{}} \PYG{o}{\PYGZlt{}}\PYG{n}{paramètres}\PYG{o}{\PYGZgt{}} \PYG{p}{;}
\end{sphinxVerbatim}

\sphinxAtStartPar
avec :
\begin{itemize}
\item {} 
\sphinxAtStartPar
\sphinxcode{\sphinxupquote{\textless{}nom\textgreater{} ::= \textless{}symbole\textgreater{}+}} est le nom du domaine de règles.

\item {} 
\sphinxAtStartPar
\textasciigrave{}\textless{}paramètres\textgreater{} ::= (: specialise )? \textasciicircum{} (: calculable)? \textasciicircum{} (: par\_defaut)?

\item {} 
\sphinxAtStartPar
 ::=  , …\textasciigrave{} représente l’ensemble des domaines de règles que spécialise ce domaine.

\end{itemize}


\subsubsection{Définition d’un domaine de vérifications}
\label{\detokenize{syntax:definition-d-un-domaine-de-verifications}}
\sphinxAtStartPar
Les \sphinxstyleemphasis{définitions de domaines de vérifications} ont la forme suivante :

\begin{sphinxVerbatim}[commandchars=\\\{\}]
\PYG{n}{domaine} \PYG{n}{verif} \PYG{o}{\PYGZlt{}}\PYG{n}{nom}\PYG{o}{\PYGZgt{}} \PYG{o}{\PYGZlt{}}\PYG{n}{paramètres}\PYG{o}{\PYGZgt{}} \PYG{p}{;}
\end{sphinxVerbatim}

\sphinxAtStartPar
avec :
\begin{itemize}
\item {} 
\sphinxAtStartPar
\sphinxcode{\sphinxupquote{\textless{}nom\textgreater{} ::= \textless{}symbole\textgreater{}+}} est le nom domaine de vérifications que spécialise ce domaine.

\item {} 
\sphinxAtStartPar
\sphinxcode{\sphinxupquote{\textless{}paramètres\textgreater{} ::= (: specialise \textless{}noms\textgreater{})? \textasciicircum{} (: autorise \textless{}catvars\textgreater{})? \textasciicircum{} (: par\_defaut)? \textasciicircum{} (: verifiable)?}}

\item {} 
\sphinxAtStartPar
\sphinxcode{\sphinxupquote{\textless{}noms\textgreater{} ::= \textless{}nom\textgreater{} , …}}

\item {} 
\sphinxAtStartPar
\sphinxcode{\sphinxupquote{\textless{}catvars\textgreater{} ::= * | saisie (* | \textless{}nom\textgreater{}) | calculee (* | base)?}} représente les catégories des variables autorisées à êtres utilisées dans les vérifications de ce domaine.

\end{itemize}


\subsubsection{Déclaration d’une sortie}
\label{\detokenize{syntax:declaration-d-une-sortie}}
\sphinxAtStartPar
Les \sphinxstyleemphasis{déclarations de sorties} ont la forme suivante :

\begin{sphinxVerbatim}[commandchars=\\\{\}]
\PYG{n}{sortie} \PYG{p}{(} \PYG{o}{\PYGZlt{}}\PYG{n}{symbole}\PYG{o}{\PYGZgt{}} \PYG{p}{)} \PYG{p}{;}
\end{sphinxVerbatim}


\subsubsection{Indices}
\label{\detokenize{syntax:indices}}
\sphinxAtStartPar
Les \sphinxstyleemphasis{indices} ont la forme suivante :

\begin{sphinxVerbatim}[commandchars=\\\{\}]
\PYGZlt{}indices\PYGZgt{} ::= \PYGZlt{}indice\PYGZgt{} ; …
\PYGZlt{}indice\PYGZgt{} ::= \PYGZlt{}minuscule\PYGZgt{} = \PYGZlt{}intervalle\PYGZgt{} , …
\end{sphinxVerbatim}

\sphinxAtStartPar
avec :
\begin{itemize}
\item {} 
\sphinxAtStartPar
\sphinxcode{\sphinxupquote{\textless{}minuscule\textgreater{} ::= {[}a\sphinxhyphen{}z{]}}}

\item {} 
\sphinxAtStartPar
\sphinxcode{\sphinxupquote{\textless{}majuscule\textgreater{} ::= {[}A\sphinxhyphen{}Z{]}}}

\item {} 
\sphinxAtStartPar
\sphinxcode{\sphinxupquote{\textless{}intervalle\textgreater{} ::= \textless{}majuscule\textgreater{}+ | \textless{}majuscule\textgreater{} .. \textless{}majuscule\textgreater{} | \textless{}naturel\textgreater{} (.. \textless{}naturel\textgreater{} | \sphinxhyphen{} \textless{}variable\textgreater{})?}}

\end{itemize}


\subsubsection{Expressions}
\label{\detokenize{syntax:expressions}}
\sphinxAtStartPar
Les \sphinxstyleemphasis{expressions atomiques} ont la forme suivante :

\begin{sphinxVerbatim}[commandchars=\\\{\}]
\PYGZlt{}atomique booléen\PYGZgt{} ::=
( \PYGZlt{}expression booléenne\PYGZgt{} )
| \PYGZlt{}expression numérique\PYGZgt{} (\PYGZlt{}= | \PYGZgt{}= | \PYGZlt{} | \PYGZgt{} | != | =) \PYGZlt{}expression numérique\PYGZgt{}
| \PYGZlt{}expression numérique\PYGZgt{} non? dans ( \PYGZlt{}intervalles numériques\PYGZgt{} )
| (present | non\PYGZus{}present) ( \PYGZlt{}symbole\PYGZgt{} ([ \PYGZlt{}expression numérique\PYGZgt{} ])? )

\PYGZlt{}atomique numérique\PYGZgt{} ::=
| \PYGZlt{}expression numerique\PYGZgt{}
| \PYGZlt{}atome réel\PYGZgt{}
| \PYGZlt{}symbole\PYGZgt{} [ \PYGZlt{}expression numérique\PYGZgt{} ]
| \PYGZlt{}symbole\PYGZgt{} ( (\PYGZlt{}expression numerique\PYGZgt{} , …)? )
| somme ( \PYGZlt{}indices\PYGZgt{} : \PYGZlt{}expression numerique\PYGZgt{} )
| si ( \PYGZlt{}expression boolenne\PYGZgt{} )
  alors ( \PYGZlt{}expression numerique\PYGZgt{} )
  (( sinon \PYGZlt{}expression numerique\PYGZgt{} ))?
  finsi
\end{sphinxVerbatim}

\sphinxAtStartPar
avec :
\begin{itemize}
\item {} 
\sphinxAtStartPar
\sphinxcode{\sphinxupquote{\textless{}expression ou\textgreater{} ::= \textless{}expression et\textgreater{} | \textless{}expression ou\textgreater{} ou \textless{}expression ou\textgreater{}}}

\item {} 
\sphinxAtStartPar
\sphinxcode{\sphinxupquote{\textless{}expression et\textgreater{} ::= \textless{}expression non\textgreater{} | \textless{}expression et\textgreater{} et \textless{}expression et\textgreater{}}}

\item {} 
\sphinxAtStartPar
\sphinxcode{\sphinxupquote{\textless{}expression non\textgreater{} ::= \textless{}atomique booléen\textgreater{} | non \textless{}expression non\textgreater{}}}

\item {} 
\sphinxAtStartPar
\sphinxcode{\sphinxupquote{\textless{}expression + \sphinxhyphen{}\textgreater{} ::= \textless{}expression * /\textgreater{} | \textless{}expression + \sphinxhyphen{}\textgreater{} (+ | \sphinxhyphen{}) \textless{}expression + \sphinxhyphen{}\textgreater{}}}

\item {} 
\sphinxAtStartPar
\sphinxcode{\sphinxupquote{\textless{}expression * /\textgreater{} ::= \textless{}expression \sphinxhyphen{}\textgreater{} | \textless{}expression * /\textgreater{} (* | /) \textless{}expression * /\textgreater{}}}

\item {} 
\sphinxAtStartPar
\sphinxcode{\sphinxupquote{\textless{}expression \sphinxhyphen{}\textgreater{} ::= \textless{}atomique numérique\textgreater{} | \sphinxhyphen{} \textless{}expression \sphinxhyphen{}\textgreater{}}}

\item {} 
\sphinxAtStartPar
\sphinxcode{\sphinxupquote{\textless{}intervalles numériques\textgreater{} ::= \textless{}intervalle numérique\textgreater{} , …}}

\item {} 
\sphinxAtStartPar
\sphinxcode{\sphinxupquote{\textless{}intervalle numérique\textgreater{} ::= \textless{}expression numérique\textgreater{} (.. \textless{}expression numérique\textgreater{})?}}

\end{itemize}

\sphinxAtStartPar
Les \sphinxcode{\sphinxupquote{formules}} ont la forme suivante :

\begin{sphinxVerbatim}[commandchars=\\\{\}]
\PYGZlt{}formule\PYGZgt{} ::=
    \PYGZlt{}symbole\PYGZgt{} ([ \PYGZlt{}expression numérique\PYGZgt{} ])? = \PYGZlt{}expression numérique\PYGZgt{}
\end{sphinxVerbatim}

\sphinxAtStartPar
Les \sphinxcode{\sphinxupquote{multi\sphinxhyphen{}formules}} ont la forme suivante :

\begin{sphinxVerbatim}[commandchars=\\\{\}]
\PYG{o}{\PYGZlt{}}\PYG{n}{multi}\PYG{o}{\PYGZhy{}}\PYG{n}{formule}\PYG{o}{\PYGZgt{}} \PYG{p}{:}\PYG{o}{:=} \PYG{n}{pour} \PYG{o}{\PYGZlt{}}\PYG{n}{indices}\PYG{o}{\PYGZgt{}} \PYG{p}{:} \PYG{o}{\PYGZlt{}}\PYG{n}{formule}\PYG{o}{\PYGZgt{}}
\end{sphinxVerbatim}


\subsubsection{Instructions}
\label{\detokenize{syntax:instructions}}
\sphinxAtStartPar
Les \sphinxstyleemphasis{instructions} ont la forme suivante :

\begin{sphinxVerbatim}[commandchars=\\\{\}]
\PYGZlt{}instruction\PYGZgt{} ::=
  \PYGZlt{}formule\PYGZgt{}
  | \PYGZlt{}multi\PYGZhy{}formule\PYGZgt{} 
  | si \PYGZlt{}expression\PYGZti{}booléenne\PYGZgt{}
    alors \PYGZlt{}instruction\PYGZgt{}+
    (sinon \PYGZlt{}instruction\PYGZgt{}+)?
    finsi 
  | calculer domaine \PYGZlt{}symbole\PYGZgt{}+;
  | calculer enchaineur \PYGZlt{}symbole\PYGZgt{};
  | calculer cible \PYGZlt{}symbole\PYGZgt{}
    (: avec \PYGZlt{}atome\PYGZti{}réel\PYGZgt{} , …)?;
  | verifier domaine \PYGZlt{}symbole\PYGZgt{}+
    (: avec \PYGZlt{}expression\PYGZti{}booléenne\PYGZgt{});
  | (afficher | afficher\PYGZus{}\PYGZus{}erreur) \PYGZlt{}affichable\PYGZgt{}+;
  | iterer 
    : variable \PYGZlt{}symbole\PYGZgt{}
    : categorie \PYGZlt{}symbole\PYGZgt{}+ , …
    (: avec \PYGZlt{}expression\PYGZti{}booléenne\PYGZgt{})? 
    : dans ( \PYGZlt{}instruction\PYGZgt{}+ ) 
  | restaurer 
    ( 
    : \PYGZlt{}symbole\PYGZgt{} , … 
     | : variable \PYGZlt{}symbole\PYGZgt{} : categorie \PYGZlt{}symbole\PYGZgt{}+ , … 
       (: avec \PYGZlt{}expression\PYGZti{}booléenne\PYGZgt{})? 
    )+ 
    : apres ( \PYGZlt{}instruction\PYGZgt{}+ ) 
  | leve\PYGZus{}\PYGZus{}erreur \PYGZlt{}symbole\PYGZgt{} \PYGZlt{}symbole\PYGZgt{}?;
  | (
      nettoie\PYGZus{}\PYGZus{}erreurs
      | exporte\PYGZus{}\PYGZus{}erreurs
      | finalise\PYGZus{}\PYGZus{}erreurs
    ) ;
  | aiguillage (nom)? (\PYGZlt{}symbole\PYGZgt{}) : (\PYGZlt{}multiple\PYGZus{}cas\PYGZgt{})
  | stop \PYGZlt{}type\PYGZus{}stop\PYGZgt{}?
\end{sphinxVerbatim}

\sphinxAtStartPar
avec :

\begin{sphinxVerbatim}[commandchars=\\\{\}]
\PYGZlt{}affichable\PYGZgt{} ::=
  \PYGZlt{}chaîne\PYGZgt{} 
  | `*(*` \PYGZlt{}expression\PYGZti{}numérique\PYGZgt{} `*)*`
    (`*:*` \PYGZlt{}atome\PYGZti{}naturel\PYGZgt{} (`*..*` \PYGZlt{}atome\PYGZti{}naturel\PYGZgt{})?)? 
  | (`*nom*` | `*alias*`) `*(*` \PYGZlt{}symbole\PYGZgt{} `*)*` 
  | `*indenter*` `*(*` \PYGZlt{}atome\PYGZti{}naturel\PYGZgt{} `*)*`

\PYGZlt{}multiple cas\PYGZgt{} ::= \PYGZlt{}cas\PYGZgt{}*

\PYGZlt{}cas\PYGZgt{} ::=
  cas \PYGZlt{}atome reel\PYGZgt{} : \PYGZlt{}instruction\PYGZgt{}+
  | cas indefini : \PYGZlt{}instruction\PYGZgt{}+
  | cas \PYGZlt{}symbole\PYGZgt{} : \PYGZlt{}instruction\PYGZgt{}+
  | par\PYGZus{}default : \PYGZlt{}instruction\PYGZgt{}+

\PYGZlt{}type\PYGZus{}stop\PYGZgt{} ::=
  application
  | cible
  | fonction
  | \PYGZlt{}symbole\PYGZgt{}
\end{sphinxVerbatim}

\sphinxAtStartPar
Les \sphinxstyleemphasis{instructions simples} ont la forme suivante :

\begin{sphinxVerbatim}[commandchars=\\\{\}]
\PYGZlt{}instruction simple\PYGZgt{} ::= 
  \PYGZlt{}formule\PYGZgt{} 
  | \PYGZlt{}multi\PYGZhy{}formule\PYGZgt{} 
  | `*si*` \PYGZlt{}expression\PYGZti{}booléenne\PYGZgt{}
    `*alors*` \PYGZlt{}instruction\PYGZti{}simple\PYGZgt{}+
    (`*sinon*` \PYGZlt{}instruction\PYGZti{}simple\PYGZgt{}+)?
    `*finsi*` 
  | (`*afficher*` | `*afficher\PYGZus{}\PYGZus{}erreur*`) \PYGZlt{}affichable\PYGZgt{}+ `*;*`
\end{sphinxVerbatim}


\subsubsection{Définition d’une règle}
\label{\detokenize{syntax:definition-d-une-regle}}
\sphinxAtStartPar
Les \sphinxstyleemphasis{règles} ont la forme suivante :

\begin{sphinxVerbatim}[commandchars=\\\{\}]
\PYGZlt{}regle\PYGZgt{} ::= 
 regle \PYGZlt{}symbole\PYGZgt{}.. \PYGZlt{}naturel\PYGZgt{} :
    ( 
      application : \PYGZlt{}symbole\PYGZgt{} , … ; 
      \PYGZca{}\PYGZca{} (enchaineur : \PYGZlt{}symbole\PYGZgt{} , … ;)? 
      \PYGZca{}\PYGZca{} (variable temporaire : \PYGZlt{}déclaration\PYGZgt{} , … ;)? 
    ) 
    \PYGZlt{}instruction\PYGZti{}simple\PYGZgt{}+
\end{sphinxVerbatim}

\sphinxAtStartPar
avec :
\begin{itemize}
\item {} 
\sphinxAtStartPar
\sphinxcode{\sphinxupquote{\textless{}déclaration\textgreater{} ::= \textless{}symbole\textgreater{} ({[} \textless{}atome\textasciitilde{}naturel\textgreater{} {]})?}}

\end{itemize}


\subsubsection{Définition d’une vérification}
\label{\detokenize{syntax:definition-d-une-verification}}
\sphinxAtStartPar
Les \sphinxstyleemphasis{vérifications} ont la forme suivante :

\begin{sphinxVerbatim}[commandchars=\\\{\}]
\PYGZlt{}vérification\PYGZgt{} ::= 
  verif \PYGZlt{}symbole\PYGZgt{} \PYGZlt{}naturel\PYGZgt{} : 
  application : \PYGZlt{}symbole\PYGZgt{} , … ; 
  si \PYGZlt{}expression\PYGZti{}booléenne\PYGZgt{}
  alors erreur \PYGZlt{}symbole\PYGZgt{} \PYGZlt{}symbole\PYGZgt{}?
  ;
\end{sphinxVerbatim}


\subsubsection{Définition d’une fonction}
\label{\detokenize{syntax:definition-d-une-fonction}}
\sphinxAtStartPar
Les \sphinxstyleemphasis{fonctions} ont la forme suivante :

\begin{sphinxVerbatim}[commandchars=\\\{\}]
\PYGZlt{}fonction\PYGZgt{} ::= 
   fonction \PYGZlt{}symbole\PYGZgt{} : 
    ( 
      application : \PYGZlt{}symbole\PYGZgt{} , … ; 
      \PYGZca{}\PYGZca{} (variable temporaire : \PYGZlt{}déclaration\PYGZgt{} , … ;)? 
      \PYGZca{}\PYGZca{} (argument : \PYGZlt{}symbole\PYGZgt{} , … ;)? 
      \PYGZca{}\PYGZca{} resultat : \PYGZlt{}symbole\PYGZgt{} ; 
    ) 
    \PYGZlt{}instruction\PYGZti{}simple\PYGZgt{}+
\end{sphinxVerbatim}

\sphinxAtStartPar
avec :
\begin{itemize}
\item {} 
\sphinxAtStartPar
\sphinxcode{\sphinxupquote{\textless{}déclaration\textgreater{} ::= \textless{}symbole\textgreater{} ({[} \textless{}atome\textasciitilde{}naturel\textgreater{} {]})?}}

\end{itemize}


\subsubsection{Définition d’une cible}
\label{\detokenize{syntax:definition-d-une-cible}}
\sphinxAtStartPar
Les \sphinxstyleemphasis{cibles} ont la forme suivante :

\begin{sphinxVerbatim}[commandchars=\\\{\}]
\PYGZlt{}cible\PYGZgt{} ::= 
  cible \PYGZlt{}symbole\PYGZgt{} : 
  ( 
    application : \PYGZlt{}symbole\PYGZgt{} , … ; 
    \PYGZca{}\PYGZca{} (enchaineur : \PYGZlt{}symbole\PYGZgt{} , … ;)? 
    \PYGZca{}\PYGZca{} (variable temporaire : \PYGZlt{}déclaration\PYGZgt{} , … ;)? 
    \PYGZca{}\PYGZca{} (argument : \PYGZlt{}symbole\PYGZgt{} , … ;)?    
  ) 
  \PYGZlt{}instruction\PYGZgt{}+
\end{sphinxVerbatim}

\sphinxAtStartPar
avec :
\begin{itemize}
\item {} 
\sphinxAtStartPar
\sphinxcode{\sphinxupquote{\textless{}déclaration\textgreater{} ::= \textless{}symbole\textgreater{} ({[} \textless{}atome\textasciitilde{}naturel\textgreater{} {]})?}}

\end{itemize}


\subsubsection{Commentaires}
\label{\detokenize{syntax:commentaires}}
\sphinxAtStartPar
Les commentaires sont précédés du caractère \sphinxcode{\sphinxupquote{\#}}.
Il est possible d’inclure des commentaires multi\sphinxhyphen{}ligne avec les délimiteurs \sphinxcode{\sphinxupquote{\#\{}} et
\sphinxcode{\sphinxupquote{\}\#}}.

\sphinxAtStartPar
Exemple:

\begin{sphinxVerbatim}[commandchars=\\\{\}]
\PYG{c+c1}{\PYGZsh{} Ceci est un commentaire sur une ligne}
\PYG{c+c1}{\PYGZsh{}\PYGZob{} Ceci est un commentaire }
   \PYG{n}{sur} \PYG{n}{plusieurs} \PYG{n}{lignes}\PYG{o}{.} \PYG{p}{\PYGZcb{}}\PYG{c+c1}{\PYGZsh{}}
\end{sphinxVerbatim}


\subsection{Les valeurs}
\label{\detokenize{syntax:les-valeurs}}\label{\detokenize{syntax:valeurs}}
\sphinxAtStartPar
Les variables en M prennent soit leur valeur sur les flottants, soit ont la
valeur \sphinxcode{\sphinxupquote{indefini}}.
Les valeurs de type booléen sont représentées par les flottants \sphinxcode{\sphinxupquote{0}} et
\sphinxcode{\sphinxupquote{1}}, respectivement pour \sphinxcode{\sphinxupquote{faux}} et \sphinxcode{\sphinxupquote{vrai}}.

\sphinxAtStartPar
Les résultats suivants peuvent être observés en exécutant le script disponible
dans l’exemple sur le {\hyperref[\detokenize{exemples/valeurs:exemples-valeurs}]{\sphinxcrossref{\DUrole{std}{\DUrole{std-ref}{Calcul des valeurs}}}}}.


\subsubsection{Calcul booléen}
\label{\detokenize{syntax:calcul-booleen}}
\sphinxAtStartPar
Les opérations booléennes standard (\sphinxcode{\sphinxupquote{et}}, \sphinxcode{\sphinxupquote{ou}}, \sphinxcode{\sphinxupquote{non}}) comportent sur les
flottants de façon standard.
Toute valeur flottante différente de \sphinxcode{\sphinxupquote{0}} sera considerée comme \sphinxcode{\sphinxupquote{vrai}}e dans
le cas d’un calcul booléen (\sphinxcode{\sphinxupquote{10 ou 0 = 1}}).
Les calculs booléens impliquant la valeur \sphinxcode{\sphinxupquote{indefini}} ont un comportement
spécifique :
\begin{itemize}
\item {} 
\sphinxAtStartPar
\sphinxcode{\sphinxupquote{indefini ou indefini = indefini}}

\item {} 
\sphinxAtStartPar
\sphinxcode{\sphinxupquote{indefini ou b = (0 si b = 0, 1 sinon)}}

\item {} 
\sphinxAtStartPar
\sphinxcode{\sphinxupquote{b ou indefini = (0 si b = 0, 1 sinon)}}

\item {} 
\sphinxAtStartPar
\sphinxcode{\sphinxupquote{indefini et indefini = indefini}}

\item {} 
\sphinxAtStartPar
\sphinxcode{\sphinxupquote{b et indefini = indefini}}

\item {} 
\sphinxAtStartPar
\sphinxcode{\sphinxupquote{indefini et b = indefini}}

\item {} 
\sphinxAtStartPar
\sphinxcode{\sphinxupquote{non indefini = indefini}}

\end{itemize}

\sphinxAtStartPar
Les comparaisons (\sphinxcode{\sphinxupquote{=}}, \sphinxcode{\sphinxupquote{\textless{}}}, \sphinxcode{\sphinxupquote{\textgreater{}}}, \sphinxcode{\sphinxupquote{\textless{}=}}, \sphinxcode{\sphinxupquote{\textgreater{}=}}) renvoient soit \sphinxcode{\sphinxupquote{0}} soit \sphinxcode{\sphinxupquote{1}}
lorsque des valeurs definies sont comparées.
Si l’une des valeurs comparée est \sphinxcode{\sphinxupquote{indéfini}}e, alors le résultat est également
\sphinxcode{\sphinxupquote{indefini}}.

\sphinxAtStartPar
\sphinxstylestrong{Note}: même l’égalité \sphinxcode{\sphinxupquote{indefini = indefini}} renvoie \sphinxcode{\sphinxupquote{indefini}}.
Pour vérifier si une valeur est définie ou non, il faut utiliser la fonction
\sphinxcode{\sphinxupquote{present}} qui renvoie \sphinxcode{\sphinxupquote{0}} si la valeur est indéfinie et \sphinxcode{\sphinxupquote{1}} sinon.


\subsubsection{Calcul numérique}
\label{\detokenize{syntax:calcul-numerique}}
\sphinxAtStartPar
Les opérations arithmétiques standard (\sphinxcode{\sphinxupquote{+}}, \sphinxcode{\sphinxupquote{\sphinxhyphen{}}}, \sphinxcode{\sphinxupquote{*}}, \sphinxcode{\sphinxupquote{/}}) se comportent sur
les flottants de façon standard, à l’exception de la division d’un flottant par
zero qui renvoie toujours \sphinxcode{\sphinxupquote{0}}.
Les calculs impliquant la valeur \sphinxcode{\sphinxupquote{indefini}} ont un comportement spécifique :
\begin{itemize}
\item {} 
\sphinxAtStartPar
\sphinxcode{\sphinxupquote{indefini + indefini = indefini}}

\item {} 
\sphinxAtStartPar
\sphinxcode{\sphinxupquote{indefini + n = n}}

\item {} 
\sphinxAtStartPar
\sphinxcode{\sphinxupquote{n + indefini = n}}

\item {} 
\sphinxAtStartPar
\sphinxcode{\sphinxupquote{indefini \sphinxhyphen{} indefini = indefini}}

\item {} 
\sphinxAtStartPar
\sphinxcode{\sphinxupquote{indefini \sphinxhyphen{} n = \sphinxhyphen{}n}}

\item {} 
\sphinxAtStartPar
\sphinxcode{\sphinxupquote{n \sphinxhyphen{} indefini = n}}

\item {} 
\sphinxAtStartPar
\sphinxcode{\sphinxupquote{indefini * indefini = indefini}}

\item {} 
\sphinxAtStartPar
\sphinxcode{\sphinxupquote{indefini * n = indefini}}

\item {} 
\sphinxAtStartPar
\sphinxcode{\sphinxupquote{n * indefini = indefini}}

\item {} 
\sphinxAtStartPar
\sphinxcode{\sphinxupquote{indefini / indefini = indefini}}

\item {} 
\sphinxAtStartPar
\sphinxcode{\sphinxupquote{indefini / n = indefini}}

\item {} 
\sphinxAtStartPar
\sphinxcode{\sphinxupquote{n / indefini = indefini}}

\end{itemize}

\sphinxAtStartPar
Pour résumer, seules les opérations d’addition et de soustraction avec une
valeur renvoient autre chose que \sphinxcode{\sphinxupquote{indefini}}.

\sphinxstepscope


\section{Les fonctions}
\label{\detokenize{fonctions:les-fonctions}}\label{\detokenize{fonctions:fonctions}}\label{\detokenize{fonctions::doc}}
\sphinxAtStartPar
Voici la liste de toutes les fonctions standard du M.
Leur utilisation est présentée dans l’exemple sur
{\hyperref[\detokenize{exemples/fonctions:exemples-fonctions}]{\sphinxcrossref{\DUrole{std}{\DUrole{std-ref}{Les fonctions}}}}}.


\subsection{abs(X)}
\label{\detokenize{fonctions:abs-x}}
\sphinxAtStartPar
Cette fonction prend un argument et renvoie sa valeur absolue.
Si \sphinxcode{\sphinxupquote{X}} est \sphinxcode{\sphinxupquote{indefini}}, alors \sphinxcode{\sphinxupquote{abs(X) = indefini}}.


\subsection{afficher(X) / afficher « texte »}
\label{\detokenize{fonctions:afficher-x-afficher-texte}}
\sphinxAtStartPar
Cette fonction affiche sur la sortie standard la valeur de l’expression en argument.


\subsection{arr(X)}
\label{\detokenize{fonctions:arr-x}}
\sphinxAtStartPar
Cette fonction prend un argument et renvoie l’entier le plus
proche, ou \sphinxcode{\sphinxupquote{indefini}} s’il est \sphinxcode{\sphinxupquote{indefini}}.


\subsection{attribut(X, attribut)}
\label{\detokenize{fonctions:attribut-x-attribut}}
\sphinxAtStartPar
Cette fonction prend un argument \sphinxcode{\sphinxupquote{X}} et un nom d’attribut \sphinxcode{\sphinxupquote{a}} et retourne
la valeur associée à l’attribut \sphinxcode{\sphinxupquote{a}} de \sphinxcode{\sphinxupquote{X}}.


\subsection{champ\_evenement(X, champ)}
\label{\detokenize{fonctions:champ-evenement-x-champ}}
\sphinxAtStartPar
Cette fonction prend un argument \sphinxcode{\sphinxupquote{X}} et un nom de champ d’événement \sphinxcode{\sphinxupquote{champ}}.
La valeur \sphinxcode{\sphinxupquote{X}} correspond à l’identifiant de l’événement. \%TODO: référence chap événement
Si le \sphinxcode{\sphinxupquote{champ}} correspond à une valeur, \sphinxcode{\sphinxupquote{champ\_evenement}} la retourne. Si
le \sphinxcode{\sphinxupquote{champ}} correspond à une variable, \sphinxcode{\sphinxupquote{champ\_evenement}} retourne sa valeur.

\sphinxAtStartPar
\sphinxstylestrong{Note} : \sphinxcode{\sphinxupquote{champ\_evenement}} peut être utilisée pour assigner des valeurs et
des références à des événéments. Exemple:

\begin{sphinxVerbatim}[commandchars=\\\{\}]
\PYG{n}{evenement}
\PYG{p}{:} \PYG{n}{valeur} \PYG{n}{ev\PYGZus{}val}
\PYG{p}{:} \PYG{n}{variable} \PYG{n}{ev\PYGZus{}var}\PYG{p}{;}

\PYG{c+c1}{\PYGZsh{} Associe la valeur 42 au champ ev\PYGZus{}val de l\PYGZsq{}événement 0.}
\PYG{n}{champ\PYGZus{}evenement} \PYG{p}{(}\PYG{l+m+mi}{0}\PYG{p}{,}\PYG{n}{ev\PYGZus{}val}\PYG{p}{)} \PYG{o}{=} \PYG{l+m+mi}{42}\PYG{p}{;}
\PYG{c+c1}{\PYGZsh{} Associe la variable X au champ ev\PYGZus{}var de l\PYGZsq{}événement 0.}
\PYG{n}{champ\PYGZus{}evenement} \PYG{p}{(}\PYG{l+m+mi}{0}\PYG{p}{,}\PYG{n}{ev\PYGZus{}var}\PYG{p}{)} \PYG{n}{reference} \PYG{n}{X}\PYG{p}{;}
\end{sphinxVerbatim}


\subsection{inf(X)}
\label{\detokenize{fonctions:inf-x}}
\sphinxAtStartPar
Cette fonction prend un argument et renvoie l’entier inferieur
le plus proche, ou \sphinxcode{\sphinxupquote{indefini}} s’il est \sphinxcode{\sphinxupquote{indefini}}.


\subsection{max(X, Y)}
\label{\detokenize{fonctions:max-x-y}}
\sphinxAtStartPar
Cette fonction prend deux arguments et renvoie la plus grande
valeur des deux.
Si les deux valeurs sont \sphinxcode{\sphinxupquote{indefini}}es, alors la fonction renvoie
\sphinxcode{\sphinxupquote{indefini}}.
Si \sphinxcode{\sphinxupquote{X}} est défini et \sphinxcode{\sphinxupquote{Y}} est \sphinxcode{\sphinxupquote{indefini}}, alors \sphinxcode{\sphinxupquote{max(X,Y) = max(X, 0)}}. De même,
si \sphinxcode{\sphinxupquote{X}} est \sphinxcode{\sphinxupquote{indéfini}} et \sphinxcode{\sphinxupquote{Y}} est defini, alors \sphinxcode{\sphinxupquote{max(X,Y) = max(0, Y)}}.


\subsection{meme\_variable(X, Y)}
\label{\detokenize{fonctions:meme-variable-x-y}}
\sphinxAtStartPar
Cette fonction prend en argument deux variables et vérifie s’il s’agit de la même
variable.
Elle est notamment utilisée dans les iterateurs de variables pour effectuer des
traitements spécifiques.


\subsection{multimax(X, TAB)}
\label{\detokenize{fonctions:multimax-x-tab}}
\sphinxAtStartPar
Cette fonction prend deux arguments \sphinxcode{\sphinxupquote{X}} une valeur et \sphinxcode{\sphinxupquote{TAB}} un tableau
de valeurs.
Elle calcule la plus grande valeur contenue dans \sphinxcode{\sphinxupquote{TAB}} entre \sphinxcode{\sphinxupquote{0}} et \sphinxcode{\sphinxupquote{X \sphinxhyphen{} 1}}.
Si X est \sphinxcode{\sphinxupquote{indefini}} ou \sphinxcode{\sphinxupquote{X \textless{}= 0}}, alors \sphinxcode{\sphinxupquote{multimax(X, TAB) = indefini}}.
\sphinxcode{\sphinxupquote{X}} peut être plus grand que la taille de \sphinxcode{\sphinxupquote{TAB}}, auquel cas \sphinxcode{\sphinxupquote{multimax}} calculera le
maximum de \sphinxcode{\sphinxupquote{TAB}}.


\subsection{min(X, Y)}
\label{\detokenize{fonctions:min-x-y}}
\sphinxAtStartPar
Cette fonction prend deux arguments et renvoie la plus petite
valeur des deux.
Si les deux valeurs sont \sphinxcode{\sphinxupquote{indefini}}es, alors la fonction renvoie
\sphinxcode{\sphinxupquote{indefini}}.
Si \sphinxcode{\sphinxupquote{X}} est défini et \sphinxcode{\sphinxupquote{Y}} est \sphinxcode{\sphinxupquote{indefini}}, alors \sphinxcode{\sphinxupquote{min(X,Y) = min(X, 0)}}. De même,
si \sphinxcode{\sphinxupquote{X}} est \sphinxcode{\sphinxupquote{indéfini}} et \sphinxcode{\sphinxupquote{Y}} est defini, alors \sphinxcode{\sphinxupquote{min(X,Y) = min(0, Y)}}.


\subsection{nb\_evenements()}
\label{\detokenize{fonctions:nb-evenements}}
\sphinxAtStartPar
Cette fonction renvoie le nombre total d’événements.


\subsection{null(X)}
\label{\detokenize{fonctions:null-x}}
\sphinxAtStartPar
Cette fonction prend un argument \sphinxcode{\sphinxupquote{X}} et renvoie \sphinxcode{\sphinxupquote{1}} si \sphinxcode{\sphinxupquote{X}} est
égal à \sphinxcode{\sphinxupquote{0}}, \sphinxcode{\sphinxupquote{0}} si \sphinxcode{\sphinxupquote{X}} est défini et différent de \sphinxcode{\sphinxupquote{0}},
et \sphinxcode{\sphinxupquote{indefini}} s’il est \sphinxcode{\sphinxupquote{indefini}}.

\sphinxAtStartPar
Cette fonction est strictement équivalente à l’expression \sphinxcode{\sphinxupquote{(X = 0)}}


\subsection{positif(X)}
\label{\detokenize{fonctions:positif-x}}
\sphinxAtStartPar
Cette fonction prend un argument \sphinxcode{\sphinxupquote{X}} et renvoie \sphinxcode{\sphinxupquote{1}} si \sphinxcode{\sphinxupquote{X}} est
\sphinxstylestrong{strictement} supérieur à \sphinxcode{\sphinxupquote{0}}, \sphinxcode{\sphinxupquote{0}} si \sphinxcode{\sphinxupquote{X}} est inferieur ou égal à \sphinxcode{\sphinxupquote{0}}, et
\sphinxcode{\sphinxupquote{indefini}} si \sphinxcode{\sphinxupquote{X}} est indefini.

\sphinxAtStartPar
Cette fonction est strictement équivalente à l’expression \sphinxcode{\sphinxupquote{X \textgreater{} 0}}.


\subsection{positif\_ou\_nul(X)}
\label{\detokenize{fonctions:positif-ou-nul-x}}
\sphinxAtStartPar
Cette fonction prend un argument \sphinxcode{\sphinxupquote{X}} et renvoie \sphinxcode{\sphinxupquote{1}} si \sphinxcode{\sphinxupquote{X}} est
supérieur ou égale à \sphinxcode{\sphinxupquote{0}}, \sphinxcode{\sphinxupquote{0}} si \sphinxcode{\sphinxupquote{X}} est strictement inferieur ou égal à
\sphinxcode{\sphinxupquote{0}}, et \sphinxcode{\sphinxupquote{indefini}} si \sphinxcode{\sphinxupquote{X}} est indefini.

\sphinxAtStartPar
Cette fonction est strictement équivalente à l’expression \sphinxcode{\sphinxupquote{X \textgreater{}= 0}}.


\subsection{present(X)}
\label{\detokenize{fonctions:present-x}}
\sphinxAtStartPar
Cette fonction prend un argument \sphinxcode{\sphinxupquote{X}} et renvoie \sphinxcode{\sphinxupquote{0}} s’il
est \sphinxcode{\sphinxupquote{indefini}} et \sphinxcode{\sphinxupquote{1}} sinon.

\sphinxAtStartPar
\sphinxstylestrong{Note} : cette fonction est la seule façon de tester si une valeur est
égale à \sphinxcode{\sphinxupquote{indefini}} car l’expression booléenne \sphinxcode{\sphinxupquote{(indefini = indefini)}}
est égale à \sphinxcode{\sphinxupquote{indefini}}.


\subsection{somme(X, Y, …)}
\label{\detokenize{fonctions:somme-x-y}}
\sphinxAtStartPar
Cette fonction n’a pas de limite d’arguments.
Elle calcule leur somme.
Si aucun argument n’est donné à la fonction \sphinxcode{\sphinxupquote{somme}}, elle renvoie \sphinxcode{\sphinxupquote{0}}.
Si tous ses arguments sont \sphinxcode{\sphinxupquote{indefini}}s, alors elle renvoie la valeur
\sphinxcode{\sphinxupquote{indefini}}.


\subsection{supzero(X)}
\label{\detokenize{fonctions:supzero-x}}
\sphinxAtStartPar
Cette fonction prend un argument \sphinxcode{\sphinxupquote{X}} et renvoie sa valeur s’il est strictement supérieur
à \sphinxcode{\sphinxupquote{0}}, sinon il renvoie \sphinxcode{\sphinxupquote{indefini}}.


\subsection{taille(TAB)}
\label{\detokenize{fonctions:taille-tab}}
\sphinxAtStartPar
Cette fonction prend en argument soit une variable, soit un tableau.
S’il s’agit d’une variable, \sphinxcode{\sphinxupquote{taille}} renvoie \sphinxcode{\sphinxupquote{1}}, sinon elle renvoie la taille du
tableau.
Elle échoue si l’argument est une constante ou une valeur.


\subsection{type(V, type)}
\label{\detokenize{fonctions:type-v-type}}
\sphinxAtStartPar
Cette fonction prend en argument une variable et un type. Elle
renvoie 1 si V est du type donné, 0 sinon.

\sphinxstepscope


\section{Les fichiers IRJ}
\label{\detokenize{syntax_irj:les-fichiers-irj}}\label{\detokenize{syntax_irj:syntax-irj}}\label{\detokenize{syntax_irj::doc}}

\subsection{Syntaxe}
\label{\detokenize{syntax_irj:syntaxe}}
\sphinxAtStartPar
Les fichiers IRJ sont des fichiers de test faisant corespondre entrées et
sorties du calcul d’un programme M.
Chaque fichier IRJ est divisé en 7 parties, plus 3 parties optionnelles,
et est terminé par la ligne \sphinxcode{\sphinxupquote{\#\#}}.
\begin{itemize}
\item {} 
\sphinxAtStartPar
\sphinxcode{\sphinxupquote{NOM}} : le nom du test.

\item {} 
\sphinxAtStartPar
\sphinxcode{\sphinxupquote{ENTREES\sphinxhyphen{}PRIMITIF}} : les valeurs des variables d’entrées au début du calcul. A chaque variable \sphinxcode{\sphinxupquote{VAR}} est associée une valeur \sphinxcode{\sphinxupquote{val}} entière ou décimale sur une ligne de la forme \sphinxcode{\sphinxupquote{VAR/val}}. Les variables d’entrées qui ne sont pas définies dans le fichier IRJ seront initialisées à \sphinxcode{\sphinxupquote{indefini}}.

\item {} 
\sphinxAtStartPar
\sphinxcode{\sphinxupquote{CONTROLES\sphinxhyphen{}PRIMITIF}} : à documenter

\item {} 
\sphinxAtStartPar
\sphinxcode{\sphinxupquote{RESULTATS\sphinxhyphen{}PRIMITIF}} : les valeurs des variables restituées à la fin du calcul. Comme pour les variables d’entrées, elles sont décrites sous la forme \sphinxcode{\sphinxupquote{VAR/val}}. Les valeurs restituées absentes du fichiers IRJ sont supposées avoir la valeur \sphinxcode{\sphinxupquote{indefini}} à la fin du calcul.

\item {} 
\sphinxAtStartPar
\sphinxcode{\sphinxupquote{ENTREES\sphinxhyphen{}CORRECTIF}} : à documenter

\item {} 
\sphinxAtStartPar
\sphinxcode{\sphinxupquote{CONTROLES\sphinxhyphen{}CORRECTIF}} : à documenter

\item {} 
\sphinxAtStartPar
\sphinxcode{\sphinxupquote{RESULTATS\sphinxhyphen{}CORRECTIF}} : à documenter

\item {} 
\sphinxAtStartPar
\sphinxcode{\sphinxupquote{ENTREES\sphinxhyphen{}RAPPELS}} : partie optionnelle, à documenter

\item {} 
\sphinxAtStartPar
\sphinxcode{\sphinxupquote{CONTROLES\sphinxhyphen{}RAPPELS}} : partie optionnelle, à documenter

\item {} 
\sphinxAtStartPar
\sphinxcode{\sphinxupquote{RESULTATS\sphinxhyphen{}RAPPELS}} : partie optionnelle, à documenter

\end{itemize}

\sphinxAtStartPar
Voici un exemple de fichier IRJ minimal

\begin{sphinxVerbatim}[commandchars=\\\{\}]
   \PYG{c+c1}{\PYGZsh{}NOM}
   \PYG{n}{MON}\PYG{o}{\PYGZhy{}}\PYG{n}{TEST}
   \PYG{c+c1}{\PYGZsh{}ENTREES\PYGZhy{}PRIMITIF}
   \PYG{n}{X}\PYG{o}{/}\PYG{l+m+mf}{0.0}
   \PYG{c+c1}{\PYGZsh{}CONTROLES\PYGZhy{}PRIMITIF}
   \PYG{c+c1}{\PYGZsh{}RESULTATS\PYGZhy{}PRIMITIF}
   \PYG{n}{Y}\PYG{o}{/}\PYG{l+m+mi}{1}
   \PYG{n}{Z}\PYG{o}{/}\PYG{l+m+mf}{2.0}
   \PYG{c+c1}{\PYGZsh{}ENTREES\PYGZhy{}CORRECTIF}
   \PYG{c+c1}{\PYGZsh{}CONTROLES\PYGZhy{}CORRECTIF}
   \PYG{c+c1}{\PYGZsh{}RESULTATS\PYGZhy{}CORRECTIF}
   \PYG{c+c1}{\PYGZsh{}\PYGZsh{}}
\end{sphinxVerbatim}

\sphinxAtStartPar
De nombreux exemples de fichiers IRJ sont disponibles dans le dossier \sphinxcode{\sphinxupquote{tests/}}.


\subsection{Utilisation}
\label{\detokenize{syntax_irj:utilisation}}
\sphinxAtStartPar
Trois utilitaires sont dédiés à l’utilisation des fichiers IRJ.


\subsubsection{IRJ checker}
\label{\detokenize{syntax_irj:irj-checker}}
\sphinxAtStartPar
Les sources de mlang mettent à disposition un code de vérification de fichiers
IRJ.
Après compilation, il peut être exécuté via la commande :

\begin{sphinxVerbatim}[commandchars=\\\{\}]
\PYGZdl{} dune exec \PYGZhy{}\PYGZhy{} irj\PYGZus{}checker
\end{sphinxVerbatim}


\subsubsection{Interpreteur}
\label{\detokenize{syntax_irj:interpreteur}}
\sphinxAtStartPar
Le binaire mlang intègre un interpreteur de M qui prend en entrée un ou
plusieurs fichiers M ainsi qu’un fichier IRJ.
Après compilation, l’interpreteur peut être exécuté via la commande :

\begin{sphinxVerbatim}[commandchars=\\\{\}]
\PYGZdl{} dune exec \PYGZhy{}\PYGZhy{} mlang test.m \PYGZhy{}A application \PYGZhy{}b interpreter \PYGZhy{}\PYGZhy{}mpp\PYGZus{}function cible\PYGZus{}dentree  \PYGZhy{}\PYGZhy{}dgfip\PYGZus{}options=\PYGZsq{}\PYGZsq{} \PYGZhy{}r test.irj
\end{sphinxVerbatim}

\sphinxAtStartPar
L’option \sphinxcode{\sphinxupquote{\sphinxhyphen{}r}} peut être remplacée par \sphinxcode{\sphinxupquote{\sphinxhyphen{}R}} pour tester l’ensemble des tests d’un
dossier complet.


\subsubsection{Fichiers C}
\label{\detokenize{syntax_irj:fichiers-c}}
\sphinxAtStartPar
Il est également possible d’utiliser les fichiers IRJ comme entrée du code C
généré à partir d’une compilation de code M par \sphinxcode{\sphinxupquote{mlang}}.
Ce traitement est effectué dans le cas du backend \sphinxcode{\sphinxupquote{dgfip\_c}} de mlang, dont
le code est disponible dans \sphinxcode{\sphinxupquote{examples/dgfip\_c/ml\_primitif}}.
La calculette d’une année donnée peut êter compilée via :

\begin{sphinxVerbatim}[commandchars=\\\{\}]
\PYGZdl{} make YEAR=year compile\PYGZus{}dgfip\PYGZus{}c\PYGZus{}backend
\end{sphinxVerbatim}

\sphinxAtStartPar
Les scripts présents dans \sphinxcode{\sphinxupquote{examples/dgfip\_c/ml\_primitif/ml\_driver}} permettent
d’interfacer les fichiers IRJ avec les valeurs \sphinxcode{\sphinxupquote{C}} de type \sphinxcode{\sphinxupquote{T\_irdata}} contenant
le TGV (Tableau Général des Variables).
La commande suivante permet de tester la calculette d’une année sur l’ensemble
des tests fuzzés de la dite année.

\begin{sphinxVerbatim}[commandchars=\\\{\}]
\PYGZdl{} make test\PYGZus{}dgfip\PYGZus{}c\PYGZus{}backend
\end{sphinxVerbatim}

\sphinxAtStartPar
L’utilisation du script \sphinxcode{\sphinxupquote{cal}} utilisé dans cette règle est documentée dans
\sphinxcode{\sphinxupquote{examples/dgfip\_c/README.md}}.

\sphinxstepscope


\section{Le compilateur MLang}
\label{\detokenize{mlang:le-compilateur-mlang}}\label{\detokenize{mlang:mlang}}\label{\detokenize{mlang::doc}}

\subsection{Installer}
\label{\detokenize{mlang:installer}}
\sphinxAtStartPar
Mlang est implanté en OCaml. L’utilisation du gestionnaire de paquets OCaml \sphinxcode{\sphinxupquote{opam}}
est fortement recommandée.
Vous pouvez l’installer via votre gestionnaire de paquet préféré s’il distribue
\sphinxcode{\sphinxupquote{opam}}, ou bien en vous referant à la \sphinxhref{https://opam.ocaml.org/doc/Install.html}{documentation d’opam}.
Mlang a également quelques autres dépendances, dont une vers la librairie de calcul
de flotants MPFR. Si vous êtes sous Debian, vous pouvez simplement utiliser la commande
suivante pour installer toutes les dépendances externes à OCaml :

\begin{sphinxVerbatim}[commandchars=\\\{\}]
\PYGZdl{} sudo apt install \PYGZbs{}
	libgmp\PYGZhy{}dev \PYGZbs{}
	libmpfr\PYGZhy{}dev \PYGZbs{}
	git \PYGZbs{}
	patch \PYGZbs{}
	unzip \PYGZbs{}
	bubblewrap \PYGZbs{}
	bzip2 \PYGZbs{}
	opam
\end{sphinxVerbatim}

\sphinxAtStartPar
Si vous n’avez jamais utilisé \sphinxcode{\sphinxupquote{opam}}, commencez par lancer :

\begin{sphinxVerbatim}[commandchars=\\\{\}]
\PYGZdl{} opam init
\PYGZdl{} opam update
\end{sphinxVerbatim}

\sphinxAtStartPar
Enfin, vous pouvre initialiser le projet \sphinxcode{\sphinxupquote{mlang}} avec

\begin{sphinxVerbatim}[commandchars=\\\{\}]
\PYGZdl{} make init
\end{sphinxVerbatim}

\sphinxAtStartPar
\sphinxstylestrong{Note pour les utilisateurs d’opam confirmés} : la commande \sphinxcode{\sphinxupquote{make init}} crée
un switch local où seront installées les dépendances OCaml de mlang.

\sphinxAtStartPar
Cette commande initialise le dossier \sphinxcode{\sphinxupquote{ir\sphinxhyphen{}calcul}} dans lequel sont poussés
le code de calcul primitif de l’impot sur le revenu.

\sphinxAtStartPar
Si besoin, la commande

\begin{sphinxVerbatim}[commandchars=\\\{\}]
\PYGZdl{} make deps
\end{sphinxVerbatim}

\sphinxAtStartPar
réinstallera les dépendances OCaml et mettra à jour le dossier \sphinxcode{\sphinxupquote{ir\sphinxhyphen{}calcul}}.

\sphinxAtStartPar
Une fois compilé, vous pouvez soit appeler mlang via la commande :

\begin{sphinxVerbatim}[commandchars=\\\{\}]
\PYGZdl{} opam exec \PYGZhy{}\PYGZhy{} mlang
\end{sphinxVerbatim}

\sphinxAtStartPar
tant que vous êtes dans le dossier depuis lequel vous avez compilé, soit
l’installer localement avec la commande :

\begin{sphinxVerbatim}[commandchars=\\\{\}]
\PYG{n}{opam} \PYG{n}{install} \PYG{o}{.}\PYG{o}{/}\PYG{n}{mlang}\PYG{o}{.}\PYG{n}{opam}
\end{sphinxVerbatim}


\subsection{Utiliser MLang}
\label{\detokenize{mlang:utiliser-mlang}}
\sphinxAtStartPar
Le binaire \sphinxcode{\sphinxupquote{mlang}} prend en argument le fichier \sphinxstyleemphasis{M} à exécuter.
Il peut également prendre en argument une liste de fichiers, auquel cas leur
traitement sera équivalent au traitement d’un seul et même fichier dans lequel
serait concatené le contenu de chaque fichier.

\sphinxAtStartPar
Les options principales sont :
\begin{itemize}
\item {} 
\sphinxAtStartPar
\sphinxcode{\sphinxupquote{\sphinxhyphen{}A}}: le nom de l’application à traiter;

\item {} 
\sphinxAtStartPar
\sphinxcode{\sphinxupquote{\sphinxhyphen{}b}}: le mode d’utilisation, ou \sphinxcode{\sphinxupquote{backend}};

\item {} 
\sphinxAtStartPar
\sphinxcode{\sphinxupquote{\sphinxhyphen{}\sphinxhyphen{}mpp\_function}}: le nom de la fonction principale à traiter.

\end{itemize}


\subsubsection{Mode interpreteur}
\label{\detokenize{mlang:mode-interpreteur}}
\sphinxAtStartPar
Le mode interpreteur de \sphinxcode{\sphinxupquote{mlang}} utilise un fichier \sphinxstyleemphasis{IRJ} (voir {\hyperref[\detokenize{syntax_irj:syntax-irj}]{\sphinxcrossref{\DUrole{std}{\DUrole{std-ref}{Les fichiers IRJ}}}}})
pour exécuter le code \sphinxstyleemphasis{M} directement depuis sa représentation abstraite.

\sphinxAtStartPar
Voici une commande simple pour invoquer l’interpreteur :

\begin{sphinxVerbatim}[commandchars=\\\{\}]
\PYGZdl{} mlang test.m \PYGZbs{}
	\PYGZhy{}A mon\PYGZus{}application \PYGZbs{}
	\PYGZhy{}b interpreter \PYGZbs{}
	\PYGZhy{}\PYGZhy{}mpp\PYGZus{}function hello\PYGZus{}world \PYGZbs{}
	\PYGZhy{}\PYGZhy{}run\PYGZus{}test test.irj \PYGZbs{}
	\PYGZhy{}\PYGZhy{}without\PYGZus{}dfgip\PYGZus{}m
\end{sphinxVerbatim}


\subsubsection{Mode transpilation}
\label{\detokenize{mlang:mode-transpilation}}
\sphinxAtStartPar
Le mode transpilation de \sphinxcode{\sphinxupquote{mlang}} permet de traduire le code \sphinxstyleemphasis{M} dans un autre
langage.
En 2025, seul le langage C est supporté.
Voici une commande simple pour traduire un fichier \sphinxstyleemphasis{M} en \sphinxstyleemphasis{C} :

\begin{sphinxVerbatim}[commandchars=\\\{\}]
\PYGZdl{} mlang test.m \PYGZbs{}
	\PYGZhy{}A mon\PYGZus{}application \PYGZbs{}
	\PYGZhy{}b dgfip\PYGZus{}c \PYGZbs{}
	\PYGZhy{}\PYGZhy{}mpp\PYGZus{}function hello\PYGZus{}world
	\PYGZhy{}\PYGZhy{}dgfip\PYGZus{}options=\PYGZsq{}\PYGZsq{} 
	\PYGZhy{}\PYGZhy{}output output/mon\PYGZhy{}test.c
\end{sphinxVerbatim}

\sphinxAtStartPar
NB: le dossier \sphinxcode{\sphinxupquote{output}} doit avoir été créé en amont.


\subsubsection{Options DGFiP}
\label{\detokenize{mlang:options-dgfip}}
\sphinxAtStartPar
Les options DGFiP sont à usage interne. Elles sont spécifiées dans
l’option \sphinxcode{\sphinxupquote{\sphinxhyphen{}\sphinxhyphen{}dgfip\_options}}.

\begin{sphinxVerbatim}[commandchars=\\\{\}]
       \PYG{o}{\PYGZhy{}}\PYG{n}{b} \PYG{n}{VAL}
           \PYG{n}{Set} \PYG{n}{application} \PYG{n}{to} \PYG{l+s+s2}{\PYGZdq{}}\PYG{l+s+s2}{batch}\PYG{l+s+s2}{\PYGZdq{}} \PYG{p}{(}\PYG{n}{b0} \PYG{o}{=} \PYG{n}{normal}\PYG{p}{,} \PYG{n}{b1} \PYG{o}{=} \PYG{k}{with} \PYG{n}{EBCDIC} \PYG{n}{sort}\PYG{p}{)}

       \PYG{o}{\PYGZhy{}}\PYG{n}{D}  \PYG{n}{Generate} \PYG{n}{labels} \PYG{k}{for} \PYG{n}{output} \PYG{n}{variables}

       \PYG{o}{\PYGZhy{}}\PYG{n}{g}  \PYG{n}{Generate} \PYG{k}{for} \PYG{n}{test} \PYG{p}{(}\PYG{n}{debug}\PYG{p}{)}

       \PYG{o}{\PYGZhy{}}\PYG{n}{I}  \PYG{n}{Generate} \PYG{n}{immediate} \PYG{n}{controls}

       \PYG{o}{\PYGZhy{}}\PYG{n}{k} \PYG{n}{VAL} \PYG{p}{(}\PYG{n}{absent}\PYG{o}{=}\PYG{l+m+mi}{0}\PYG{p}{)}
           \PYG{n}{Number} \PYG{n}{of} \PYG{n}{debug} \PYG{n}{files}

       \PYG{o}{\PYGZhy{}}\PYG{n}{L}  \PYG{n}{Generate} \PYG{n}{calls} \PYG{n}{to} \PYG{n}{ticket} \PYG{n}{function}

       \PYG{o}{\PYGZhy{}}\PYG{n}{m} \PYG{n}{VAL} \PYG{p}{(}\PYG{n}{absent}\PYG{o}{=}\PYG{l+m+mi}{1991}\PYG{p}{)}
           \PYG{n}{Income} \PYG{n}{year}

       \PYG{o}{\PYGZhy{}}\PYG{n}{O}  \PYG{n}{Optimize} \PYG{n}{generated} \PYG{n}{code} \PYG{p}{(}\PYG{n}{inline} \PYG{n}{min\PYGZus{}max} \PYG{n}{function}\PYG{p}{)}

       \PYG{o}{\PYGZhy{}}\PYG{n}{o}  \PYG{n}{Generate} \PYG{n}{overlays}

       \PYG{o}{\PYGZhy{}}\PYG{n}{P}  \PYG{n}{Primitive} \PYG{n}{calculation} \PYG{n}{only}

       \PYG{o}{\PYGZhy{}}\PYG{n}{r}  \PYG{n}{Pass} \PYG{n}{TGV} \PYG{n}{pointer} \PYG{k}{as} \PYG{n}{register} \PYG{o+ow}{in} \PYG{n}{rules}

       \PYG{o}{\PYGZhy{}}\PYG{n}{R}  \PYG{n}{Set} \PYG{n}{application} \PYG{n}{to} \PYG{n}{both} \PYG{l+s+s2}{\PYGZdq{}}\PYG{l+s+s2}{iliad}\PYG{l+s+s2}{\PYGZdq{}} \PYG{o+ow}{and} \PYG{l+s+s2}{\PYGZdq{}}\PYG{l+s+s2}{pro}\PYG{l+s+s2}{\PYGZdq{}}

       \PYG{o}{\PYGZhy{}}\PYG{n}{s}  \PYG{n}{Strip} \PYG{n}{comments} \PYG{k+kn}{from}\PYG{+w}{ }\PYG{n+nn}{generated} \PYG{n}{output}

       \PYG{o}{\PYGZhy{}}\PYG{n}{S}  \PYG{n}{Generate} \PYG{n}{separate} \PYG{n}{controls}

       \PYG{o}{\PYGZhy{}}\PYG{n}{t}  \PYG{n}{Generate} \PYG{n}{trace} \PYG{n}{code}

       \PYG{o}{\PYGZhy{}}\PYG{n}{U}  \PYG{n}{Set} \PYG{n}{application} \PYG{n}{to} \PYG{l+s+s2}{\PYGZdq{}}\PYG{l+s+s2}{cfir}\PYG{l+s+s2}{\PYGZdq{}}

       \PYG{o}{\PYGZhy{}}\PYG{n}{x}  \PYG{n}{Generate} \PYG{n}{cross} \PYG{n}{references}

       \PYG{o}{\PYGZhy{}}\PYG{n}{X}  \PYG{n}{Generate} \PYG{k}{global} \PYG{n}{extraction}

       \PYG{o}{\PYGZhy{}}\PYG{n}{Z}  \PYG{n}{Colored} \PYG{n}{output} \PYG{o+ow}{in} \PYG{n}{chainings}
\end{sphinxVerbatim}


\subsection{Comportement de Mlang}
\label{\detokenize{mlang:comportement-de-mlang}}
\sphinxAtStartPar
Le compilateur Mlang effectue son traitement en quatre étapes :
\begin{itemize}
\item {} 
\sphinxAtStartPar
la traduction dans un format abstrait interne;

\item {} 
\sphinxAtStartPar
un pré\sphinxhyphen{}traitement pour le simplifier;

\item {} 
\sphinxAtStartPar
une vérification pour analyser la cohérence du code;

\item {} 
\sphinxAtStartPar
le traitement du code M, que ce soit son interprétation ou sa compilation.

\end{itemize}

\sphinxAtStartPar
Le module \sphinxcode{\sphinxupquote{Driver}} (et plus précisément la fonction \sphinxcode{\sphinxupquote{Driver.main}}) correspond au
point d’entrée de mlang.


\subsubsection{Traduction}
\label{\detokenize{mlang:traduction}}
\sphinxAtStartPar
Le langage M est parsé selon les règles spécifiées dans {\hyperref[\detokenize{syntax:syntax}]{\sphinxcrossref{\DUrole{std}{\DUrole{std-ref}{La syntaxe du M}}}}}.
Elle est effecuée par les modules \sphinxcode{\sphinxupquote{Mparser}}, \sphinxcode{\sphinxupquote{Mlexer}} et \sphinxcode{\sphinxupquote{Parse\_utils}}.


\subsubsection{Pré\sphinxhyphen{}traitement}
\label{\detokenize{mlang:pre-traitement}}
\sphinxAtStartPar
Le prétraitement est une opération purement syntaxique. Elle est effectuée par
les modules \sphinxcode{\sphinxupquote{Expander}} et \sphinxcode{\sphinxupquote{Mir}}. Son but est triple :
\begin{itemize}
\item {} 
\sphinxAtStartPar
éliminer les constructions relatives aux applications non\sphinxhyphen{}sélectionnées ;

\item {} 
\sphinxAtStartPar
remplacer les constantes par leur valeur numérique ;

\item {} 
\sphinxAtStartPar
remplacer les expressions numériques débutant par \sphinxcode{\sphinxupquote{somme}} avec des
additions ;

\item {} 
\sphinxAtStartPar
éliminer les \sphinxcode{\sphinxupquote{\textless{}multi\sphinxhyphen{}formule\textgreater{}}}s en les remplaçant par des séries de \sphinxcode{\sphinxupquote{\textless{}formule\textgreater{}}}s.

\end{itemize}

\sphinxAtStartPar
Toute substitution transformant un programme M syntaxiquement valide en un
texte ne correspondant à aucun programme M provoque l’échec du traitement.


\paragraph{Cas des applications}
\label{\detokenize{mlang:cas-des-applications}}
\sphinxAtStartPar
On élimine du programme M:
\begin{itemize}
\item {} 
\sphinxAtStartPar
les déclarations des applications non\sphinxhyphen{}sélectionnées ;

\item {} 
\sphinxAtStartPar
les déclarations des enchaîneurs ne spécifiant pas une application
sélectionnée ;

\item {} 
\sphinxAtStartPar
les déclarations des règles ne spécifiant pas une application
sélectionnée ;

\item {} 
\sphinxAtStartPar
les déclarations des vérifications ne spécifiant pas une application
sélectionnée ;

\item {} 
\sphinxAtStartPar
les déclarations des cibles ne spécifiant pas une application sélectionnée.

\end{itemize}

\sphinxAtStartPar
Dans les déclarations des règles spécifiant une application sélectionnée, on
élimine les enchaîneurs qui ne sont plus déclarés dans le programme M obtenu.


\paragraph{Cas des constantes}
\label{\detokenize{mlang:cas-des-constantes}}
\sphinxAtStartPar
Pour prétraiter un programme, on le parcourt du début à la fin. Pour chaque
\sphinxcode{\sphinxupquote{\textless{}variable\textgreater{}}} rencontrée, si elle correspond à une contante défini précédemment,
alors il est remplacé dans le programme par la valeur numérique correspondante.

\sphinxAtStartPar
Un intervalle de la forme \sphinxcode{\sphinxupquote{\textless{}naturel:début\textgreater{}..\textless{}symbole:const\textgreater{}}} est
converti par substitution de la constante \sphinxcode{\sphinxupquote{const}} en l’intervalle
\sphinxcode{\sphinxupquote{\textless{}naturel:début\textgreater{}..\textless{}naturel:fin\textgreater{}}}, avec \sphinxcode{\sphinxupquote{fin}} le naturel correspondant
à \sphinxcode{\sphinxupquote{const}}.

\sphinxAtStartPar
\sphinxstyleemphasis{Exemple} : considérons le programme suivant :

\begin{sphinxVerbatim}[commandchars=\\\{\}]
fin : const = 10;
… 
pour i = 9\PYGZhy{}fin :
  Bi = Bi + fin;
\end{sphinxVerbatim}

\sphinxAtStartPar
Par substitution de la constante \sphinxcode{\sphinxupquote{fin}}, il est remplacé dans un premier temps
remplacée par le programme :

\begin{sphinxVerbatim}[commandchars=\\\{\}]
fin : const = 10;`
… 
pour i = 9..10 :
  Bi = Bi + fin;
\end{sphinxVerbatim}

\sphinxAtStartPar
puis dans un second temps par le programme :

\begin{sphinxVerbatim}[commandchars=\\\{\}]
fin : const = 10;
…
pour i = 9..10 :
  Bi = Bi + 10;
\end{sphinxVerbatim}


\paragraph{Interprétation des indices}
\label{\detokenize{mlang:interpretation-des-indices}}
\sphinxAtStartPar
Pour rappel, les \sphinxstyleemphasis{indices} ont la forme suivante :

\begin{sphinxVerbatim}[commandchars=\\\{\}]
\PYGZlt{}indices\PYGZgt{} ::= \PYGZlt{}indice\PYGZgt{} ; …
\PYGZlt{}indice\PYGZgt{} ::= \PYGZlt{}minuscule\PYGZgt{} = \PYGZlt{}intervalle\PYGZgt{} , …
\end{sphinxVerbatim}

\sphinxAtStartPar
avec :
\begin{itemize}
\item {} 
\sphinxAtStartPar
\sphinxcode{\sphinxupquote{\textless{}minuscule\textgreater{} ::= {[}a\sphinxhyphen{}z{]}}}

\item {} 
\sphinxAtStartPar
\sphinxcode{\sphinxupquote{\textless{}majuscule\textgreater{} ::= {[}A\sphinxhyphen{}Z{]}}}

\item {} 
\sphinxAtStartPar
\sphinxcode{\sphinxupquote{\textless{}intervalle\textgreater{} ::= \textless{}majuscule\textgreater{}+ | \textless{}majuscule\textgreater{} .. \textless{}majuscule\textgreater{} | \textless{}naturel\textgreater{} (.. \textless{}naturel\textgreater{} | \sphinxhyphen{} \textless{}variable\textgreater{})?}}

\end{itemize}

\sphinxAtStartPar
Chaque \sphinxcode{\sphinxupquote{\textless{}indice\textgreater{}}} associe à une lettre minuscule une série de chaînes de
caractères de même taille.
Cette série est composée de la succession de chaque
 à droite du symbole \sphinxcode{\sphinxupquote{=}}.

\sphinxAtStartPar
Les séries associées aux  sont définis comme suit :
\begin{itemize}
\item {} 
\sphinxAtStartPar
\sphinxcode{\sphinxupquote{\textless{}majuscule\textgreater{}+}} est la série composée de chacune des majuscules prises
séparément, donc de taille 1 (\sphinxstyleemphasis{par exemple} : \sphinxcode{\sphinxupquote{AXF}} représente la série \sphinxcode{\sphinxupquote{A}}, \sphinxcode{\sphinxupquote{X}}, \sphinxcode{\sphinxupquote{F}});

\item {} 
\sphinxAtStartPar
\sphinxcode{\sphinxupquote{\textless{}majuscule:/début/\textgreater{}****\textless{}majuscule:/fin/\textgreater{}}} est la série composée de toutes
les majuscules comprises entre /début/ et /fin/, bornes comprises, suivant
l’ordre alphabétique, donc de taille 1 (\sphinxstyleemphasis{par exemple} : \sphinxcode{\sphinxupquote{A..D}} représente la série \sphinxcode{\sphinxupquote{A}}, \sphinxcode{\sphinxupquote{B}}, \sphinxcode{\sphinxupquote{C}}, \sphinxcode{\sphinxupquote{D}});

\item {} 
\sphinxAtStartPar
\sphinxcode{\sphinxupquote{\textless{}naturel:début\textgreater{}..\textless{}naturel:fin\textgreater{}}} est la série composée de tous les
nombres naturels entre \sphinxcode{\sphinxupquote{début}} et \sphinxcode{\sphinxupquote{fin}}, bornes comprises, la taille des
éléments étant égale à la taille du plus grand naturel en base 10; les
naturels trop petits pour avoir la taille requise  sont complétés par des 0 à
gauche (\sphinxstyleemphasis{par exemple} : \sphinxcode{\sphinxupquote{9..11}} représente la série \sphinxcode{\sphinxupquote{09}}, \sphinxcode{\sphinxupquote{10}}, \sphinxcode{\sphinxupquote{11}});

\item {} 
\sphinxAtStartPar
\sphinxcode{\sphinxupquote{\textless{}naturel\textgreater{}\sphinxhyphen{}\textless{}variable\textgreater{}}} est converti lors du prétraitement des constantes
en un intervalle de la forme \sphinxcode{\sphinxupquote{\textless{}naturel\textgreater{}..\textless{}naturel\textgreater{}}}.

\end{itemize}


\paragraph{Cas des expressions \sphinxstyleliteralintitle{\sphinxupquote{somme(…)}}}
\label{\detokenize{mlang:cas-des-expressions-somme}}
\sphinxAtStartPar
Pour une expression numérique
\sphinxcode{\sphinxupquote{somme ( \textless{}indice:ind\textgreater{}: \textless{}expression numerique:expr\textgreater{} )}},
l’expression \sphinxstyleemphasis{expr} sera remplacée par la somme des expressions construites
ainsi : pour chaque chaîne de caractères de la série introduite par \sphinxstyleemphasis{ind}, les
symboles apparaîssant dans \sphinxstyleemphasis{expr} sont remplacés par ceux dans lesquels la
lettre minuscule de /ind/ est substituéé par la chaine de caractère.

\sphinxAtStartPar
La somme ainsi générée est parenthésée.

\sphinxAtStartPar
Une expression numérique
\sphinxcode{\sphinxupquote{somme ( \textless{}indice:ind\textgreater{} ; \textless{}indices:inds\textgreater{} : \textless{}expression numérique:expr\textgreater{} )}},
est remplacée par la somme des expressions
\sphinxstyleemphasis{\sphinxcode{\sphinxupquote{somme}}}\sphinxcode{\sphinxupquote{ \textless{}indices:inds\textgreater{} : expr\textquotesingle{}\textgreater{} )}} avec \sphinxcode{\sphinxupquote{expr\textquotesingle{}}} prenant sa valeur
dans la série \sphinxcode{\sphinxupquote{sub(ind, expr)}}.

\sphinxAtStartPar
La procédure est ensuite appliquée récursivement à cette somme.

\sphinxAtStartPar
La somme ainsi générée est parenthésée.

\sphinxAtStartPar
Notons que les sous\sphinxhyphen{}sommes récursivement produites n’ont pas besoin de l’être
car l’addition est associative.

\sphinxAtStartPar
On applique cette transformation à toutes les expressions \sphinxcode{\sphinxupquote{somme(…)}} tant
qu’il en existe dans le programme.

\sphinxAtStartPar
\sphinxstylestrong{Exemple}. Considérons l’expression suivante :
\begin{itemize}
\item {} 
\sphinxAtStartPar
\sphinxcode{\sphinxupquote{somme(i = XZ ; j = 9..10 : Bi + Bj)}}

\end{itemize}

\sphinxAtStartPar
Elle est dans un premier temps remplacée par la somme suivante :
\begin{itemize}
\item {} 
\sphinxAtStartPar
\sphinxcode{\sphinxupquote{(somme(j = 9..10 : BX + Bj) + somme(j = 9..10 : BZ + Bj))}}

\end{itemize}

\sphinxAtStartPar
puis dans un second temps par la somme :
\begin{itemize}
\item {} 
\sphinxAtStartPar
\sphinxcode{\sphinxupquote{(BX + B09 + BX + B10 + BZ + B09 + BZ + B10)}}

\end{itemize}

\sphinxAtStartPar
On remarque que seule l’expression globale est parenthésée car l’addition est
associative.


\paragraph{Cas des multi\sphinxhyphen{}formules}
\label{\detokenize{mlang:cas-des-multi-formules}}
\sphinxAtStartPar
On commence par substituer toutes les expressions \sphinxcode{\sphinxupquote{somme(…)}} apparaissant dans
les multi\sphinxhyphen{}formules.

\sphinxAtStartPar
Puis on transforme les multi\sphinxhyphen{}formules en séries de formules en suivant la
méthode décrite ci\sphinxhyphen{}dessous.

\sphinxAtStartPar
Pour chaque multi\sphinxhyphen{}formule rencontrée, de la forme \sphinxcode{\sphinxupquote{pour \textless{}indices\textgreater{} : \textless{}formule\textgreater{}}},
dès que les constantes apparaîssant dans les  ont été substitués, on
remplace cette multi\sphinxhyphen{}formule par la série de formules générée de la manière
suivante.

\sphinxAtStartPar
Pour une multi\sphinxhyphen{}formule \sphinxcode{\sphinxupquote{pour \textless{}indice:ind\textgreater{} : \textless{}formule:f\textgreater{}}}, la
formule \sphinxcode{\sphinxupquote{f}} sera remplacée par la série des formules construites ainsi : pour
chaque chaîne de caractères de la série introduite par \sphinxcode{\sphinxupquote{ind}}, les symboles
apparaîssant dans \sphinxcode{\sphinxupquote{f}} sont remplacés par ceux dans lesquels la lettre
minuscule de \sphinxcode{\sphinxupquote{ind}} est substituéé par la chaïne de caractère.

\sphinxAtStartPar
On notera \sphinxcode{\sphinxupquote{sub(ind, f)}} la série constituée des formules ainsi régérées.

\sphinxAtStartPar
Une multi\sphinxhyphen{}formule \sphinxcode{\sphinxupquote{pour \textless{}indice:ind\textgreater{} ; \textless{}indices:inds\textgreater{} : \textless{}formule:f\textgreater{}}},
est remplacée par la série des multi\sphinxhyphen{}formules
\sphinxcode{\sphinxupquote{pour \textless{}indices:inds\textgreater{} : f\textquotesingle{}\textgreater{}}} avec \sphinxcode{\sphinxupquote{f\textquotesingle{}}} prenant sa valeur
dans la série \sphinxcode{\sphinxupquote{sub(ind, f)}}.

\sphinxAtStartPar
La procédure est ensuite appliquée récursivement
à ces multi\sphinxhyphen{}formules.

\sphinxAtStartPar
\sphinxstylestrong{Exemple}. Considérons la multi\sphinxhyphen{}formule suivante :

\begin{sphinxVerbatim}[commandchars=\\\{\}]
\PYG{n}{pour} \PYG{n}{i} \PYG{o}{=} \PYG{n}{XZ} \PYG{p}{;} \PYG{n}{j} \PYG{o}{=} \PYG{l+m+mf}{9.}\PYG{l+m+mf}{.10} \PYG{p}{:}
  \PYG{n}{Bi} \PYG{o}{=} \PYG{n}{Bi} \PYG{o}{+} \PYG{n}{Bj}\PYG{p}{;}
\end{sphinxVerbatim}

\sphinxAtStartPar
Elle est dans un premier temps remplacée par la série suivante :

\begin{sphinxVerbatim}[commandchars=\\\{\}]
\PYG{n}{pour} \PYG{n}{j} \PYG{o}{=} \PYG{l+m+mf}{9.}\PYG{l+m+mf}{.10} \PYG{p}{:}
  \PYG{n}{BX} \PYG{o}{=} \PYG{n}{BX} \PYG{o}{+} \PYG{n}{Bj}\PYG{p}{;}
 
\PYG{n}{pour} \PYG{n}{j} \PYG{o}{=} \PYG{l+m+mf}{9.}\PYG{l+m+mf}{.10} \PYG{p}{:}
  \PYG{n}{BZ} \PYG{o}{=} \PYG{n}{BZ} \PYG{o}{+} \PYG{n}{Bj}\PYG{p}{;}
\end{sphinxVerbatim}

\sphinxAtStartPar
Puis dans un second temps par la série des formules :

\begin{sphinxVerbatim}[commandchars=\\\{\}]
\PYG{n}{BX} \PYG{o}{=} \PYG{n}{BX} \PYG{o}{+} \PYG{n}{B09}\PYG{p}{;}
\PYG{n}{BX} \PYG{o}{=} \PYG{n}{BX} \PYG{o}{+} \PYG{n}{B10}\PYG{p}{;}
\PYG{n}{BZ} \PYG{o}{=} \PYG{n}{BZ} \PYG{o}{+} \PYG{n}{B09}\PYG{p}{;}
\PYG{n}{BZ} \PYG{o}{=} \PYG{n}{BZ} \PYG{o}{+} \PYG{n}{B10}\PYG{p}{;}
\end{sphinxVerbatim}


\subsubsection{Vérification de cohérence}
\label{\detokenize{mlang:verification-de-coherence}}
\sphinxAtStartPar
De nombreuses constructions sont valides à la traduction, mais brisent
certains invariants nécessaires à la bonne exécution du code : double
déclaration d’attributs, nom de variable déjà utilisé, variable mal
typée…
L’ensemble de ces vérifications est accessible dans le module
\sphinxcode{\sphinxupquote{Validator}} du frontend. Le sous module \sphinxcode{\sphinxupquote{Validator.Err}} définit l’ensemble
des erreurs levées par cette étape de vérification.

\sphinxAtStartPar
\sphinxstylestrong{NB} : seule la première erreur rencontrée par le validateur est levée.


\subsubsection{Traitement}
\label{\detokenize{mlang:traitement}}

\paragraph{Interpreteur}
\label{\detokenize{mlang:interpreteur}}
\sphinxAtStartPar
L’interpréteur utilise la représentation interne du code M
pour lancer le calcul à partir d’un fichier IRJ
(voir {\hyperref[\detokenize{syntax_irj:syntax-irj}]{\sphinxcrossref{\DUrole{std}{\DUrole{std-ref}{Les fichiers IRJ}}}}}).
L’interprétation est effecuée par le module \sphinxcode{\sphinxupquote{Test\_interpreter}}.


\paragraph{Transpilation}
\label{\detokenize{mlang:transpilation}}
\sphinxAtStartPar
La transpilation traduit le code dans le langage spécifié (en 2025, seul le C
est transpilable). La transpilation est effecutée dans le module
\sphinxcode{\sphinxupquote{Bir\_to\_dgfip\_c}}.


\chapter{Exemples}
\label{\detokenize{index:exemples}}
\sphinxstepscope


\section{Calcul des valeurs}
\label{\detokenize{exemples/valeurs:calcul-des-valeurs}}\label{\detokenize{exemples/valeurs:exemples-valeurs}}\label{\detokenize{exemples/valeurs::doc}}

\subsection{Fichier test : test.m}
\label{\detokenize{exemples/valeurs:fichier-test-test-m}}
\begin{sphinxVerbatim}[commandchars=\\\{\}]
\PYG{n}{espace\PYGZus{}variables} \PYG{n}{GLOBAL} \PYG{p}{:} \PYG{n}{par\PYGZus{}defaut}\PYG{p}{;}
\PYG{n}{domaine} \PYG{n}{regle} \PYG{n}{mon\PYGZus{}domaine\PYGZus{}de\PYGZus{}regle}\PYG{p}{:} \PYG{n}{par\PYGZus{}defaut}\PYG{p}{;}
\PYG{n}{domaine} \PYG{n}{verif} \PYG{n}{mon\PYGZus{}domaine\PYGZus{}de\PYGZus{}verifications}\PYG{p}{:} \PYG{n}{par\PYGZus{}defaut}\PYG{p}{;}
\PYG{n}{application} \PYG{n}{mon\PYGZus{}application}\PYG{p}{;}
\PYG{n}{variable} \PYG{n}{saisie} \PYG{p}{:} \PYG{n}{attribut} \PYG{n}{mon\PYGZus{}attribut}\PYG{p}{;}
\PYG{n}{variable} \PYG{n}{calculee} \PYG{p}{:} \PYG{n}{attribut} \PYG{n}{mon\PYGZus{}attribut}\PYG{p}{;}
\PYG{n}{X} \PYG{p}{:} \PYG{n}{saisie} \PYG{n}{mon\PYGZus{}attribut} \PYG{o}{=} \PYG{l+m+mi}{0} \PYG{n}{alias} \PYG{n}{AX} \PYG{p}{:} \PYG{l+s+s2}{\PYGZdq{}}\PYG{l+s+s2}{\PYGZdq{}}\PYG{p}{;}

\PYG{n}{cible} \PYG{n}{indefini\PYGZus{}test}\PYG{p}{:}
\PYG{n}{application} \PYG{p}{:} \PYG{n}{mon\PYGZus{}application}\PYG{p}{;}
\PYG{n}{afficher} \PYG{l+s+s2}{\PYGZdq{}}\PYG{l+s+s2}{Bonjour, monde, X = }\PYG{l+s+s2}{\PYGZdq{}}\PYG{p}{;}
\PYG{n}{afficher} \PYG{p}{(}\PYG{n}{X}\PYG{p}{)}\PYG{p}{;}
\PYG{n}{afficher} \PYG{l+s+s2}{\PYGZdq{}}\PYG{l+s+s2}{ !}\PYG{l+s+se}{\PYGZbs{}n}\PYG{l+s+s2}{\PYGZdq{}}\PYG{p}{;}
\PYG{c+c1}{\PYGZsh{} Addition}
\PYG{n}{afficher} \PYG{l+s+s2}{\PYGZdq{}}\PYG{l+s+s2}{X + 1 = }\PYG{l+s+s2}{\PYGZdq{}}\PYG{p}{;}
\PYG{n}{afficher} \PYG{p}{(}\PYG{n}{X} \PYG{o}{+} \PYG{l+m+mi}{1}\PYG{p}{)}\PYG{p}{;}
\PYG{n}{afficher} \PYG{l+s+s2}{\PYGZdq{}}\PYG{l+s+s2}{ !}\PYG{l+s+se}{\PYGZbs{}n}\PYG{l+s+s2}{\PYGZdq{}}\PYG{p}{;}
\PYG{n}{afficher} \PYG{l+s+s2}{\PYGZdq{}}\PYG{l+s+s2}{1 + X = }\PYG{l+s+s2}{\PYGZdq{}}\PYG{p}{;}
\PYG{n}{afficher} \PYG{p}{(}\PYG{l+m+mi}{1} \PYG{o}{+} \PYG{n}{X}\PYG{p}{)}\PYG{p}{;}
\PYG{n}{afficher} \PYG{l+s+s2}{\PYGZdq{}}\PYG{l+s+s2}{ !}\PYG{l+s+se}{\PYGZbs{}n}\PYG{l+s+s2}{\PYGZdq{}}\PYG{p}{;}
\PYG{n}{afficher} \PYG{l+s+s2}{\PYGZdq{}}\PYG{l+s+s2}{X + X = }\PYG{l+s+s2}{\PYGZdq{}}\PYG{p}{;}
\PYG{n}{afficher} \PYG{p}{(}\PYG{n}{X} \PYG{o}{+} \PYG{n}{X}\PYG{p}{)}\PYG{p}{;}
\PYG{n}{afficher} \PYG{l+s+s2}{\PYGZdq{}}\PYG{l+s+s2}{ !}\PYG{l+s+se}{\PYGZbs{}n}\PYG{l+s+s2}{\PYGZdq{}}\PYG{p}{;}
\PYG{c+c1}{\PYGZsh{} Soustraction}
\PYG{n}{afficher} \PYG{l+s+s2}{\PYGZdq{}}\PYG{l+s+s2}{X \PYGZhy{} 1 = }\PYG{l+s+s2}{\PYGZdq{}}\PYG{p}{;}
\PYG{n}{afficher} \PYG{p}{(}\PYG{n}{X} \PYG{o}{\PYGZhy{}} \PYG{l+m+mi}{1}\PYG{p}{)}\PYG{p}{;}
\PYG{n}{afficher} \PYG{l+s+s2}{\PYGZdq{}}\PYG{l+s+s2}{ !}\PYG{l+s+se}{\PYGZbs{}n}\PYG{l+s+s2}{\PYGZdq{}}\PYG{p}{;}
\PYG{n}{afficher} \PYG{l+s+s2}{\PYGZdq{}}\PYG{l+s+s2}{1 \PYGZhy{} X = }\PYG{l+s+s2}{\PYGZdq{}}\PYG{p}{;}
\PYG{n}{afficher} \PYG{p}{(}\PYG{l+m+mi}{1} \PYG{o}{\PYGZhy{}} \PYG{n}{X}\PYG{p}{)}\PYG{p}{;}
\PYG{n}{afficher} \PYG{l+s+s2}{\PYGZdq{}}\PYG{l+s+s2}{ !}\PYG{l+s+se}{\PYGZbs{}n}\PYG{l+s+s2}{\PYGZdq{}}\PYG{p}{;}
\PYG{n}{afficher} \PYG{l+s+s2}{\PYGZdq{}}\PYG{l+s+s2}{X \PYGZhy{} X = }\PYG{l+s+s2}{\PYGZdq{}}\PYG{p}{;}
\PYG{n}{afficher} \PYG{p}{(}\PYG{n}{X} \PYG{o}{\PYGZhy{}} \PYG{n}{X}\PYG{p}{)}\PYG{p}{;}
\PYG{n}{afficher} \PYG{l+s+s2}{\PYGZdq{}}\PYG{l+s+s2}{ !}\PYG{l+s+se}{\PYGZbs{}n}\PYG{l+s+s2}{\PYGZdq{}}\PYG{p}{;}
\PYG{c+c1}{\PYGZsh{} Multiplication}
\PYG{n}{afficher} \PYG{l+s+s2}{\PYGZdq{}}\PYG{l+s+s2}{X * 1 = }\PYG{l+s+s2}{\PYGZdq{}}\PYG{p}{;}
\PYG{n}{afficher} \PYG{p}{(}\PYG{n}{X} \PYG{o}{*} \PYG{l+m+mi}{1}\PYG{p}{)}\PYG{p}{;}
\PYG{n}{afficher} \PYG{l+s+s2}{\PYGZdq{}}\PYG{l+s+s2}{ !}\PYG{l+s+se}{\PYGZbs{}n}\PYG{l+s+s2}{\PYGZdq{}}\PYG{p}{;}
\PYG{n}{afficher} \PYG{l+s+s2}{\PYGZdq{}}\PYG{l+s+s2}{1 * X = }\PYG{l+s+s2}{\PYGZdq{}}\PYG{p}{;}
\PYG{n}{afficher} \PYG{p}{(}\PYG{l+m+mi}{1} \PYG{o}{*} \PYG{n}{X}\PYG{p}{)}\PYG{p}{;}
\PYG{n}{afficher} \PYG{l+s+s2}{\PYGZdq{}}\PYG{l+s+s2}{ !}\PYG{l+s+se}{\PYGZbs{}n}\PYG{l+s+s2}{\PYGZdq{}}\PYG{p}{;}
\PYG{n}{afficher} \PYG{l+s+s2}{\PYGZdq{}}\PYG{l+s+s2}{X * X = }\PYG{l+s+s2}{\PYGZdq{}}\PYG{p}{;}
\PYG{n}{afficher} \PYG{p}{(}\PYG{n}{X} \PYG{o}{*} \PYG{n}{X}\PYG{p}{)}\PYG{p}{;}
\PYG{n}{afficher} \PYG{l+s+s2}{\PYGZdq{}}\PYG{l+s+s2}{ !}\PYG{l+s+se}{\PYGZbs{}n}\PYG{l+s+s2}{\PYGZdq{}}\PYG{p}{;}
\PYG{c+c1}{\PYGZsh{} Division}
\PYG{n}{afficher} \PYG{l+s+s2}{\PYGZdq{}}\PYG{l+s+s2}{X / 1 = }\PYG{l+s+s2}{\PYGZdq{}}\PYG{p}{;}
\PYG{n}{afficher} \PYG{p}{(}\PYG{n}{X} \PYG{o}{/} \PYG{l+m+mi}{1}\PYG{p}{)}\PYG{p}{;}
\PYG{n}{afficher} \PYG{l+s+s2}{\PYGZdq{}}\PYG{l+s+s2}{ !}\PYG{l+s+se}{\PYGZbs{}n}\PYG{l+s+s2}{\PYGZdq{}}\PYG{p}{;}
\PYG{n}{afficher} \PYG{l+s+s2}{\PYGZdq{}}\PYG{l+s+s2}{1 / X = }\PYG{l+s+s2}{\PYGZdq{}}\PYG{p}{;}
\PYG{n}{afficher} \PYG{p}{(}\PYG{l+m+mi}{1} \PYG{o}{/} \PYG{n}{X}\PYG{p}{)}\PYG{p}{;}
\PYG{n}{afficher} \PYG{l+s+s2}{\PYGZdq{}}\PYG{l+s+s2}{ !}\PYG{l+s+se}{\PYGZbs{}n}\PYG{l+s+s2}{\PYGZdq{}}\PYG{p}{;}
\PYG{n}{afficher} \PYG{l+s+s2}{\PYGZdq{}}\PYG{l+s+s2}{1 / 0 = }\PYG{l+s+s2}{\PYGZdq{}}\PYG{p}{;}
\PYG{n}{afficher} \PYG{p}{(}\PYG{l+m+mi}{1} \PYG{o}{/} \PYG{l+m+mi}{0}\PYG{p}{)}\PYG{p}{;}
\PYG{n}{afficher} \PYG{l+s+s2}{\PYGZdq{}}\PYG{l+s+s2}{ !}\PYG{l+s+se}{\PYGZbs{}n}\PYG{l+s+s2}{\PYGZdq{}}\PYG{p}{;}
\PYG{n}{afficher} \PYG{l+s+s2}{\PYGZdq{}}\PYG{l+s+s2}{X / 0 = }\PYG{l+s+s2}{\PYGZdq{}}\PYG{p}{;}
\PYG{n}{afficher} \PYG{p}{(}\PYG{n}{X} \PYG{o}{/} \PYG{l+m+mi}{0}\PYG{p}{)}\PYG{p}{;}
\PYG{n}{afficher} \PYG{l+s+s2}{\PYGZdq{}}\PYG{l+s+s2}{ !}\PYG{l+s+se}{\PYGZbs{}n}\PYG{l+s+s2}{\PYGZdq{}}\PYG{p}{;}
\PYG{n}{afficher} \PYG{l+s+s2}{\PYGZdq{}}\PYG{l+s+s2}{X / X = }\PYG{l+s+s2}{\PYGZdq{}}\PYG{p}{;}
\PYG{n}{afficher} \PYG{p}{(}\PYG{n}{X} \PYG{o}{/} \PYG{n}{X}\PYG{p}{)}\PYG{p}{;}
\PYG{n}{afficher} \PYG{l+s+s2}{\PYGZdq{}}\PYG{l+s+s2}{ !}\PYG{l+s+se}{\PYGZbs{}n}\PYG{l+s+s2}{\PYGZdq{}}\PYG{p}{;}
\PYG{c+c1}{\PYGZsh{} Comparaisons}
\PYG{n}{afficher} \PYG{l+s+s2}{\PYGZdq{}}\PYG{l+s+s2}{(X = 0) = }\PYG{l+s+s2}{\PYGZdq{}}\PYG{p}{;}
\PYG{n}{afficher} \PYG{p}{(}\PYG{n}{X} \PYG{o}{=} \PYG{l+m+mi}{0}\PYG{p}{)}\PYG{p}{;}
\PYG{n}{afficher} \PYG{l+s+s2}{\PYGZdq{}}\PYG{l+s+s2}{ !}\PYG{l+s+se}{\PYGZbs{}n}\PYG{l+s+s2}{\PYGZdq{}}\PYG{p}{;}
\PYG{n}{afficher} \PYG{l+s+s2}{\PYGZdq{}}\PYG{l+s+s2}{(X = X) = }\PYG{l+s+s2}{\PYGZdq{}}\PYG{p}{;}
\PYG{n}{afficher} \PYG{p}{(}\PYG{n}{X} \PYG{o}{=} \PYG{n}{X}\PYG{p}{)}\PYG{p}{;}
\PYG{n}{afficher} \PYG{l+s+s2}{\PYGZdq{}}\PYG{l+s+s2}{ !}\PYG{l+s+se}{\PYGZbs{}n}\PYG{l+s+s2}{\PYGZdq{}}\PYG{p}{;}
\PYG{n}{afficher} \PYG{l+s+s2}{\PYGZdq{}}\PYG{l+s+s2}{(X \PYGZgt{}= 0) = }\PYG{l+s+s2}{\PYGZdq{}}\PYG{p}{;}
\PYG{n}{afficher} \PYG{p}{(}\PYG{n}{X} \PYG{o}{\PYGZgt{}}\PYG{o}{=} \PYG{l+m+mi}{0}\PYG{p}{)}\PYG{p}{;}
\PYG{n}{afficher} \PYG{l+s+s2}{\PYGZdq{}}\PYG{l+s+s2}{ !}\PYG{l+s+se}{\PYGZbs{}n}\PYG{l+s+s2}{\PYGZdq{}}\PYG{p}{;}
\PYG{n}{afficher} \PYG{l+s+s2}{\PYGZdq{}}\PYG{l+s+s2}{(X \PYGZgt{}= X) = }\PYG{l+s+s2}{\PYGZdq{}}\PYG{p}{;}
\PYG{n}{afficher} \PYG{p}{(}\PYG{n}{X} \PYG{o}{\PYGZgt{}}\PYG{o}{=} \PYG{n}{X}\PYG{p}{)}\PYG{p}{;}
\PYG{n}{afficher} \PYG{l+s+s2}{\PYGZdq{}}\PYG{l+s+s2}{ !}\PYG{l+s+se}{\PYGZbs{}n}\PYG{l+s+s2}{\PYGZdq{}}\PYG{p}{;}
\PYG{n}{afficher} \PYG{l+s+s2}{\PYGZdq{}}\PYG{l+s+s2}{(X \PYGZgt{} 0) = }\PYG{l+s+s2}{\PYGZdq{}}\PYG{p}{;}
\PYG{n}{afficher} \PYG{p}{(}\PYG{n}{X} \PYG{o}{\PYGZgt{}} \PYG{l+m+mi}{0}\PYG{p}{)}\PYG{p}{;}
\PYG{n}{afficher} \PYG{l+s+s2}{\PYGZdq{}}\PYG{l+s+s2}{ !}\PYG{l+s+se}{\PYGZbs{}n}\PYG{l+s+s2}{\PYGZdq{}}\PYG{p}{;}
\PYG{n}{afficher} \PYG{l+s+s2}{\PYGZdq{}}\PYG{l+s+s2}{(X \PYGZgt{} X) = }\PYG{l+s+s2}{\PYGZdq{}}\PYG{p}{;}
\PYG{n}{afficher} \PYG{p}{(}\PYG{n}{X} \PYG{o}{\PYGZgt{}} \PYG{n}{X}\PYG{p}{)}\PYG{p}{;}
\PYG{n}{afficher} \PYG{l+s+s2}{\PYGZdq{}}\PYG{l+s+s2}{ !}\PYG{l+s+se}{\PYGZbs{}n}\PYG{l+s+s2}{\PYGZdq{}}\PYG{p}{;}
\PYG{n}{afficher} \PYG{l+s+s2}{\PYGZdq{}}\PYG{l+s+s2}{(X \PYGZlt{}= 0) = }\PYG{l+s+s2}{\PYGZdq{}}\PYG{p}{;}
\PYG{n}{afficher} \PYG{p}{(}\PYG{n}{X} \PYG{o}{\PYGZlt{}}\PYG{o}{=} \PYG{l+m+mi}{0}\PYG{p}{)}\PYG{p}{;}
\PYG{n}{afficher} \PYG{l+s+s2}{\PYGZdq{}}\PYG{l+s+s2}{ !}\PYG{l+s+se}{\PYGZbs{}n}\PYG{l+s+s2}{\PYGZdq{}}\PYG{p}{;}
\PYG{n}{afficher} \PYG{l+s+s2}{\PYGZdq{}}\PYG{l+s+s2}{(X \PYGZlt{} 0) = }\PYG{l+s+s2}{\PYGZdq{}}\PYG{p}{;}
\PYG{n}{afficher} \PYG{p}{(}\PYG{n}{X} \PYG{o}{\PYGZlt{}} \PYG{l+m+mi}{0}\PYG{p}{)}\PYG{p}{;}
\PYG{n}{afficher} \PYG{l+s+s2}{\PYGZdq{}}\PYG{l+s+s2}{ !}\PYG{l+s+se}{\PYGZbs{}n}\PYG{l+s+s2}{\PYGZdq{}}\PYG{p}{;}
\PYG{n}{afficher} \PYG{l+s+s2}{\PYGZdq{}}\PYG{l+s+s2}{(X \PYGZlt{}= X) = }\PYG{l+s+s2}{\PYGZdq{}}\PYG{p}{;}
\PYG{n}{afficher} \PYG{p}{(}\PYG{n}{X} \PYG{o}{\PYGZlt{}}\PYG{o}{=} \PYG{n}{X}\PYG{p}{)}\PYG{p}{;}
\PYG{n}{afficher} \PYG{l+s+s2}{\PYGZdq{}}\PYG{l+s+s2}{ !}\PYG{l+s+se}{\PYGZbs{}n}\PYG{l+s+s2}{\PYGZdq{}}\PYG{p}{;}
\PYG{n}{afficher} \PYG{l+s+s2}{\PYGZdq{}}\PYG{l+s+s2}{(X \PYGZlt{} X) = }\PYG{l+s+s2}{\PYGZdq{}}\PYG{p}{;}
\PYG{n}{afficher} \PYG{p}{(}\PYG{n}{X} \PYG{o}{\PYGZlt{}} \PYG{n}{X}\PYG{p}{)}\PYG{p}{;}
\PYG{n}{afficher} \PYG{l+s+s2}{\PYGZdq{}}\PYG{l+s+s2}{ !}\PYG{l+s+se}{\PYGZbs{}n}\PYG{l+s+s2}{\PYGZdq{}}\PYG{p}{;}
\PYG{c+c1}{\PYGZsh{} Opérations booléennes}
\PYG{n}{afficher} \PYG{l+s+s2}{\PYGZdq{}}\PYG{l+s+s2}{(X et X) = }\PYG{l+s+s2}{\PYGZdq{}}\PYG{p}{;}
\PYG{n}{afficher} \PYG{p}{(}\PYG{n}{X} \PYG{n}{et} \PYG{n}{X}\PYG{p}{)}\PYG{p}{;}
\PYG{n}{afficher} \PYG{l+s+s2}{\PYGZdq{}}\PYG{l+s+s2}{ !}\PYG{l+s+se}{\PYGZbs{}n}\PYG{l+s+s2}{\PYGZdq{}}\PYG{p}{;}
\PYG{n}{afficher} \PYG{l+s+s2}{\PYGZdq{}}\PYG{l+s+s2}{(X et 1) = }\PYG{l+s+s2}{\PYGZdq{}}\PYG{p}{;}
\PYG{n}{afficher} \PYG{p}{(}\PYG{n}{X} \PYG{n}{et} \PYG{l+m+mi}{1}\PYG{p}{)}\PYG{p}{;}
\PYG{n}{afficher} \PYG{l+s+s2}{\PYGZdq{}}\PYG{l+s+s2}{ !}\PYG{l+s+se}{\PYGZbs{}n}\PYG{l+s+s2}{\PYGZdq{}}\PYG{p}{;}
\PYG{n}{afficher} \PYG{l+s+s2}{\PYGZdq{}}\PYG{l+s+s2}{(X ou X) = }\PYG{l+s+s2}{\PYGZdq{}}\PYG{p}{;}
\PYG{n}{afficher} \PYG{p}{(}\PYG{n}{X} \PYG{n}{ou} \PYG{n}{X}\PYG{p}{)}\PYG{p}{;}
\PYG{n}{afficher} \PYG{l+s+s2}{\PYGZdq{}}\PYG{l+s+s2}{ !}\PYG{l+s+se}{\PYGZbs{}n}\PYG{l+s+s2}{\PYGZdq{}}\PYG{p}{;}
\PYG{n}{afficher} \PYG{l+s+s2}{\PYGZdq{}}\PYG{l+s+s2}{(X ou 1) = }\PYG{l+s+s2}{\PYGZdq{}}\PYG{p}{;}
\PYG{n}{afficher} \PYG{p}{(}\PYG{n}{X} \PYG{n}{ou} \PYG{l+m+mi}{1}\PYG{p}{)}\PYG{p}{;}
\PYG{n}{afficher} \PYG{l+s+s2}{\PYGZdq{}}\PYG{l+s+s2}{ !}\PYG{l+s+se}{\PYGZbs{}n}\PYG{l+s+s2}{\PYGZdq{}}\PYG{p}{;}
\PYG{n}{afficher} \PYG{l+s+s2}{\PYGZdq{}}\PYG{l+s+s2}{non X = }\PYG{l+s+s2}{\PYGZdq{}}\PYG{p}{;}
\PYG{n}{afficher} \PYG{p}{(}\PYG{n}{non} \PYG{n}{X}\PYG{p}{)}\PYG{p}{;}
\PYG{n}{afficher} \PYG{l+s+s2}{\PYGZdq{}}\PYG{l+s+s2}{ !}\PYG{l+s+se}{\PYGZbs{}n}\PYG{l+s+s2}{\PYGZdq{}}\PYG{p}{;}
\PYG{c+c1}{\PYGZsh{} Vérifier la définition d\PYGZsq{}une valeur}
\PYG{n}{afficher} \PYG{l+s+s2}{\PYGZdq{}}\PYG{l+s+s2}{present(X) = }\PYG{l+s+s2}{\PYGZdq{}}\PYG{p}{;}
\PYG{n}{afficher} \PYG{p}{(}\PYG{n}{present}\PYG{p}{(}\PYG{n}{X}\PYG{p}{)}\PYG{p}{)}\PYG{p}{;}
\PYG{n}{afficher} \PYG{l+s+s2}{\PYGZdq{}}\PYG{l+s+s2}{ !}\PYG{l+s+se}{\PYGZbs{}n}\PYG{l+s+s2}{\PYGZdq{}}\PYG{p}{;}
\end{sphinxVerbatim}


\subsection{Fichier irj : test.irj}
\label{\detokenize{exemples/valeurs:fichier-irj-test-irj}}
\begin{sphinxVerbatim}[commandchars=\\\{\}]
\PYG{c+c1}{\PYGZsh{}NOM}
\PYG{n}{MON}\PYG{o}{\PYGZhy{}}\PYG{n}{TEST}
\PYG{c+c1}{\PYGZsh{}ENTREES\PYGZhy{}PRIMITIF}
\PYG{c+c1}{\PYGZsh{}CONTROLES\PYGZhy{}PRIMITIF}
\PYG{c+c1}{\PYGZsh{}RESULTATS\PYGZhy{}PRIMITIF}
\PYG{c+c1}{\PYGZsh{}ENTREES\PYGZhy{}CORRECTIF}
\PYG{c+c1}{\PYGZsh{}CONTROLES\PYGZhy{}CORRECTIF}
\PYG{c+c1}{\PYGZsh{}RESULTATS\PYGZhy{}CORRECTIF}
\PYG{c+c1}{\PYGZsh{}\PYGZsh{}}
\end{sphinxVerbatim}


\subsection{Commande}
\label{\detokenize{exemples/valeurs:commande}}
\begin{sphinxVerbatim}[commandchars=\\\{\}]
\PYG{n}{mlang} \PYG{n}{test}\PYG{o}{.}\PYG{n}{m} \PYGZbs{}
	\PYG{o}{\PYGZhy{}}\PYG{o}{\PYGZhy{}}\PYG{n}{without\PYGZus{}dfgip\PYGZus{}m} \PYGZbs{}
	\PYG{o}{\PYGZhy{}}\PYG{n}{A} \PYG{n}{mon\PYGZus{}application} \PYGZbs{}
	\PYG{o}{\PYGZhy{}}\PYG{o}{\PYGZhy{}}\PYG{n}{mpp\PYGZus{}function} \PYG{n}{indefini\PYGZus{}test} \PYGZbs{}
	\PYG{o}{\PYGZhy{}}\PYG{o}{\PYGZhy{}}\PYG{n}{run\PYGZus{}test} \PYG{n}{test}\PYG{o}{.}\PYG{n}{irj}
\end{sphinxVerbatim}

\sphinxstepscope


\section{Les fonctions}
\label{\detokenize{exemples/fonctions:les-fonctions}}\label{\detokenize{exemples/fonctions:exemples-fonctions}}\label{\detokenize{exemples/fonctions::doc}}

\subsection{Fichier test : test.m}
\label{\detokenize{exemples/fonctions:fichier-test-test-m}}
\begin{sphinxVerbatim}[commandchars=\\\{\}]
\PYG{n}{espace\PYGZus{}variables} \PYG{n}{GLOBAL} \PYG{p}{:} \PYG{n}{par\PYGZus{}defaut}\PYG{p}{;}
\PYG{n}{domaine} \PYG{n}{regle} \PYG{n}{mon\PYGZus{}domaine\PYGZus{}de\PYGZus{}regle}\PYG{p}{:} \PYG{n}{par\PYGZus{}defaut}\PYG{p}{;}
\PYG{n}{domaine} \PYG{n}{verif} \PYG{n}{mon\PYGZus{}domaine\PYGZus{}de\PYGZus{}verifications}\PYG{p}{:} \PYG{n}{par\PYGZus{}defaut}\PYG{p}{;}
\PYG{n}{application} \PYG{n}{mon\PYGZus{}application}\PYG{p}{;}
\PYG{n}{variable} \PYG{n}{calculee} \PYG{p}{:} \PYG{n}{attribut} \PYG{n}{mon\PYGZus{}attribut}\PYG{p}{;}

\PYG{n}{evenement}
\PYG{p}{:} \PYG{n}{valeur} \PYG{n}{ev\PYGZus{}val}
\PYG{p}{:} \PYG{n}{variable} \PYG{n}{ev\PYGZus{}var}\PYG{p}{;}

\PYG{n}{X} \PYG{p}{:} \PYG{n}{calculee} \PYG{n}{mon\PYGZus{}attribut} \PYG{o}{=} \PYG{l+m+mi}{0} \PYG{p}{:} \PYG{l+s+s2}{\PYGZdq{}}\PYG{l+s+s2}{\PYGZdq{}} \PYG{n+nb}{type} \PYG{n}{REEL}\PYG{p}{;}
\PYG{n}{Y} \PYG{p}{:} \PYG{n}{calculee} \PYG{n}{mon\PYGZus{}attribut} \PYG{o}{=} \PYG{l+m+mi}{1} \PYG{p}{:} \PYG{l+s+s2}{\PYGZdq{}}\PYG{l+s+s2}{\PYGZdq{}}\PYG{p}{;}
\PYG{n}{TAB} \PYG{p}{:} \PYG{n}{tableau}\PYG{p}{[}\PYG{l+m+mi}{10}\PYG{p}{]} \PYG{n}{calculee} \PYG{n}{mon\PYGZus{}attribut} \PYG{o}{=} \PYG{l+m+mi}{2} \PYG{p}{:} \PYG{l+s+s2}{\PYGZdq{}}\PYG{l+s+s2}{\PYGZdq{}} \PYG{n+nb}{type} \PYG{n}{ENTIER}\PYG{p}{;}

\PYG{n}{cible} \PYG{n}{init\PYGZus{}tab}\PYG{p}{:}
\PYG{n}{application} \PYG{p}{:} \PYG{n}{mon\PYGZus{}application}\PYG{p}{;}
\PYG{n}{iterer} \PYG{p}{:} \PYG{n}{variable} \PYG{n}{I} \PYG{p}{:} \PYG{n}{entre} \PYG{l+m+mf}{0.}\PYG{o}{.}\PYG{p}{(}\PYG{n}{taille}\PYG{p}{(}\PYG{n}{TAB}\PYG{p}{)} \PYG{o}{\PYGZhy{}} \PYG{l+m+mi}{1}\PYG{p}{)} \PYG{n}{increment} \PYG{l+m+mi}{1} \PYG{p}{:} \PYG{n}{dans} \PYG{p}{(}
  \PYG{n}{TAB}\PYG{p}{[}\PYG{n}{I}\PYG{p}{]} \PYG{o}{=} \PYG{n}{I}\PYG{p}{;}
\PYG{p}{)}

\PYG{n}{cible} \PYG{n}{reinit\PYGZus{}tab}\PYG{p}{:}
\PYG{n}{application} \PYG{p}{:} \PYG{n}{mon\PYGZus{}application}\PYG{p}{;}
\PYG{n}{iterer} \PYG{p}{:} \PYG{n}{variable} \PYG{n}{I} \PYG{p}{:} \PYG{n}{entre} \PYG{l+m+mf}{0.}\PYG{o}{.}\PYG{p}{(}\PYG{n}{taille}\PYG{p}{(}\PYG{n}{TAB}\PYG{p}{)} \PYG{o}{\PYGZhy{}} \PYG{l+m+mi}{1}\PYG{p}{)} \PYG{n}{increment} \PYG{l+m+mi}{1} \PYG{p}{:} \PYG{n}{dans} \PYG{p}{(}
  \PYG{n}{TAB}\PYG{p}{[}\PYG{n}{I}\PYG{p}{]} \PYG{o}{=} \PYG{n}{indefini}\PYG{p}{;}
\PYG{p}{)}

\PYG{n}{cible} \PYG{n}{test\PYGZus{}abs}\PYG{p}{:}
\PYG{n}{application} \PYG{p}{:} \PYG{n}{mon\PYGZus{}application}\PYG{p}{;}
\PYG{n}{afficher} \PYG{l+s+s2}{\PYGZdq{}}\PYG{l+s+se}{\PYGZbs{}n}\PYG{l+s+s2}{\PYGZus{}\PYGZus{}ABS\PYGZus{}\PYGZus{}}\PYG{l+s+s2}{\PYGZdq{}}\PYG{p}{;}
\PYG{n}{afficher} \PYG{n}{indenter}\PYG{p}{(}\PYG{l+m+mi}{2}\PYG{p}{)}\PYG{p}{;}
\PYG{n}{afficher} \PYG{l+s+s2}{\PYGZdq{}}\PYG{l+s+se}{\PYGZbs{}n}\PYG{l+s+s2}{abs(indefini) = }\PYG{l+s+s2}{\PYGZdq{}}\PYG{p}{;}
\PYG{n}{afficher} \PYG{p}{(}\PYG{n+nb}{abs}\PYG{p}{(}\PYG{n}{indefini}\PYG{p}{)}\PYG{p}{)}\PYG{p}{;}
\PYG{n}{afficher} \PYG{l+s+s2}{\PYGZdq{}}\PYG{l+s+se}{\PYGZbs{}n}\PYG{l+s+s2}{abs(1) = }\PYG{l+s+s2}{\PYGZdq{}}\PYG{p}{;}
\PYG{n}{afficher}  \PYG{p}{(}\PYG{n+nb}{abs}\PYG{p}{(}\PYG{l+m+mi}{1}\PYG{p}{)}\PYG{p}{)}\PYG{p}{;}
\PYG{n}{afficher} \PYG{l+s+s2}{\PYGZdq{}}\PYG{l+s+se}{\PYGZbs{}n}\PYG{l+s+s2}{abs(\PYGZhy{}1) = }\PYG{l+s+s2}{\PYGZdq{}}\PYG{p}{;}
\PYG{n}{afficher}  \PYG{p}{(}\PYG{n+nb}{abs}\PYG{p}{(}\PYG{o}{\PYGZhy{}}\PYG{l+m+mi}{1}\PYG{p}{)}\PYG{p}{)}\PYG{p}{;}
\PYG{n}{afficher} \PYG{l+s+s2}{\PYGZdq{}}\PYG{l+s+se}{\PYGZbs{}n}\PYG{l+s+s2}{\PYGZdq{}}\PYG{p}{;}
\PYG{n}{afficher} \PYG{n}{indenter}\PYG{p}{(}\PYG{o}{\PYGZhy{}}\PYG{l+m+mi}{2}\PYG{p}{)}\PYG{p}{;}

\PYG{n}{cible} \PYG{n}{test\PYGZus{}arr}\PYG{p}{:}
\PYG{n}{application} \PYG{p}{:} \PYG{n}{mon\PYGZus{}application}\PYG{p}{;}
\PYG{n}{afficher} \PYG{l+s+s2}{\PYGZdq{}}\PYG{l+s+se}{\PYGZbs{}n}\PYG{l+s+s2}{\PYGZus{}\PYGZus{}ARR\PYGZus{}\PYGZus{}}\PYG{l+s+s2}{\PYGZdq{}}\PYG{p}{;}
\PYG{n}{afficher} \PYG{n}{indenter}\PYG{p}{(}\PYG{l+m+mi}{2}\PYG{p}{)}\PYG{p}{;}
\PYG{n}{afficher} \PYG{l+s+s2}{\PYGZdq{}}\PYG{l+s+se}{\PYGZbs{}n}\PYG{l+s+s2}{arr(indefini) = }\PYG{l+s+s2}{\PYGZdq{}}\PYG{p}{;}
\PYG{n}{afficher}  \PYG{p}{(}\PYG{n}{arr}\PYG{p}{(}\PYG{n}{indefini}\PYG{p}{)}\PYG{p}{)}\PYG{p}{;}
\PYG{n}{afficher} \PYG{l+s+s2}{\PYGZdq{}}\PYG{l+s+se}{\PYGZbs{}n}\PYG{l+s+s2}{arr(1.8) = }\PYG{l+s+s2}{\PYGZdq{}}\PYG{p}{;}
\PYG{n}{afficher}  \PYG{p}{(}\PYG{n}{arr}\PYG{p}{(}\PYG{l+m+mf}{1.8}\PYG{p}{)}\PYG{p}{)}\PYG{p}{;}
\PYG{n}{afficher} \PYG{l+s+s2}{\PYGZdq{}}\PYG{l+s+s2}{\PYGZbs{}}\PYG{l+s+s2}{ arr(\PYGZhy{}1.7) = }\PYG{l+s+s2}{\PYGZdq{}}\PYG{p}{;}
\PYG{n}{afficher}  \PYG{p}{(}\PYG{n}{arr}\PYG{p}{(}\PYG{o}{\PYGZhy{}}\PYG{l+m+mf}{1.7}\PYG{p}{)}\PYG{p}{)}\PYG{p}{;}
\PYG{n}{afficher} \PYG{l+s+s2}{\PYGZdq{}}\PYG{l+s+se}{\PYGZbs{}n}\PYG{l+s+s2}{\PYGZdq{}}\PYG{p}{;}
\PYG{n}{afficher} \PYG{n}{indenter}\PYG{p}{(}\PYG{o}{\PYGZhy{}}\PYG{l+m+mi}{2}\PYG{p}{)}\PYG{p}{;}

\PYG{n}{cible} \PYG{n}{test\PYGZus{}attribut}\PYG{p}{:}
\PYG{n}{application} \PYG{p}{:} \PYG{n}{mon\PYGZus{}application}\PYG{p}{;}
\PYG{n}{afficher} \PYG{l+s+s2}{\PYGZdq{}}\PYG{l+s+se}{\PYGZbs{}n}\PYG{l+s+s2}{\PYGZus{}\PYGZus{}ATTRIBUT\PYGZus{}\PYGZus{}}\PYG{l+s+s2}{\PYGZdq{}}\PYG{p}{;}
\PYG{n}{afficher} \PYG{n}{indenter}\PYG{p}{(}\PYG{l+m+mi}{2}\PYG{p}{)}\PYG{p}{;}
\PYG{n}{afficher} \PYG{l+s+s2}{\PYGZdq{}}\PYG{l+s+se}{\PYGZbs{}n}\PYG{l+s+s2}{attribut(X, mon\PYGZus{}attribut) = }\PYG{l+s+s2}{\PYGZdq{}}\PYG{p}{;}
\PYG{n}{afficher}  \PYG{p}{(}\PYG{n}{attribut}\PYG{p}{(}\PYG{n}{X}\PYG{p}{,} \PYG{n}{mon\PYGZus{}attribut}\PYG{p}{)}\PYG{p}{)}\PYG{p}{;}
\PYG{n}{afficher} \PYG{l+s+s2}{\PYGZdq{}}\PYG{l+s+se}{\PYGZbs{}n}\PYG{l+s+s2}{attribut(TAB, mon\PYGZus{}attribut) = }\PYG{l+s+s2}{\PYGZdq{}}\PYG{p}{;}
\PYG{n}{afficher}  \PYG{p}{(}\PYG{n}{attribut}\PYG{p}{(}\PYG{n}{TAB}\PYG{p}{,} \PYG{n}{mon\PYGZus{}attribut}\PYG{p}{)}\PYG{p}{)}\PYG{p}{;}
\PYG{n}{afficher} \PYG{l+s+s2}{\PYGZdq{}}\PYG{l+s+se}{\PYGZbs{}n}\PYG{l+s+s2}{\PYGZdq{}}\PYG{p}{;}
\PYG{n}{afficher} \PYG{n}{indenter}\PYG{p}{(}\PYG{o}{\PYGZhy{}}\PYG{l+m+mi}{2}\PYG{p}{)}\PYG{p}{;}

\PYG{n}{cible} \PYG{n}{test\PYGZus{}champ\PYGZus{}evenement\PYGZus{}base}\PYG{p}{:}
\PYG{n}{application} \PYG{p}{:} \PYG{n}{mon\PYGZus{}application}\PYG{p}{;}
\PYG{n}{afficher} \PYG{n}{indenter}\PYG{p}{(}\PYG{l+m+mi}{2}\PYG{p}{)}\PYG{p}{;}
\PYG{n}{afficher} \PYG{l+s+s2}{\PYGZdq{}}\PYG{l+s+se}{\PYGZbs{}n}\PYG{l+s+s2}{champ\PYGZus{}evenement(0, ev\PYGZus{}val) = }\PYG{l+s+s2}{\PYGZdq{}}\PYG{p}{;}
\PYG{n}{afficher} \PYG{p}{(}\PYG{n}{champ\PYGZus{}evenement} \PYG{p}{(}\PYG{l+m+mi}{0}\PYG{p}{,}\PYG{n}{ev\PYGZus{}val}\PYG{p}{)}\PYG{p}{)}\PYG{p}{;}
\PYG{n}{afficher} \PYG{l+s+s2}{\PYGZdq{}}\PYG{l+s+se}{\PYGZbs{}n}\PYG{l+s+s2}{champ\PYGZus{}evenement(0, ev\PYGZus{}var) = }\PYG{l+s+s2}{\PYGZdq{}}\PYG{p}{;}
\PYG{n}{afficher} \PYG{p}{(}\PYG{n}{champ\PYGZus{}evenement} \PYG{p}{(}\PYG{l+m+mi}{0}\PYG{p}{,}\PYG{n}{ev\PYGZus{}var}\PYG{p}{)}\PYG{p}{)}\PYG{p}{;}
\PYG{n}{afficher} \PYG{l+s+s2}{\PYGZdq{}}\PYG{l+s+se}{\PYGZbs{}n}\PYG{l+s+s2}{\PYGZdq{}}\PYG{p}{;}
\PYG{n}{afficher} \PYG{n}{indenter}\PYG{p}{(}\PYG{o}{\PYGZhy{}}\PYG{l+m+mi}{2}\PYG{p}{)}\PYG{p}{;}

\PYG{n}{cible} \PYG{n}{test\PYGZus{}champ\PYGZus{}evenement}\PYG{p}{:}
\PYG{n}{application} \PYG{p}{:} \PYG{n}{mon\PYGZus{}application}\PYG{p}{;}
\PYG{n}{afficher} \PYG{l+s+s2}{\PYGZdq{}}\PYG{l+s+se}{\PYGZbs{}n}\PYG{l+s+s2}{\PYGZus{}\PYGZus{}CHAMP\PYGZus{}EVENEMENT\PYGZus{}\PYGZus{}}\PYG{l+s+s2}{\PYGZdq{}}\PYG{p}{;}
\PYG{n}{afficher} \PYG{n}{indenter}\PYG{p}{(}\PYG{l+m+mi}{2}\PYG{p}{)}\PYG{p}{;}
\PYG{n}{arranger\PYGZus{}evenements}
  \PYG{p}{:} \PYG{n}{ajouter} \PYG{l+m+mi}{1}
  \PYG{p}{:} \PYG{n}{dans} \PYG{p}{(}
      \PYG{n}{afficher} \PYG{l+s+s2}{\PYGZdq{}}\PYG{l+s+se}{\PYGZbs{}n}\PYG{l+s+s2}{Avant d}\PYG{l+s+s2}{\PYGZsq{}}\PYG{l+s+s2}{initialiser les champs de l}\PYG{l+s+s2}{\PYGZsq{}}\PYG{l+s+s2}{événement:}\PYG{l+s+s2}{\PYGZdq{}}\PYG{p}{;}
      \PYG{n}{calculer} \PYG{n}{cible} \PYG{n}{test\PYGZus{}champ\PYGZus{}evenement\PYGZus{}base}\PYG{p}{;}
      \PYG{n}{champ\PYGZus{}evenement} \PYG{p}{(}\PYG{l+m+mi}{0}\PYG{p}{,}\PYG{n}{ev\PYGZus{}var}\PYG{p}{)} \PYG{n}{reference} \PYG{n}{X}\PYG{p}{;}      
      \PYG{n}{champ\PYGZus{}evenement} \PYG{p}{(}\PYG{l+m+mi}{0}\PYG{p}{,}\PYG{n}{ev\PYGZus{}val}\PYG{p}{)} \PYG{o}{=} \PYG{l+m+mi}{42}\PYG{p}{;}
      \PYG{n}{X} \PYG{o}{=} \PYG{l+m+mi}{2}\PYG{p}{;}
      \PYG{n}{afficher} \PYG{l+s+s2}{\PYGZdq{}}\PYG{l+s+se}{\PYGZbs{}n}\PYG{l+s+s2}{Après avoir initialisé les champs de l}\PYG{l+s+s2}{\PYGZsq{}}\PYG{l+s+s2}{événement:}\PYG{l+s+s2}{\PYGZdq{}}\PYG{p}{;}
      \PYG{n}{calculer} \PYG{n}{cible} \PYG{n}{test\PYGZus{}champ\PYGZus{}evenement\PYGZus{}base}\PYG{p}{;}
      \PYG{n}{X} \PYG{o}{=} \PYG{n}{indefini}\PYG{p}{;}
    \PYG{p}{)}
\PYG{n}{afficher} \PYG{n}{indenter}\PYG{p}{(}\PYG{o}{\PYGZhy{}}\PYG{l+m+mi}{2}\PYG{p}{)}\PYG{p}{;}

\PYG{n}{cible} \PYG{n}{test\PYGZus{}inf}\PYG{p}{:}
\PYG{n}{application} \PYG{p}{:} \PYG{n}{mon\PYGZus{}application}\PYG{p}{;}
\PYG{n}{afficher} \PYG{l+s+s2}{\PYGZdq{}}\PYG{l+s+se}{\PYGZbs{}n}\PYG{l+s+s2}{\PYGZus{}\PYGZus{}INF\PYGZus{}\PYGZus{}}\PYG{l+s+s2}{\PYGZdq{}}\PYG{p}{;}
\PYG{n}{afficher} \PYG{n}{indenter}\PYG{p}{(}\PYG{l+m+mi}{2}\PYG{p}{)}\PYG{p}{;}
\PYG{n}{afficher} \PYG{l+s+s2}{\PYGZdq{}}\PYG{l+s+se}{\PYGZbs{}n}\PYG{l+s+s2}{inf(indefini) = }\PYG{l+s+s2}{\PYGZdq{}}\PYG{p}{;}
\PYG{n}{afficher}  \PYG{p}{(}\PYG{n}{inf}\PYG{p}{(}\PYG{n}{indefini}\PYG{p}{)}\PYG{p}{)}\PYG{p}{;}
\PYG{n}{afficher} \PYG{l+s+s2}{\PYGZdq{}}\PYG{l+s+se}{\PYGZbs{}n}\PYG{l+s+s2}{inf(1.8) = }\PYG{l+s+s2}{\PYGZdq{}}\PYG{p}{;}
\PYG{n}{afficher}  \PYG{p}{(}\PYG{n}{inf}\PYG{p}{(}\PYG{l+m+mf}{1.8}\PYG{p}{)}\PYG{p}{)}\PYG{p}{;}
\PYG{n}{afficher} \PYG{l+s+s2}{\PYGZdq{}}\PYG{l+s+se}{\PYGZbs{}n}\PYG{l+s+s2}{inf(\PYGZhy{}1.7) = }\PYG{l+s+s2}{\PYGZdq{}}\PYG{p}{;}
\PYG{n}{afficher}  \PYG{p}{(}\PYG{n}{inf}\PYG{p}{(}\PYG{o}{\PYGZhy{}}\PYG{l+m+mf}{1.7}\PYG{p}{)}\PYG{p}{)}\PYG{p}{;}
\PYG{n}{afficher} \PYG{l+s+s2}{\PYGZdq{}}\PYG{l+s+se}{\PYGZbs{}n}\PYG{l+s+s2}{\PYGZdq{}}\PYG{p}{;}
\PYG{n}{afficher} \PYG{n}{indenter}\PYG{p}{(}\PYG{o}{\PYGZhy{}}\PYG{l+m+mi}{2}\PYG{p}{)}\PYG{p}{;}

\PYG{n}{cible} \PYG{n}{test\PYGZus{}max}\PYG{p}{:}
\PYG{n}{application} \PYG{p}{:} \PYG{n}{mon\PYGZus{}application}\PYG{p}{;}
\PYG{n}{afficher} \PYG{l+s+s2}{\PYGZdq{}}\PYG{l+s+se}{\PYGZbs{}n}\PYG{l+s+s2}{\PYGZus{}\PYGZus{}MAX\PYGZus{}\PYGZus{}}\PYG{l+s+s2}{\PYGZdq{}}\PYG{p}{;}
\PYG{n}{afficher} \PYG{n}{indenter}\PYG{p}{(}\PYG{l+m+mi}{2}\PYG{p}{)}\PYG{p}{;}
\PYG{n}{afficher} \PYG{l+s+s2}{\PYGZdq{}}\PYG{l+s+se}{\PYGZbs{}n}\PYG{l+s+s2}{max(indefini, indefini) = }\PYG{l+s+s2}{\PYGZdq{}}\PYG{p}{;}
\PYG{n}{afficher}   \PYG{p}{(}\PYG{n+nb}{max}\PYG{p}{(}\PYG{n}{indefini}\PYG{p}{,} \PYG{n}{indefini}\PYG{p}{)}\PYG{p}{)}\PYG{p}{;}
\PYG{n}{afficher} \PYG{l+s+s2}{\PYGZdq{}}\PYG{l+s+se}{\PYGZbs{}n}\PYG{l+s+s2}{max(\PYGZhy{}1, indefini) = }\PYG{l+s+s2}{\PYGZdq{}}\PYG{p}{;}
\PYG{n}{afficher}   \PYG{p}{(}\PYG{n+nb}{max}\PYG{p}{(}\PYG{o}{\PYGZhy{}}\PYG{l+m+mi}{1}\PYG{p}{,} \PYG{n}{indefini}\PYG{p}{)}\PYG{p}{)}\PYG{p}{;}
\PYG{n}{afficher} \PYG{l+s+s2}{\PYGZdq{}}\PYG{l+s+se}{\PYGZbs{}n}\PYG{l+s+s2}{max(indefini, \PYGZhy{}1) = }\PYG{l+s+s2}{\PYGZdq{}}\PYG{p}{;}
\PYG{n}{afficher}   \PYG{p}{(}\PYG{n+nb}{max}\PYG{p}{(}\PYG{n}{indefini}\PYG{p}{,} \PYG{o}{\PYGZhy{}}\PYG{l+m+mi}{1}\PYG{p}{)}\PYG{p}{)}\PYG{p}{;}
\PYG{n}{afficher} \PYG{l+s+s2}{\PYGZdq{}}\PYG{l+s+se}{\PYGZbs{}n}\PYG{l+s+s2}{max(1, indefini) = }\PYG{l+s+s2}{\PYGZdq{}}\PYG{p}{;}
\PYG{n}{afficher}   \PYG{p}{(}\PYG{n+nb}{max}\PYG{p}{(}\PYG{l+m+mi}{1}\PYG{p}{,} \PYG{n}{indefini}\PYG{p}{)}\PYG{p}{)}\PYG{p}{;}
\PYG{n}{afficher} \PYG{l+s+s2}{\PYGZdq{}}\PYG{l+s+se}{\PYGZbs{}n}\PYG{l+s+s2}{\PYGZdq{}}\PYG{p}{;}
\PYG{n}{afficher} \PYG{n}{indenter}\PYG{p}{(}\PYG{o}{\PYGZhy{}}\PYG{l+m+mi}{2}\PYG{p}{)}\PYG{p}{;}

\PYG{n}{cible} \PYG{n}{test\PYGZus{}meme\PYGZus{}variable}\PYG{p}{:}
\PYG{n}{application} \PYG{p}{:} \PYG{n}{mon\PYGZus{}application}\PYG{p}{;}
\PYG{n}{afficher} \PYG{l+s+s2}{\PYGZdq{}}\PYG{l+s+se}{\PYGZbs{}n}\PYG{l+s+s2}{\PYGZus{}\PYGZus{}MEME\PYGZus{}VARIABLE\PYGZus{}\PYGZus{}}\PYG{l+s+s2}{\PYGZdq{}}\PYG{p}{;}
\PYG{n}{afficher} \PYG{n}{indenter}\PYG{p}{(}\PYG{l+m+mi}{2}\PYG{p}{)}\PYG{p}{;}
\PYG{n}{afficher} \PYG{l+s+s2}{\PYGZdq{}}\PYG{l+s+se}{\PYGZbs{}n}\PYG{l+s+s2}{meme\PYGZus{}variable(X,X) = }\PYG{l+s+s2}{\PYGZdq{}}\PYG{p}{;}
\PYG{n}{afficher} \PYG{p}{(}\PYG{n}{meme\PYGZus{}variable}\PYG{p}{(}\PYG{n}{X}\PYG{p}{,}\PYG{n}{X}\PYG{p}{)}\PYG{p}{)}\PYG{p}{;}
\PYG{n}{afficher} \PYG{l+s+s2}{\PYGZdq{}}\PYG{l+s+se}{\PYGZbs{}n}\PYG{l+s+s2}{meme\PYGZus{}variable(X,TAB) = }\PYG{l+s+s2}{\PYGZdq{}}\PYG{p}{;}
\PYG{n}{afficher} \PYG{p}{(}\PYG{n}{meme\PYGZus{}variable}\PYG{p}{(}\PYG{n}{X}\PYG{p}{,}\PYG{n}{TAB}\PYG{p}{)}\PYG{p}{)}\PYG{p}{;}
\PYG{n}{afficher} \PYG{l+s+s2}{\PYGZdq{}}\PYG{l+s+se}{\PYGZbs{}n}\PYG{l+s+s2}{meme\PYGZus{}variable(TAB[0],TAB) = }\PYG{l+s+s2}{\PYGZdq{}}\PYG{p}{;}
\PYG{n}{afficher} \PYG{p}{(}\PYG{n}{meme\PYGZus{}variable}\PYG{p}{(}\PYG{n}{TAB}\PYG{p}{[}\PYG{l+m+mi}{0}\PYG{p}{]}\PYG{p}{,}\PYG{n}{TAB}\PYG{p}{)}\PYG{p}{)}\PYG{p}{;}
\PYG{n}{afficher} \PYG{l+s+s2}{\PYGZdq{}}\PYG{l+s+se}{\PYGZbs{}n}\PYG{l+s+s2}{\PYGZdq{}}\PYG{p}{;}
\PYG{n}{afficher} \PYG{n}{indenter}\PYG{p}{(}\PYG{o}{\PYGZhy{}}\PYG{l+m+mi}{2}\PYG{p}{)}\PYG{p}{;}

\PYG{n}{cible} \PYG{n}{test\PYGZus{}min}\PYG{p}{:}
\PYG{n}{application} \PYG{p}{:} \PYG{n}{mon\PYGZus{}application}\PYG{p}{;}
\PYG{n}{afficher} \PYG{l+s+s2}{\PYGZdq{}}\PYG{l+s+se}{\PYGZbs{}n}\PYG{l+s+s2}{\PYGZus{}\PYGZus{}MIN\PYGZus{}\PYGZus{}}\PYG{l+s+s2}{\PYGZdq{}}\PYG{p}{;}
\PYG{n}{afficher} \PYG{n}{indenter}\PYG{p}{(}\PYG{l+m+mi}{2}\PYG{p}{)}\PYG{p}{;}
\PYG{n}{afficher} \PYG{l+s+s2}{\PYGZdq{}}\PYG{l+s+se}{\PYGZbs{}n}\PYG{l+s+s2}{min(indefini, indefini) = }\PYG{l+s+s2}{\PYGZdq{}}\PYG{p}{;}
\PYG{n}{afficher} \PYG{p}{(}\PYG{n+nb}{min}\PYG{p}{(}\PYG{n}{indefini}\PYG{p}{,} \PYG{n}{indefini}\PYG{p}{)}\PYG{p}{)}\PYG{p}{;}
\PYG{n}{afficher} \PYG{l+s+s2}{\PYGZdq{}}\PYG{l+s+se}{\PYGZbs{}n}\PYG{l+s+s2}{min(1, indefini) = }\PYG{l+s+s2}{\PYGZdq{}}\PYG{p}{;}
\PYG{n}{afficher} \PYG{p}{(}\PYG{n+nb}{min}\PYG{p}{(}\PYG{l+m+mi}{1}\PYG{p}{,} \PYG{n}{indefini}\PYG{p}{)}\PYG{p}{)}\PYG{p}{;}
\PYG{n}{afficher} \PYG{l+s+s2}{\PYGZdq{}}\PYG{l+s+se}{\PYGZbs{}n}\PYG{l+s+s2}{min(indefini, 1) = }\PYG{l+s+s2}{\PYGZdq{}}\PYG{p}{;}
\PYG{n}{afficher} \PYG{p}{(}\PYG{n+nb}{min}\PYG{p}{(}\PYG{n}{indefini}\PYG{p}{,} \PYG{l+m+mi}{1}\PYG{p}{)}\PYG{p}{)}\PYG{p}{;}
\PYG{n}{afficher} \PYG{l+s+s2}{\PYGZdq{}}\PYG{l+s+se}{\PYGZbs{}n}\PYG{l+s+s2}{min(\PYGZhy{}1, indefini) = }\PYG{l+s+s2}{\PYGZdq{}}\PYG{p}{;}
\PYG{n}{afficher} \PYG{p}{(}\PYG{n+nb}{min}\PYG{p}{(}\PYG{o}{\PYGZhy{}}\PYG{l+m+mi}{1}\PYG{p}{,} \PYG{n}{indefini}\PYG{p}{)}\PYG{p}{)}\PYG{p}{;}
\PYG{n}{afficher} \PYG{l+s+s2}{\PYGZdq{}}\PYG{l+s+se}{\PYGZbs{}n}\PYG{l+s+s2}{\PYGZdq{}}\PYG{p}{;}
\PYG{n}{afficher} \PYG{n}{indenter}\PYG{p}{(}\PYG{o}{\PYGZhy{}}\PYG{l+m+mi}{2}\PYG{p}{)}\PYG{p}{;}

\PYG{n}{cible} \PYG{n}{test\PYGZus{}multimax\PYGZus{}base}\PYG{p}{:}
\PYG{n}{application} \PYG{p}{:} \PYG{n}{mon\PYGZus{}application}\PYG{p}{;}
\PYG{n}{afficher} \PYG{n}{indenter}\PYG{p}{(}\PYG{l+m+mi}{2}\PYG{p}{)}\PYG{p}{;}
\PYG{n}{afficher} \PYG{l+s+s2}{\PYGZdq{}}\PYG{l+s+se}{\PYGZbs{}n}\PYG{l+s+s2}{multimax(indefini, TAB) = }\PYG{l+s+s2}{\PYGZdq{}}\PYG{p}{;}
\PYG{n}{afficher} \PYG{p}{(}\PYG{n}{multimax}\PYG{p}{(}\PYG{n}{indefini}\PYG{p}{,} \PYG{n}{TAB}\PYG{p}{)}\PYG{p}{)}\PYG{p}{;}
\PYG{n}{afficher} \PYG{l+s+s2}{\PYGZdq{}}\PYG{l+s+se}{\PYGZbs{}n}\PYG{l+s+s2}{multimax(7, TAB) = }\PYG{l+s+s2}{\PYGZdq{}}\PYG{p}{;}
\PYG{n}{afficher} \PYG{p}{(}\PYG{n}{multimax}\PYG{p}{(}\PYG{l+m+mi}{7}\PYG{p}{,} \PYG{n}{TAB}\PYG{p}{)}\PYG{p}{)}\PYG{p}{;}
\PYG{n}{afficher} \PYG{l+s+s2}{\PYGZdq{}}\PYG{l+s+se}{\PYGZbs{}n}\PYG{l+s+s2}{multimax(taille(TAB) + 1, TAB) = }\PYG{l+s+s2}{\PYGZdq{}}\PYG{p}{;}  
\PYG{n}{afficher} \PYG{p}{(}\PYG{n}{multimax}\PYG{p}{(}\PYG{n}{taille}\PYG{p}{(}\PYG{n}{TAB}\PYG{p}{)} \PYG{o}{+} \PYG{l+m+mi}{1}\PYG{p}{,} \PYG{n}{TAB}\PYG{p}{)}\PYG{p}{)}\PYG{p}{;}
\PYG{n}{afficher} \PYG{l+s+s2}{\PYGZdq{}}\PYG{l+s+se}{\PYGZbs{}n}\PYG{l+s+s2}{multimax(0, TAB) = }\PYG{l+s+s2}{\PYGZdq{}}\PYG{p}{;}  
\PYG{n}{afficher} \PYG{p}{(}\PYG{n}{multimax}\PYG{p}{(}\PYG{l+m+mi}{0}\PYG{p}{,} \PYG{n}{TAB}\PYG{p}{)}\PYG{p}{)}\PYG{p}{;}
\PYG{n}{afficher} \PYG{l+s+s2}{\PYGZdq{}}\PYG{l+s+se}{\PYGZbs{}n}\PYG{l+s+s2}{multimax(\PYGZhy{}1, TAB) = }\PYG{l+s+s2}{\PYGZdq{}}\PYG{p}{;}  
\PYG{n}{afficher} \PYG{p}{(}\PYG{n}{multimax}\PYG{p}{(}\PYG{o}{\PYGZhy{}}\PYG{l+m+mi}{1}\PYG{p}{,} \PYG{n}{TAB}\PYG{p}{)}\PYG{p}{)}\PYG{p}{;}
\PYG{n}{afficher} \PYG{n}{indenter}\PYG{p}{(}\PYG{o}{\PYGZhy{}}\PYG{l+m+mi}{2}\PYG{p}{)}\PYG{p}{;}

\PYG{n}{cible} \PYG{n}{test\PYGZus{}multimax}\PYG{p}{:}
\PYG{n}{application} \PYG{p}{:} \PYG{n}{mon\PYGZus{}application}\PYG{p}{;}
\PYG{n}{afficher} \PYG{l+s+s2}{\PYGZdq{}}\PYG{l+s+se}{\PYGZbs{}n}\PYG{l+s+s2}{\PYGZus{}\PYGZus{}MULTIMAX\PYGZus{}\PYGZus{}}\PYG{l+s+s2}{\PYGZdq{}}\PYG{p}{;}
\PYG{n}{afficher} \PYG{n}{indenter}\PYG{p}{(}\PYG{l+m+mi}{2}\PYG{p}{)}\PYG{p}{;}
\PYG{n}{afficher} \PYG{l+s+s2}{\PYGZdq{}}\PYG{l+s+se}{\PYGZbs{}n}\PYG{l+s+s2}{Avant initialisation du tableau :}\PYG{l+s+s2}{\PYGZdq{}}\PYG{p}{;}
\PYG{n}{calculer} \PYG{n}{cible} \PYG{n}{test\PYGZus{}multimax\PYGZus{}base}\PYG{p}{;}
\PYG{n}{calculer} \PYG{n}{cible} \PYG{n}{init\PYGZus{}tab}\PYG{p}{;}
\PYG{n}{afficher} \PYG{l+s+s2}{\PYGZdq{}}\PYG{l+s+se}{\PYGZbs{}n}\PYG{l+s+s2}{Après initialisation du tableau :}\PYG{l+s+s2}{\PYGZdq{}}\PYG{p}{;}
\PYG{n}{calculer} \PYG{n}{cible} \PYG{n}{test\PYGZus{}multimax\PYGZus{}base}\PYG{p}{;}
\PYG{n}{calculer} \PYG{n}{cible} \PYG{n}{reinit\PYGZus{}tab}\PYG{p}{;}
\PYG{n}{afficher} \PYG{l+s+s2}{\PYGZdq{}}\PYG{l+s+se}{\PYGZbs{}n}\PYG{l+s+s2}{\PYGZdq{}}\PYG{p}{;}
\PYG{n}{afficher} \PYG{n}{indenter}\PYG{p}{(}\PYG{o}{\PYGZhy{}}\PYG{l+m+mi}{2}\PYG{p}{)}\PYG{p}{;}

\PYG{n}{cible} \PYG{n}{test\PYGZus{}nb\PYGZus{}evenements\PYGZus{}base}\PYG{p}{:}
\PYG{n}{application} \PYG{p}{:} \PYG{n}{mon\PYGZus{}application}\PYG{p}{;}
\PYG{n}{afficher} \PYG{n}{indenter}\PYG{p}{(}\PYG{l+m+mi}{2}\PYG{p}{)}\PYG{p}{;}
\PYG{n}{afficher} \PYG{l+s+s2}{\PYGZdq{}}\PYG{l+s+se}{\PYGZbs{}n}\PYG{l+s+s2}{nb\PYGZus{}evenements() = }\PYG{l+s+s2}{\PYGZdq{}}\PYG{p}{;}
\PYG{n}{afficher} \PYG{p}{(}\PYG{n}{nb\PYGZus{}evenements} \PYG{p}{(}\PYG{p}{)}\PYG{p}{)}\PYG{p}{;}
\PYG{n}{afficher} \PYG{l+s+s2}{\PYGZdq{}}\PYG{l+s+se}{\PYGZbs{}n}\PYG{l+s+s2}{\PYGZdq{}}\PYG{p}{;}
\PYG{n}{afficher} \PYG{n}{indenter}\PYG{p}{(}\PYG{o}{\PYGZhy{}}\PYG{l+m+mi}{2}\PYG{p}{)}\PYG{p}{;}

\PYG{n}{cible} \PYG{n}{test\PYGZus{}nb\PYGZus{}evenements}\PYG{p}{:}
\PYG{n}{application} \PYG{p}{:} \PYG{n}{mon\PYGZus{}application}\PYG{p}{;}
\PYG{n}{afficher} \PYG{l+s+s2}{\PYGZdq{}}\PYG{l+s+se}{\PYGZbs{}n}\PYG{l+s+s2}{\PYGZus{}\PYGZus{}NB\PYGZus{}EVENEMENTS\PYGZus{}\PYGZus{}}\PYG{l+s+s2}{\PYGZdq{}}\PYG{p}{;}
\PYG{n}{afficher} \PYG{n}{indenter}\PYG{p}{(}\PYG{l+m+mi}{2}\PYG{p}{)}\PYG{p}{;}
\PYG{n}{afficher} \PYG{l+s+s2}{\PYGZdq{}}\PYG{l+s+se}{\PYGZbs{}n}\PYG{l+s+s2}{Avant la définition d}\PYG{l+s+s2}{\PYGZsq{}}\PYG{l+s+s2}{un événement :}\PYG{l+s+s2}{\PYGZdq{}}\PYG{p}{;}
\PYG{n}{calculer} \PYG{n}{cible} \PYG{n}{test\PYGZus{}nb\PYGZus{}evenements\PYGZus{}base}\PYG{p}{;}
\PYG{n}{arranger\PYGZus{}evenements}
  \PYG{p}{:} \PYG{n}{ajouter} \PYG{l+m+mi}{1}
  \PYG{p}{:} \PYG{n}{dans} \PYG{p}{(}
      \PYG{n}{afficher} \PYG{l+s+s2}{\PYGZdq{}}\PYG{l+s+se}{\PYGZbs{}n}\PYG{l+s+s2}{Pendant la définition d}\PYG{l+s+s2}{\PYGZsq{}}\PYG{l+s+s2}{un événement :}\PYG{l+s+s2}{\PYGZdq{}}\PYG{p}{;}
      \PYG{n}{calculer} \PYG{n}{cible} \PYG{n}{test\PYGZus{}nb\PYGZus{}evenements\PYGZus{}base}\PYG{p}{;}
    \PYG{p}{)}
\PYG{n}{afficher} \PYG{l+s+s2}{\PYGZdq{}}\PYG{l+s+se}{\PYGZbs{}n}\PYG{l+s+s2}{Après la définition d}\PYG{l+s+s2}{\PYGZsq{}}\PYG{l+s+s2}{un événement :}\PYG{l+s+s2}{\PYGZdq{}}\PYG{p}{;}
\PYG{n}{calculer} \PYG{n}{cible} \PYG{n}{test\PYGZus{}nb\PYGZus{}evenements\PYGZus{}base}\PYG{p}{;}
\PYG{n}{afficher} \PYG{n}{indenter}\PYG{p}{(}\PYG{o}{\PYGZhy{}}\PYG{l+m+mi}{2}\PYG{p}{)}\PYG{p}{;}

\PYG{n}{cible} \PYG{n}{test\PYGZus{}null}\PYG{p}{:}
\PYG{n}{application} \PYG{p}{:} \PYG{n}{mon\PYGZus{}application}\PYG{p}{;}
\PYG{n}{afficher} \PYG{l+s+s2}{\PYGZdq{}}\PYG{l+s+se}{\PYGZbs{}n}\PYG{l+s+s2}{\PYGZus{}\PYGZus{}NULL\PYGZus{}\PYGZus{}}\PYG{l+s+s2}{\PYGZdq{}}\PYG{p}{;}
\PYG{n}{afficher} \PYG{n}{indenter}\PYG{p}{(}\PYG{l+m+mi}{2}\PYG{p}{)}\PYG{p}{;}
\PYG{n}{afficher} \PYG{l+s+s2}{\PYGZdq{}}\PYG{l+s+se}{\PYGZbs{}n}\PYG{l+s+s2}{null(indefini) = }\PYG{l+s+s2}{\PYGZdq{}}\PYG{p}{;}
\PYG{n}{afficher}  \PYG{p}{(}\PYG{n}{null}\PYG{p}{(}\PYG{n}{indefini}\PYG{p}{)}\PYG{p}{)}\PYG{p}{;}
\PYG{n}{afficher} \PYG{l+s+s2}{\PYGZdq{}}\PYG{l+s+se}{\PYGZbs{}n}\PYG{l+s+s2}{null(0) = }\PYG{l+s+s2}{\PYGZdq{}}\PYG{p}{;}
\PYG{n}{afficher}  \PYG{p}{(}\PYG{n}{null}\PYG{p}{(}\PYG{l+m+mi}{0}\PYG{p}{)}\PYG{p}{)}\PYG{p}{;}
\PYG{n}{afficher} \PYG{l+s+s2}{\PYGZdq{}}\PYG{l+s+se}{\PYGZbs{}n}\PYG{l+s+s2}{null(1) = }\PYG{l+s+s2}{\PYGZdq{}}\PYG{p}{;}
\PYG{n}{afficher}  \PYG{p}{(}\PYG{n}{null}\PYG{p}{(}\PYG{l+m+mi}{1}\PYG{p}{)}\PYG{p}{)}\PYG{p}{;}
\PYG{n}{afficher} \PYG{l+s+s2}{\PYGZdq{}}\PYG{l+s+se}{\PYGZbs{}n}\PYG{l+s+s2}{\PYGZdq{}}\PYG{p}{;}
\PYG{n}{afficher} \PYG{n}{indenter}\PYG{p}{(}\PYG{o}{\PYGZhy{}}\PYG{l+m+mi}{2}\PYG{p}{)}\PYG{p}{;}

\PYG{n}{cible} \PYG{n}{test\PYGZus{}positif}\PYG{p}{:}
\PYG{n}{application} \PYG{p}{:} \PYG{n}{mon\PYGZus{}application}\PYG{p}{;}
\PYG{n}{afficher} \PYG{l+s+s2}{\PYGZdq{}}\PYG{l+s+se}{\PYGZbs{}n}\PYG{l+s+s2}{\PYGZus{}\PYGZus{}POSITIF\PYGZus{}\PYGZus{}}\PYG{l+s+s2}{\PYGZdq{}}\PYG{p}{;}
\PYG{n}{afficher} \PYG{n}{indenter}\PYG{p}{(}\PYG{l+m+mi}{2}\PYG{p}{)}\PYG{p}{;}
\PYG{n}{afficher} \PYG{l+s+s2}{\PYGZdq{}}\PYG{l+s+se}{\PYGZbs{}n}\PYG{l+s+s2}{positif(indefini) = }\PYG{l+s+s2}{\PYGZdq{}}\PYG{p}{;}
\PYG{n}{afficher} \PYG{p}{(}\PYG{n}{positif}\PYG{p}{(}\PYG{n}{indefini}\PYG{p}{)}\PYG{p}{)}\PYG{p}{;}
\PYG{n}{afficher} \PYG{l+s+s2}{\PYGZdq{}}\PYG{l+s+se}{\PYGZbs{}n}\PYG{l+s+s2}{positif(0) = }\PYG{l+s+s2}{\PYGZdq{}}\PYG{p}{;}
\PYG{n}{afficher} \PYG{p}{(}\PYG{n}{positif}\PYG{p}{(}\PYG{l+m+mi}{0}\PYG{p}{)}\PYG{p}{)}\PYG{p}{;}
\PYG{n}{afficher} \PYG{l+s+s2}{\PYGZdq{}}\PYG{l+s+se}{\PYGZbs{}n}\PYG{l+s+s2}{positif(1) = }\PYG{l+s+s2}{\PYGZdq{}}\PYG{p}{;}
\PYG{n}{afficher} \PYG{p}{(}\PYG{n}{positif}\PYG{p}{(}\PYG{l+m+mi}{1}\PYG{p}{)}\PYG{p}{)}\PYG{p}{;}
\PYG{n}{afficher} \PYG{l+s+s2}{\PYGZdq{}}\PYG{l+s+se}{\PYGZbs{}n}\PYG{l+s+s2}{positif(\PYGZhy{}1) = }\PYG{l+s+s2}{\PYGZdq{}}\PYG{p}{;}
\PYG{n}{afficher} \PYG{p}{(}\PYG{n}{positif}\PYG{p}{(}\PYG{o}{\PYGZhy{}}\PYG{l+m+mi}{1}\PYG{p}{)}\PYG{p}{)}\PYG{p}{;}
\PYG{n}{afficher} \PYG{l+s+s2}{\PYGZdq{}}\PYG{l+s+se}{\PYGZbs{}n}\PYG{l+s+s2}{\PYGZdq{}}\PYG{p}{;}
\PYG{n}{afficher} \PYG{n}{indenter}\PYG{p}{(}\PYG{o}{\PYGZhy{}}\PYG{l+m+mi}{2}\PYG{p}{)}\PYG{p}{;}

\PYG{n}{cible} \PYG{n}{test\PYGZus{}positif\PYGZus{}ou\PYGZus{}nul}\PYG{p}{:}
\PYG{n}{application} \PYG{p}{:} \PYG{n}{mon\PYGZus{}application}\PYG{p}{;}
\PYG{n}{afficher} \PYG{l+s+s2}{\PYGZdq{}}\PYG{l+s+se}{\PYGZbs{}n}\PYG{l+s+s2}{\PYGZus{}\PYGZus{}POSITIF OU NUL\PYGZus{}\PYGZus{}}\PYG{l+s+s2}{\PYGZdq{}}\PYG{p}{;}
\PYG{n}{afficher} \PYG{n}{indenter}\PYG{p}{(}\PYG{l+m+mi}{2}\PYG{p}{)}\PYG{p}{;}
\PYG{n}{afficher} \PYG{l+s+s2}{\PYGZdq{}}\PYG{l+s+se}{\PYGZbs{}n}\PYG{l+s+s2}{positif\PYGZus{}ou\PYGZus{}nul(indefini) = }\PYG{l+s+s2}{\PYGZdq{}}\PYG{p}{;}
\PYG{n}{afficher} \PYG{p}{(}\PYG{n}{positif\PYGZus{}ou\PYGZus{}nul}\PYG{p}{(}\PYG{n}{indefini}\PYG{p}{)}\PYG{p}{)}\PYG{p}{;}
\PYG{n}{afficher} \PYG{l+s+s2}{\PYGZdq{}}\PYG{l+s+se}{\PYGZbs{}n}\PYG{l+s+s2}{positif\PYGZus{}ou\PYGZus{}nul(0) = }\PYG{l+s+s2}{\PYGZdq{}}\PYG{p}{;}
\PYG{n}{afficher} \PYG{p}{(}\PYG{n}{positif\PYGZus{}ou\PYGZus{}nul}\PYG{p}{(}\PYG{l+m+mi}{0}\PYG{p}{)}\PYG{p}{)}\PYG{p}{;}
\PYG{n}{afficher} \PYG{l+s+s2}{\PYGZdq{}}\PYG{l+s+se}{\PYGZbs{}n}\PYG{l+s+s2}{positif\PYGZus{}ou\PYGZus{}nul(1) = }\PYG{l+s+s2}{\PYGZdq{}}\PYG{p}{;}
\PYG{n}{afficher} \PYG{p}{(}\PYG{n}{positif\PYGZus{}ou\PYGZus{}nul}\PYG{p}{(}\PYG{l+m+mi}{1}\PYG{p}{)}\PYG{p}{)}\PYG{p}{;}
\PYG{n}{afficher} \PYG{l+s+s2}{\PYGZdq{}}\PYG{l+s+se}{\PYGZbs{}n}\PYG{l+s+s2}{positif\PYGZus{}ou\PYGZus{}nul(\PYGZhy{}1) = }\PYG{l+s+s2}{\PYGZdq{}}\PYG{p}{;}
\PYG{n}{afficher} \PYG{p}{(}\PYG{n}{positif\PYGZus{}ou\PYGZus{}nul}\PYG{p}{(}\PYG{o}{\PYGZhy{}}\PYG{l+m+mi}{1}\PYG{p}{)}\PYG{p}{)}\PYG{p}{;}
\PYG{n}{afficher} \PYG{l+s+s2}{\PYGZdq{}}\PYG{l+s+se}{\PYGZbs{}n}\PYG{l+s+s2}{\PYGZdq{}}\PYG{p}{;}
\PYG{n}{afficher} \PYG{n}{indenter}\PYG{p}{(}\PYG{o}{\PYGZhy{}}\PYG{l+m+mi}{2}\PYG{p}{)}\PYG{p}{;}

\PYG{n}{cible} \PYG{n}{test\PYGZus{}present}\PYG{p}{:}
\PYG{n}{application} \PYG{p}{:} \PYG{n}{mon\PYGZus{}application}\PYG{p}{;}
\PYG{n}{afficher} \PYG{l+s+s2}{\PYGZdq{}}\PYG{l+s+se}{\PYGZbs{}n}\PYG{l+s+s2}{\PYGZus{}\PYGZus{}PRESENT\PYGZus{}\PYGZus{}}\PYG{l+s+s2}{\PYGZdq{}}\PYG{p}{;}
\PYG{n}{afficher} \PYG{n}{indenter}\PYG{p}{(}\PYG{l+m+mi}{2}\PYG{p}{)}\PYG{p}{;}
\PYG{n}{afficher} \PYG{l+s+s2}{\PYGZdq{}}\PYG{l+s+se}{\PYGZbs{}n}\PYG{l+s+s2}{present(indefini) = }\PYG{l+s+s2}{\PYGZdq{}}\PYG{p}{;}
\PYG{n}{afficher}  \PYG{p}{(}\PYG{n}{present}\PYG{p}{(}\PYG{n}{indefini}\PYG{p}{)}\PYG{p}{)}\PYG{p}{;}
\PYG{n}{afficher} \PYG{l+s+s2}{\PYGZdq{}}\PYG{l+s+se}{\PYGZbs{}n}\PYG{l+s+s2}{present(0) = }\PYG{l+s+s2}{\PYGZdq{}}\PYG{p}{;}
\PYG{n}{afficher}  \PYG{p}{(}\PYG{n}{present}\PYG{p}{(}\PYG{l+m+mi}{0}\PYG{p}{)}\PYG{p}{)}\PYG{p}{;}
\PYG{n}{afficher} \PYG{l+s+s2}{\PYGZdq{}}\PYG{l+s+se}{\PYGZbs{}n}\PYG{l+s+s2}{present(1) = }\PYG{l+s+s2}{\PYGZdq{}}\PYG{p}{;}
\PYG{n}{afficher}  \PYG{p}{(}\PYG{n}{present}\PYG{p}{(}\PYG{l+m+mi}{1}\PYG{p}{)}\PYG{p}{)}\PYG{p}{;}
\PYG{n}{afficher} \PYG{l+s+s2}{\PYGZdq{}}\PYG{l+s+se}{\PYGZbs{}n}\PYG{l+s+s2}{\PYGZdq{}}\PYG{p}{;}
\PYG{n}{afficher} \PYG{n}{indenter}\PYG{p}{(}\PYG{o}{\PYGZhy{}}\PYG{l+m+mi}{2}\PYG{p}{)}\PYG{p}{;}

\PYG{n}{cible} \PYG{n}{test\PYGZus{}somme}\PYG{p}{:}
\PYG{n}{application} \PYG{p}{:} \PYG{n}{mon\PYGZus{}application}\PYG{p}{;}
\PYG{n}{afficher} \PYG{l+s+s2}{\PYGZdq{}}\PYG{l+s+se}{\PYGZbs{}n}\PYG{l+s+s2}{\PYGZus{}\PYGZus{}SOMME\PYGZus{}\PYGZus{}}\PYG{l+s+s2}{\PYGZdq{}}\PYG{p}{;}
\PYG{n}{afficher} \PYG{n}{indenter}\PYG{p}{(}\PYG{l+m+mi}{2}\PYG{p}{)}\PYG{p}{;}
\PYG{n}{afficher} \PYG{l+s+s2}{\PYGZdq{}}\PYG{l+s+se}{\PYGZbs{}n}\PYG{l+s+s2}{somme() = }\PYG{l+s+s2}{\PYGZdq{}}\PYG{p}{;}
\PYG{n}{afficher} \PYG{p}{(}\PYG{n}{somme}\PYG{p}{(}\PYG{p}{)}\PYG{p}{)}\PYG{p}{;}
\PYG{n}{afficher} \PYG{l+s+s2}{\PYGZdq{}}\PYG{l+s+se}{\PYGZbs{}n}\PYG{l+s+s2}{somme(indefini) = }\PYG{l+s+s2}{\PYGZdq{}}\PYG{p}{;}
\PYG{n}{afficher} \PYG{p}{(}\PYG{n}{somme}\PYG{p}{(}\PYG{n}{indefini}\PYG{p}{)}\PYG{p}{)}\PYG{p}{;}
\PYG{n}{afficher} \PYG{l+s+s2}{\PYGZdq{}}\PYG{l+s+se}{\PYGZbs{}n}\PYG{l+s+s2}{somme(1, indefini) = }\PYG{l+s+s2}{\PYGZdq{}}\PYG{p}{;}
\PYG{n}{afficher} \PYG{p}{(}\PYG{n}{somme}\PYG{p}{(}\PYG{l+m+mi}{1}\PYG{p}{,} \PYG{n}{indefini}\PYG{p}{)}\PYG{p}{)}\PYG{p}{;}
\PYG{n}{afficher} \PYG{l+s+s2}{\PYGZdq{}}\PYG{l+s+se}{\PYGZbs{}n}\PYG{l+s+s2}{somme(1, 2, 3, 4, 5) = }\PYG{l+s+s2}{\PYGZdq{}}\PYG{p}{;}
\PYG{n}{afficher} \PYG{p}{(}\PYG{n}{somme}\PYG{p}{(}\PYG{l+m+mi}{1}\PYG{p}{,} \PYG{l+m+mi}{2}\PYG{p}{,} \PYG{l+m+mi}{3}\PYG{p}{,} \PYG{l+m+mi}{4}\PYG{p}{,} \PYG{l+m+mi}{5}\PYG{p}{)}\PYG{p}{)}\PYG{p}{;}
\PYG{n}{afficher} \PYG{l+s+s2}{\PYGZdq{}}\PYG{l+s+se}{\PYGZbs{}n}\PYG{l+s+s2}{\PYGZdq{}}\PYG{p}{;}
\PYG{n}{afficher} \PYG{n}{indenter}\PYG{p}{(}\PYG{o}{\PYGZhy{}}\PYG{l+m+mi}{2}\PYG{p}{)}\PYG{p}{;}

\PYG{n}{cible} \PYG{n}{test\PYGZus{}supzero}\PYG{p}{:}
\PYG{n}{application} \PYG{p}{:} \PYG{n}{mon\PYGZus{}application}\PYG{p}{;}
\PYG{n}{afficher} \PYG{l+s+s2}{\PYGZdq{}}\PYG{l+s+se}{\PYGZbs{}n}\PYG{l+s+s2}{\PYGZus{}\PYGZus{}SUPZERO\PYGZus{}\PYGZus{}}\PYG{l+s+s2}{\PYGZdq{}}\PYG{p}{;}
\PYG{n}{afficher} \PYG{n}{indenter}\PYG{p}{(}\PYG{l+m+mi}{2}\PYG{p}{)}\PYG{p}{;}
\PYG{n}{afficher} \PYG{l+s+s2}{\PYGZdq{}}\PYG{l+s+se}{\PYGZbs{}n}\PYG{l+s+s2}{supzero(indefini) = }\PYG{l+s+s2}{\PYGZdq{}}\PYG{p}{;}
\PYG{n}{afficher} \PYG{p}{(}\PYG{n}{supzero}\PYG{p}{(}\PYG{n}{indefini}\PYG{p}{)}\PYG{p}{)}\PYG{p}{;}
\PYG{n}{afficher} \PYG{l+s+s2}{\PYGZdq{}}\PYG{l+s+se}{\PYGZbs{}n}\PYG{l+s+s2}{supzero(42) = }\PYG{l+s+s2}{\PYGZdq{}}\PYG{p}{;}
\PYG{n}{afficher} \PYG{p}{(}\PYG{n}{supzero}\PYG{p}{(}\PYG{l+m+mi}{42}\PYG{p}{)}\PYG{p}{)}\PYG{p}{;}
\PYG{n}{afficher} \PYG{l+s+s2}{\PYGZdq{}}\PYG{l+s+se}{\PYGZbs{}n}\PYG{l+s+s2}{supzero(\PYGZhy{}1) = }\PYG{l+s+s2}{\PYGZdq{}}\PYG{p}{;}
\PYG{n}{afficher} \PYG{p}{(}\PYG{n}{supzero}\PYG{p}{(}\PYG{o}{\PYGZhy{}}\PYG{l+m+mi}{1}\PYG{p}{)}\PYG{p}{)}\PYG{p}{;}
\PYG{n}{afficher} \PYG{l+s+s2}{\PYGZdq{}}\PYG{l+s+se}{\PYGZbs{}n}\PYG{l+s+s2}{supzero(0) = }\PYG{l+s+s2}{\PYGZdq{}}\PYG{p}{;}
\PYG{n}{afficher} \PYG{p}{(}\PYG{n}{supzero}\PYG{p}{(}\PYG{l+m+mi}{0}\PYG{p}{)}\PYG{p}{)}\PYG{p}{;}
\PYG{n}{afficher} \PYG{l+s+s2}{\PYGZdq{}}\PYG{l+s+se}{\PYGZbs{}n}\PYG{l+s+s2}{\PYGZdq{}}\PYG{p}{;}
\PYG{n}{afficher} \PYG{n}{indenter}\PYG{p}{(}\PYG{o}{\PYGZhy{}}\PYG{l+m+mi}{2}\PYG{p}{)}\PYG{p}{;}

\PYG{n}{cible} \PYG{n}{test\PYGZus{}taille}\PYG{p}{:}
\PYG{n}{application} \PYG{p}{:} \PYG{n}{mon\PYGZus{}application}\PYG{p}{;}
\PYG{n}{afficher} \PYG{l+s+s2}{\PYGZdq{}}\PYG{l+s+se}{\PYGZbs{}n}\PYG{l+s+s2}{\PYGZus{}\PYGZus{}TAILLE\PYGZus{}\PYGZus{}}\PYG{l+s+s2}{\PYGZdq{}}\PYG{p}{;}
\PYG{n}{afficher} \PYG{n}{indenter}\PYG{p}{(}\PYG{l+m+mi}{2}\PYG{p}{)}\PYG{p}{;}
\PYG{n}{afficher} \PYG{l+s+s2}{\PYGZdq{}}\PYG{l+s+se}{\PYGZbs{}n}\PYG{l+s+s2}{taille(TAB) = }\PYG{l+s+s2}{\PYGZdq{}}\PYG{p}{;}
\PYG{n}{afficher} \PYG{p}{(}\PYG{n}{taille}\PYG{p}{(}\PYG{n}{TAB}\PYG{p}{)}\PYG{p}{)}\PYG{p}{;}
\PYG{n}{afficher} \PYG{l+s+s2}{\PYGZdq{}}\PYG{l+s+se}{\PYGZbs{}n}\PYG{l+s+s2}{taille(X) = }\PYG{l+s+s2}{\PYGZdq{}}\PYG{p}{;}
\PYG{n}{afficher} \PYG{p}{(}\PYG{n}{taille}\PYG{p}{(}\PYG{n}{X}\PYG{p}{)}\PYG{p}{)}\PYG{p}{;}
\PYG{n}{afficher} \PYG{l+s+s2}{\PYGZdq{}}\PYG{l+s+se}{\PYGZbs{}n}\PYG{l+s+s2}{\PYGZdq{}}\PYG{p}{;}
\PYG{n}{afficher} \PYG{n}{indenter}\PYG{p}{(}\PYG{o}{\PYGZhy{}}\PYG{l+m+mi}{2}\PYG{p}{)}\PYG{p}{;}

\PYG{n}{cible} \PYG{n}{test\PYGZus{}type}\PYG{p}{:}
\PYG{n}{application} \PYG{p}{:} \PYG{n}{mon\PYGZus{}application}\PYG{p}{;}
\PYG{n}{afficher} \PYG{l+s+s2}{\PYGZdq{}}\PYG{l+s+se}{\PYGZbs{}n}\PYG{l+s+s2}{\PYGZus{}\PYGZus{}TYPE\PYGZus{}\PYGZus{}}\PYG{l+s+s2}{\PYGZdq{}}\PYG{p}{;}
\PYG{n}{afficher} \PYG{n}{indenter}\PYG{p}{(}\PYG{l+m+mi}{2}\PYG{p}{)}\PYG{p}{;}
\PYG{n}{afficher} \PYG{l+s+s2}{\PYGZdq{}}\PYG{l+s+se}{\PYGZbs{}n}\PYG{l+s+s2}{type(X, REEL) = }\PYG{l+s+s2}{\PYGZdq{}}\PYG{p}{;}
\PYG{n}{afficher} \PYG{p}{(}\PYG{n+nb}{type}\PYG{p}{(}\PYG{n}{X}\PYG{p}{,} \PYG{n}{REEL}\PYG{p}{)}\PYG{p}{)}\PYG{p}{;}
\PYG{n}{afficher} \PYG{l+s+s2}{\PYGZdq{}}\PYG{l+s+se}{\PYGZbs{}n}\PYG{l+s+s2}{type(X, ENTIER) = }\PYG{l+s+s2}{\PYGZdq{}}\PYG{p}{;}
\PYG{n}{afficher} \PYG{p}{(}\PYG{n+nb}{type}\PYG{p}{(}\PYG{n}{X}\PYG{p}{,} \PYG{n}{ENTIER}\PYG{p}{)}\PYG{p}{)}\PYG{p}{;}
\PYG{n}{afficher} \PYG{l+s+s2}{\PYGZdq{}}\PYG{l+s+se}{\PYGZbs{}n}\PYG{l+s+s2}{type(TAB, ENTIER) = }\PYG{l+s+s2}{\PYGZdq{}}\PYG{p}{;}
\PYG{n}{afficher} \PYG{p}{(}\PYG{n+nb}{type}\PYG{p}{(}\PYG{n}{TAB}\PYG{p}{,} \PYG{n}{ENTIER}\PYG{p}{)}\PYG{p}{)}\PYG{p}{;}
\PYG{n}{afficher} \PYG{l+s+s2}{\PYGZdq{}}\PYG{l+s+se}{\PYGZbs{}n}\PYG{l+s+s2}{type(Y, ENTIER) = }\PYG{l+s+s2}{\PYGZdq{}}\PYG{p}{;}
\PYG{n}{afficher} \PYG{p}{(}\PYG{n+nb}{type}\PYG{p}{(}\PYG{n}{Y}\PYG{p}{,} \PYG{n}{ENTIER}\PYG{p}{)}\PYG{p}{)}\PYG{p}{;}
\PYG{n}{afficher} \PYG{l+s+s2}{\PYGZdq{}}\PYG{l+s+se}{\PYGZbs{}n}\PYG{l+s+s2}{type(Y, REEL) = }\PYG{l+s+s2}{\PYGZdq{}}\PYG{p}{;}
\PYG{n}{afficher} \PYG{p}{(}\PYG{n+nb}{type}\PYG{p}{(}\PYG{n}{Y}\PYG{p}{,} \PYG{n}{REEL}\PYG{p}{)}\PYG{p}{)}\PYG{p}{;}
\PYG{n}{afficher} \PYG{l+s+s2}{\PYGZdq{}}\PYG{l+s+se}{\PYGZbs{}n}\PYG{l+s+s2}{type(Y, BOOLEEN) = }\PYG{l+s+s2}{\PYGZdq{}}\PYG{p}{;}
\PYG{n}{afficher} \PYG{p}{(}\PYG{n+nb}{type}\PYG{p}{(}\PYG{n}{Y}\PYG{p}{,} \PYG{n}{BOOLEEN}\PYG{p}{)}\PYG{p}{)}\PYG{p}{;}
\PYG{n}{afficher} \PYG{l+s+s2}{\PYGZdq{}}\PYG{l+s+se}{\PYGZbs{}n}\PYG{l+s+s2}{type(Y, DATE\PYGZus{}AAAA) = }\PYG{l+s+s2}{\PYGZdq{}}\PYG{p}{;}
\PYG{n}{afficher} \PYG{p}{(}\PYG{n+nb}{type}\PYG{p}{(}\PYG{n}{Y}\PYG{p}{,} \PYG{n}{DATE\PYGZus{}AAAA}\PYG{p}{)}\PYG{p}{)}\PYG{p}{;}
\PYG{n}{afficher} \PYG{l+s+s2}{\PYGZdq{}}\PYG{l+s+se}{\PYGZbs{}n}\PYG{l+s+s2}{type(Y, DATE\PYGZus{}JJMMAAAA) = }\PYG{l+s+s2}{\PYGZdq{}}\PYG{p}{;}
\PYG{n}{afficher} \PYG{p}{(}\PYG{n+nb}{type}\PYG{p}{(}\PYG{n}{Y}\PYG{p}{,} \PYG{n}{DATE\PYGZus{}JJMMAAAA}\PYG{p}{)}\PYG{p}{)}\PYG{p}{;}
\PYG{n}{afficher} \PYG{l+s+s2}{\PYGZdq{}}\PYG{l+s+se}{\PYGZbs{}n}\PYG{l+s+s2}{type(Y, DATE\PYGZus{}MM) = }\PYG{l+s+s2}{\PYGZdq{}}\PYG{p}{;}
\PYG{n}{afficher} \PYG{p}{(}\PYG{n+nb}{type}\PYG{p}{(}\PYG{n}{Y}\PYG{p}{,} \PYG{n}{DATE\PYGZus{}MM}\PYG{p}{)}\PYG{p}{)}\PYG{p}{;}
\PYG{n}{afficher} \PYG{l+s+s2}{\PYGZdq{}}\PYG{l+s+se}{\PYGZbs{}n}\PYG{l+s+s2}{\PYGZdq{}}\PYG{p}{;}
\PYG{n}{afficher} \PYG{n}{indenter}\PYG{p}{(}\PYG{o}{\PYGZhy{}}\PYG{l+m+mi}{2}\PYG{p}{)}\PYG{p}{;}

\PYG{n}{cible} \PYG{n}{fun\PYGZus{}test}\PYG{p}{:}
\PYG{n}{application} \PYG{p}{:} \PYG{n}{mon\PYGZus{}application}\PYG{p}{;}
\PYG{c+c1}{\PYGZsh{} Abs}
\PYG{n}{calculer} \PYG{n}{cible} \PYG{n}{test\PYGZus{}abs}\PYG{p}{;}
\PYG{c+c1}{\PYGZsh{} Arr}
\PYG{n}{calculer} \PYG{n}{cible} \PYG{n}{test\PYGZus{}arr}\PYG{p}{;}
\PYG{c+c1}{\PYGZsh{} Attribut}
\PYG{n}{calculer} \PYG{n}{cible} \PYG{n}{test\PYGZus{}attribut}\PYG{p}{;}
\PYG{c+c1}{\PYGZsh{} Champ evenement}
\PYG{n}{calculer} \PYG{n}{cible} \PYG{n}{test\PYGZus{}champ\PYGZus{}evenement}\PYG{p}{;}
\PYG{c+c1}{\PYGZsh{} Inf}
\PYG{n}{calculer} \PYG{n}{cible} \PYG{n}{test\PYGZus{}inf}\PYG{p}{;}
\PYG{c+c1}{\PYGZsh{} Max}
\PYG{n}{calculer} \PYG{n}{cible} \PYG{n}{test\PYGZus{}max}\PYG{p}{;}
\PYG{c+c1}{\PYGZsh{} Meme variable}
\PYG{n}{calculer} \PYG{n}{cible} \PYG{n}{test\PYGZus{}meme\PYGZus{}variable}\PYG{p}{;}
\PYG{c+c1}{\PYGZsh{} Min}
\PYG{n}{calculer} \PYG{n}{cible} \PYG{n}{test\PYGZus{}min}\PYG{p}{;}
\PYG{c+c1}{\PYGZsh{} Multimax}
\PYG{n}{calculer} \PYG{n}{cible} \PYG{n}{test\PYGZus{}multimax}\PYG{p}{;}
\PYG{c+c1}{\PYGZsh{} Null}
\PYG{n}{calculer} \PYG{n}{cible} \PYG{n}{test\PYGZus{}nb\PYGZus{}evenements}\PYG{p}{;}
\PYG{c+c1}{\PYGZsh{} Null}
\PYG{n}{calculer} \PYG{n}{cible} \PYG{n}{test\PYGZus{}null}\PYG{p}{;}
\PYG{c+c1}{\PYGZsh{} Positif}
\PYG{n}{calculer} \PYG{n}{cible} \PYG{n}{test\PYGZus{}positif}\PYG{p}{;}
\PYG{c+c1}{\PYGZsh{} Positif ou nul}
\PYG{n}{calculer} \PYG{n}{cible} \PYG{n}{test\PYGZus{}positif\PYGZus{}ou\PYGZus{}nul}\PYG{p}{;}
\PYG{c+c1}{\PYGZsh{} Présent}
\PYG{n}{calculer} \PYG{n}{cible} \PYG{n}{test\PYGZus{}present}\PYG{p}{;}
\PYG{c+c1}{\PYGZsh{} Somme}
\PYG{n}{calculer} \PYG{n}{cible} \PYG{n}{test\PYGZus{}somme}\PYG{p}{;}
\PYG{c+c1}{\PYGZsh{} Supzero}
\PYG{n}{calculer} \PYG{n}{cible} \PYG{n}{test\PYGZus{}supzero}\PYG{p}{;}
\PYG{c+c1}{\PYGZsh{} Taille}
\PYG{n}{calculer} \PYG{n}{cible} \PYG{n}{test\PYGZus{}taille}\PYG{p}{;}
\PYG{c+c1}{\PYGZsh{} Type}
\PYG{n}{calculer} \PYG{n}{cible} \PYG{n}{test\PYGZus{}type}\PYG{p}{;}
\end{sphinxVerbatim}


\subsection{Fichier irj : test.irj}
\label{\detokenize{exemples/fonctions:fichier-irj-test-irj}}
\begin{sphinxVerbatim}[commandchars=\\\{\}]
\PYG{c+c1}{\PYGZsh{}NOM}
\PYG{n}{MON}\PYG{o}{\PYGZhy{}}\PYG{n}{TEST}
\PYG{c+c1}{\PYGZsh{}ENTREES\PYGZhy{}PRIMITIF}
\PYG{c+c1}{\PYGZsh{}CONTROLES\PYGZhy{}PRIMITIF}
\PYG{c+c1}{\PYGZsh{}RESULTATS\PYGZhy{}PRIMITIF}
\PYG{c+c1}{\PYGZsh{}ENTREES\PYGZhy{}CORRECTIF}
\PYG{c+c1}{\PYGZsh{}CONTROLES\PYGZhy{}CORRECTIF}
\PYG{c+c1}{\PYGZsh{}RESULTATS\PYGZhy{}CORRECTIF}
\PYG{c+c1}{\PYGZsh{}\PYGZsh{}}
\end{sphinxVerbatim}


\subsection{Commande}
\label{\detokenize{exemples/fonctions:commande}}
\begin{sphinxVerbatim}[commandchars=\\\{\}]
\PYG{n}{mlang} \PYG{n}{test}\PYG{o}{.}\PYG{n}{m} \PYGZbs{}
	\PYG{o}{\PYGZhy{}}\PYG{o}{\PYGZhy{}}\PYG{n}{without\PYGZus{}dfgip\PYGZus{}m} \PYGZbs{}
	\PYG{o}{\PYGZhy{}}\PYG{n}{A} \PYG{n}{mon\PYGZus{}application} \PYGZbs{}
	\PYG{o}{\PYGZhy{}}\PYG{o}{\PYGZhy{}}\PYG{n}{mpp\PYGZus{}function} \PYG{n}{fun\PYGZus{}test} \PYGZbs{}
	\PYG{o}{\PYGZhy{}}\PYG{o}{\PYGZhy{}}\PYG{n}{run\PYGZus{}test} \PYG{n}{test}\PYG{o}{.}\PYG{n}{irj}
\end{sphinxVerbatim}


\chapter{Documentation développeur}
\label{\detokenize{index:documentation-developpeur}}\begin{itemize}
\item {} 
\sphinxAtStartPar
Index

\end{itemize}



\renewcommand{\indexname}{Index}
\printindex
\end{document}